% Présentation du langage de définition des algorithmes — Informatique 3ème — Cours signé OnBuch+
% Programme officiel MINESEC. Compiler avec : tectonic N13-presentation-langage-definition-algorithmes.tex
\newcommand{\DOCMATIERE}{Informatique}
\newcommand{\DOCNIVEAU}{3ème}
\newcommand{\DOCMODULE}{Informatique – classe de 3ème}
\newcommand{\DOCLECON}{N13}
\newcommand{\DOCTITRE}{Présentation du langage de définition des algorithmes}
% OnBuch+ — préambule « cours signés » ENRICHI (compilé avec tectonic / XeLaTeX)
% Système visuel « Pop » d'OnBuch+ : fond crème (jamais blanc), bordures encre
% épaisses, ombres « dures » décalées, titres Archivo Black en capitales.
% Version MATHÉMATIQUES : graphes (pgfplots/tikz), boîtes pédagogiques
% (définition, propriété, méthode, attention, le savais-tu, exercice résolu,
% exercice, TP, bilan), page de garde et signature OnBuch+ ancrée en bas.
%
% Macros attendues dans main.tex AVANT \input{preamble} :
%   \DOCMATIERE  \DOCNIVEAU  \DOCMODULE  \DOCLECON (numéro)  \DOCTITRE
\documentclass[11pt]{article}

\usepackage{fontspec}
\usepackage[french]{babel}
\usepackage[a4paper,margin=2cm,top=3.4cm,bottom=2.8cm,headheight=1.4cm,footskip=1.3cm]{geometry}
\usepackage{amsmath,amssymb,amsfonts}
\usepackage{mathtools} % \xmapsto, \coloneqq... souvent utilisés par les modèles
\usepackage{cancel} % simplifications barrées (bilans de forces)
% \abs{z} (module, valeur absolue) et \norm{u} (norme), utilisés par les modèles
\providecommand{\abs}{}\let\abs\relax\DeclarePairedDelimiter{\abs}{\lvert}{\rvert}
\providecommand{\norm}{}\let\norm\relax\DeclarePairedDelimiter{\norm}{\lVert}{\rVert}
\usepackage[table]{xcolor} % [table] : fournit aussi \rowcolors (alternance de lignes)
\usepackage{pagecolor}
\usepackage{tikz}
\usetikzlibrary{arrows.meta,positioning,calc,shapes.geometric,shapes.misc,shapes.arrows,shapes.symbols,decorations.pathmorphing,decorations.pathreplacing,decorations.markings,patterns,shadows,mindmap,trees,fit,backgrounds,matrix,intersections,angles,quotes}
\usepackage{pgfplots}
\usepackage[version=4]{mhchem} % équations nucléaires \ce{^{235}_{92}U}
\usepackage[european,siunitx]{circuitikz} % schémas électriques
\pgfplotsset{compat=1.18}
\usepgfplotslibrary{fillbetween,groupplots} % aires entre courbes (calcul intégral)
\usepackage{enumitem}
\usepackage{fancyhdr}
% [most] charge toutes les bibliothèques tcolorbox (listings, minted, raster,
% documentation...) et gonfle inutilement la mémoire TeX sur un document long
% (~300 boîtes) : on ne charge que ce qui est réellement utilisé.
\usepackage[skins,breakable]{tcolorbox}
\usepackage{graphicx}
\usepackage{colortbl}
\usepackage{booktabs}
\usepackage{tabularx}
\usepackage{diagbox} % case d'en-tête diagonale des tableaux à double entrée
% Colonnes extensibles centrée / gauche / droite (C, L, R), souvent utilisées par les modèles
\newcolumntype{C}{>{\centering\arraybackslash}X}
\newcolumntype{L}{>{\raggedright\arraybackslash}X}
\newcolumntype{R}{>{\raggedleft\arraybackslash}X}
\usepackage{array}
\usepackage{multirow}
\usepackage{multicol}
\usepackage{siunitx}
% parse-numbers=false : accepte aussi \SI{2,5 \times 10^{-3}}{...}
\sisetup{locale=FR,output-decimal-marker={,},group-minimum-digits=4,parse-numbers=false}
\DeclareSIUnit\ohm{\ensuremath{\symup{\Omega}}} % U+2126 absent de la police
\DeclareSIUnit\division{div} % réglages d'oscilloscope (ms/div, V/div)
\DeclareSIUnit\tr{tr} % tours (vitesse de rotation, ex. \SI{1500}{\tr\per\minute})
\DeclareSIUnit\bit{bit}
\DeclareSIUnit\byte{o} % octet
\DeclareSIUnit\inch{po} % pouce
\DeclareSIUnit\ppp{ppp} % points par pouce (résolution d'image)
\usepackage{qrcode}
% Blocs de code source (PHP, C, HTML, JavaScript, SQL) — spécifique à la
% matière Informatique : listings + habillage aux couleurs OnBuch+ (voir le
% style \lstset ci-dessous, à côté des autres boîtes pédagogiques).
\usepackage{listings}
% \degree : commande fréquemment utilisée par les modèles pour un angle ou une
% température en dehors de \SI{}{} (non fournie par les paquets chargés).
\providecommand{\degree}{^{\circ}}
% \qed (fin de démonstration), utilisé par les modèles sans amsthm
\providecommand{\qed}{\hfill$\square$}
% \barre : double barre des valeurs interdites dans un tableau de variations (tkz-tab)
\providecommand{\barre}{\ensuremath{\|}}
% environnement proof (amsthm non chargé) utilisé par les modèles
\ifdefined\proof\else\newenvironment{proof}[1][Démonstration]{\par\noindent\textit{#1.}\ }{\qed\par}\fi
\usepackage{lastpage}
\usepackage{hyperref}
\hypersetup{hidelinks}

% ============================================================
% Palette « Pop » OnBuch+ — mêmes hex que la classe P (plus_ui.dart)
% ============================================================
\definecolor{poporange}{HTML}{F59321}
\definecolor{popdark}{HTML}{E07A0C}
\definecolor{popcream}{HTML}{FFF6E8}
\definecolor{popink}{HTML}{1C1714}
\definecolor{poppurple}{HTML}{7A5AE0}
\definecolor{popgreen}{HTML}{2E9E6A}
\definecolor{poppink}{HTML}{E0457B}
\definecolor{popblue}{HTML}{2F7DE1}
\definecolor{popgold}{HTML}{C98A12}
\definecolor{popmuted}{HTML}{8A7F76}
\definecolor{popline}{HTML}{EADFD2}
\definecolor{popgray}{HTML}{8A7F76}
\colorlet{popgrey}{popgray}
% Alias tolérants (le modèle nomme parfois les couleurs différemment)
\colorlet{popwhite}{white}
\colorlet{popblack}{popink}
\colorlet{popred}{poppink}
\colorlet{popyellow}{popgold}
% Teintes claires pour fonds de tableaux / zones
\colorlet{popblueL}{popblue!10!white}
\colorlet{poporangeL}{poporange!14!white}
\colorlet{popgreenL}{popgreen!12!white}
\colorlet{poppurpleL}{poppurple!12!white}
\colorlet{poppinkL}{poppink!12!white}
\colorlet{popgoldL}{popgold!14!white}

\pagecolor{popcream}

% ============================================================
% Typographie
% ============================================================
\defaultfontfeatures{Ligatures=TeX}
\setmainfont{PlusJakartaSans-Medium.ttf}[Path=../../../fonts/,BoldFont=PlusJakartaSans-Bold.ttf]
% Maths en sans-serif (Fira Math) avec chiffres et lettres droites de Plus
% Jakarta, pour que les formules se fondent dans le texte.
\usepackage[math-style=ISO,bold-style=ISO]{unicode-math}
\setmathfont{FiraMath-Regular.otf}[Path=../../../fonts/]
\setmathfont{PlusJakartaSans-Medium.ttf}[Path=../../../fonts/,range={up/num,up/{Latin,latin},\mathup}]
\usepackage{icomma} % virgule décimale sans espace en mode math
% Caractères mathématiques Unicode absents des polices de texte (Plus Jakarta,
% Archivo Black) : rendus par leur équivalent mathématique, en texte comme en
% formule — sinon ils s'affichent en carré vide (ex. « Arithmétique dans ℤ »).
\usepackage{newunicodechar}
\newunicodechar{ℝ}{\ensuremath{\mathbb{R}}}
\newunicodechar{ℤ}{\ensuremath{\mathbb{Z}}}
\newunicodechar{ℂ}{\ensuremath{\mathbb{C}}}
\newunicodechar{ℕ}{\ensuremath{\mathbb{N}}}
\newunicodechar{ℚ}{\ensuremath{\mathbb{Q}}}
\newunicodechar{𝔻}{\ensuremath{\mathbb{D}}}
\newunicodechar{α}{\ensuremath{\alpha}}
\newunicodechar{β}{\ensuremath{\beta}}
\newunicodechar{γ}{\ensuremath{\gamma}}
\newunicodechar{δ}{\ensuremath{\delta}}
\newunicodechar{ε}{\ensuremath{\varepsilon}}
\newunicodechar{θ}{\ensuremath{\theta}}
\newunicodechar{λ}{\ensuremath{\lambda}}
\newunicodechar{μ}{\ensuremath{\mu}}
\newunicodechar{σ}{\ensuremath{\sigma}}
\newunicodechar{τ}{\ensuremath{\tau}}
\newunicodechar{φ}{\ensuremath{\varphi}}
\newunicodechar{ψ}{\ensuremath{\psi}}
\newunicodechar{ω}{\ensuremath{\omega}}
\newunicodechar{Σ}{\ensuremath{\Sigma}}
% majuscules produites par \MakeUppercase dans les titres (α-aminés -> Α)
\newunicodechar{Α}{\ensuremath{\alpha}}
\newunicodechar{Β}{\ensuremath{\beta}}
\newunicodechar{Γ}{\ensuremath{\Gamma}}
\newunicodechar{Δ}{\ensuremath{\Delta}}
\newunicodechar{Ε}{\ensuremath{\varepsilon}}
\newunicodechar{Θ}{\ensuremath{\Theta}}
\newunicodechar{Λ}{\ensuremath{\Lambda}}
\newunicodechar{Μ}{\ensuremath{\mu}}
\newunicodechar{Τ}{\ensuremath{\tau}}
\newunicodechar{Φ}{\ensuremath{\Phi}}
\newunicodechar{Ψ}{\ensuremath{\Psi}}
\newunicodechar{Ω}{\ensuremath{\Omega}}
\newunicodechar{↦}{\ensuremath{\mapsto}}
\newunicodechar{∈}{\ensuremath{\in}}
\newunicodechar{∉}{\ensuremath{\notin}}
\newunicodechar{∧}{\ensuremath{\wedge}}
\newunicodechar{⇔}{\ensuremath{\Leftrightarrow}}
\newunicodechar{⟺}{\ensuremath{\Longleftrightarrow}}
\newunicodechar{⇒}{\ensuremath{\Rightarrow}}
\newunicodechar{⟹}{\ensuremath{\Longrightarrow}}
\newunicodechar{∘}{\ensuremath{\circ}}
\newunicodechar{∪}{\ensuremath{\cup}}
\newunicodechar{∩}{\ensuremath{\cap}}
\newunicodechar{≡}{\ensuremath{\equiv}}
\newunicodechar{⊂}{\ensuremath{\subset}}
\newunicodechar{⊥}{\ensuremath{\perp}}
\newunicodechar{∃}{\ensuremath{\exists}}
\newunicodechar{∀}{\ensuremath{\forall}}
\newunicodechar{∛}{\ensuremath{\sqrt[3]{\,}}}
\newunicodechar{∜}{\ensuremath{\sqrt[4]{\,}}}
\newunicodechar{⃗}{\ensuremath{{}^{\to}}}
\newunicodechar{ⁿ}{\ifmmode^{n}\else\textsuperscript{\ensuremath{n}}\fi}
\newunicodechar{ˣ}{\ifmmode^{x}\else\textsuperscript{\ensuremath{x}}\fi}
\newunicodechar{ᵇ}{\ifmmode^{b}\else\textsuperscript{\ensuremath{b}}\fi}
\newunicodechar{ʳ}{\ifmmode^{r}\else\textsuperscript{\ensuremath{r}}\fi}
\newunicodechar{ᵘ}{\ifmmode^{u}\else\textsuperscript{\ensuremath{u}}\fi}
\newunicodechar{ʸ}{\ifmmode^{y}\else\textsuperscript{\ensuremath{y}}\fi}
\newunicodechar{ᵃ}{\ifmmode^{a}\else\textsuperscript{\ensuremath{a}}\fi}
\newunicodechar{ᵉ}{\ifmmode^{e}\else\textsuperscript{\ensuremath{e}}\fi}
\newunicodechar{ᵏ}{\ifmmode^{k}\else\textsuperscript{\ensuremath{k}}\fi}
\newunicodechar{ᵖ}{\ifmmode^{p}\else\textsuperscript{\ensuremath{p}}\fi}
\newunicodechar{ᴺ}{\ifmmode^{N}\else\textsuperscript{\ensuremath{N}}\fi}
\newunicodechar{ᶜ}{\ifmmode^{c}\else\textsuperscript{\ensuremath{c}}\fi}
\newunicodechar{ʰ}{\ifmmode^{h}\else\textsuperscript{\ensuremath{h}}\fi}
\newunicodechar{ᵠ}{\ifmmode^{q}\else\textsuperscript{\ensuremath{q}}\fi}
\newunicodechar{ₙ}{\ifmmode_{n}\else\textsubscript{\ensuremath{n}}\fi}
\newunicodechar{ₐ}{\ifmmode_{a}\else\textsubscript{\ensuremath{a}}\fi}
\newunicodechar{ᵢ}{\ifmmode_{i}\else\textsubscript{\ensuremath{i}}\fi}
\newunicodechar{ᵣ}{\ifmmode_{r}\else\textsubscript{\ensuremath{r}}\fi}
\newunicodechar{ₖ}{\ifmmode_{k}\else\textsubscript{\ensuremath{k}}\fi}
\newunicodechar{ᵦ}{\ifmmode_{\beta}\else\textsubscript{\ensuremath{\beta}}\fi}
\newunicodechar{ₓ}{\ifmmode_{x}\else\textsubscript{\ensuremath{x}}\fi}
\newfontfamily\popdisplay{ArchivoBlack-Regular.ttf}[Path=../../../fonts/]
\newfontfamily\poplogo{SpaceGrotesk-Bold.ttf}[Path=../../../fonts/]
\newfontfamily\popsemi{PlusJakartaSans-SemiBold.ttf}[Path=../../../fonts/]
\newfontfamily\popxbold{PlusJakartaSans-ExtraBold.ttf}[Path=../../../fonts/]
\newfontfamily\popxboldls{PlusJakartaSans-ExtraBold.ttf}[Path=../../../fonts/,LetterSpace=3.0]

\setlist[enumerate]{leftmargin=1.7em,itemsep=5pt,topsep=4pt}
\setlist[itemize]{leftmargin=1.7em,itemsep=4pt,topsep=4pt,label={\color{poporange}\textbullet}}
\setlength{\parindent}{0pt}
\setlength{\parskip}{5pt}
\renewcommand{\arraystretch}{1.25}

% pgfplots : style Pop par défaut
\pgfplotsset{
  every axis/.append style={axis line style={popink,line width=1pt},tick style={popink},
    grid style={popline,line width=0.5pt},label style={font=\small},tick label style={font=\footnotesize}},
  popcurve/.style={poporange,line width=1.8pt,no markers,smooth},
}
% Même style disponible dans un \draw TikZ simple (hors pgfplots)
\tikzset{popcurve/.style={poporange,line width=1.8pt}}
\tikzset{popgrid/.style={popline,very thin}}
\colorlet{popcurve}{poporange} % le nom du style est parfois utilisé comme couleur

% ============================================================
% Pill / squiggle
% ============================================================
\newcommand{\poppill}[3]{%
  \tikz[baseline=-0.55ex]\node[draw=popink,line width=1.2pt,rounded corners=2.6pt,%
    fill=#2,inner xsep=6.5pt,inner ysep=3pt]{%
    \popxboldls\fontsize{7}{7}\selectfont\color{#3}\MakeUppercase{#1}};%
}
\newcommand{\popsquiggle}[1]{%
  \tikz[baseline=-0.4ex]{
    \draw[#1,line width=2.6pt,line cap=round]
      plot [smooth] coordinates {(0,0.09) (0.34,0.01) (0.68,0.21) (1.02,0.01) (1.36,0.10)};
  }%
}
% Mot-clé surligné (marqueur orange sous le texte)
\newcommand{\cle}[1]{{\popxbold\color{popdark}#1}}

% ============================================================
% En-tête / pied
% ============================================================
\pagestyle{fancy}
\fancyhf{}
\renewcommand{\headrulewidth}{0pt}
\renewcommand{\footrulewidth}{0pt}
\lhead{\raisebox{-0.15ex}{\poplogo\fontsize{13}{13}\selectfont\color{popink}OnBuch\color{poporange}+}}
\rhead{\poppill{\DOCMATIERE\ \textbullet\ \DOCNIVEAU\ \textbullet\ Leçon \DOCLECON}{white}{popink}}
\lfoot{\popxbold\fontsize{7.6}{7.6}\selectfont\color{popmuted}\MakeUppercase{Cours signé OnBuch+}\ \textbullet\ {\popsemi\DOCTITRE}}
\rfoot{\tikz[baseline=-0.6ex]\node[rounded corners=8.5pt,fill=popink,minimum height=17pt,minimum width=17pt,inner xsep=5pt,inner sep=0]{\popxbold\fontsize{8}{8}\selectfont\color{popcream}\thepage\,/\,\pageref*{LastPage}};}

% ============================================================
% Titres
% ============================================================
\newcommand{\chaptitre}[2]{%
  \begin{tcolorbox}[enhanced,colback=poporange,colframe=popink,boxrule=2.6pt,arc=11pt,
      drop shadow=popink,left=15pt,right=15pt,top=12pt,bottom=11pt,before skip=3pt,after skip=18pt]
    \poppill{#2}{popink}{popcream}\par\vspace{9pt}
    {\raggedright\hyphenpenalty=10000\exhyphenpenalty=10000\popdisplay\fontsize{22}{24}\selectfont\color{popcream}\MakeUppercase{#1}\par}
  \end{tcolorbox}
}
% Section numérotée : pastille orange avec le numéro + titre Archivo
\newcounter{popsec}
\newcommand{\coursec}[1]{%
  \stepcounter{popsec}\setcounter{popsub}{0}%
  \par\vspace{16pt}\noindent
  \begin{minipage}{\linewidth}
  \tikz[baseline=-0.7ex]\node[draw=popink,line width=1.6pt,fill=poporange,rounded corners=4pt,minimum size=22pt,inner sep=0]{\popdisplay\fontsize{12}{12}\selectfont\color{popcream}\thepopsec};%
  \hspace{8pt}{\popdisplay\fontsize{15}{17}\selectfont\color{popink}\MakeUppercase{#1}}\par
  \vspace{4pt}\noindent{\color{poporange}\rule{36pt}{3.6pt}}
  \end{minipage}\par\nopagebreak\vspace{8pt}
}
\newcounter{popsub}
\newcommand{\courssub}[1]{%
  \stepcounter{popsub}%
  \par\vspace{9pt}\noindent{\popxbold\fontsize{11.5}{13}\selectfont\color{popink}{\color{poporange}\thepopsec.\thepopsub}\hspace{6pt}#1}\par\nopagebreak\vspace{4pt}
}

% ============================================================
% Boîtes pédagogiques — même recette : fond blanc, bordure épaisse colorée,
% ombre dure encre, badge plein. Titre optionnel [#1].
% ============================================================
\tcbset{popbase/.style={breakable,enhanced,colback=white,boxrule=2.3pt,arc=9pt,
  shadow={3pt}{-3pt}{0pt}{popink},left=14pt,right=14pt,top=11pt,bottom=11pt,before skip=11pt,after skip=15pt,
  fontupper=\popsemi\fontsize{10.6}{14.5}\selectfont\color{popink},
  fontlower=\popsemi\fontsize{10.6}{14.5}\selectfont\color{popink}}}
% Coche dessinée (le glyphe ✓ n'existe pas dans les polices du document)
\newcommand{\popcheck}[1]{\tikz[baseline=-0.5ex]\draw[#1,line width=1.6pt,line cap=round,line join=round](-0.13,0.0)--(-0.04,-0.09)--(0.13,0.1);}
\providecommand{\coche}{\popcheck{popgreen}}
\providecommand{\resultat}[1]{\par\smallskip\noindent\fcolorbox{popgreen}{popgreenL}{\parbox{\dimexpr\linewidth-2\fboxsep-2\fboxrule\relax}{\textbf{Résultat.} #1}}\par\smallskip}
\newcommand{\popboxhead}[3]{\poppill{#1}{#2}{white}\ifx\relax#3\relax\else\hspace{6pt}{\popxbold\fontsize{10.6}{12}\selectfont\color{popink}#3}\fi\par\vspace{7pt}}

\newtcolorbox{aretenir}[1][]{popbase,colframe=popgreen,before upper={\popboxhead{À retenir}{popgreen}{#1}}}
\newtcolorbox{exemplebox}[1][]{popbase,colframe=poporange,before upper={\popboxhead{Exemple}{poporange}{#1}}}
\newtcolorbox{definition}[1][]{popbase,colframe=popblue,before upper={\popboxhead{Définition}{popblue}{#1}}}
\newtcolorbox{propriete}[1][]{popbase,colframe=poppurple,before upper={\popboxhead{Propriété}{poppurple}{#1}}}
\newtcolorbox{methode}[1][]{popbase,colframe=popgold,before upper={\popboxhead{Méthode}{popgold}{#1}}}
\newtcolorbox{attention}[1][]{popbase,colframe=poppink,before upper={\popboxhead{Attention}{poppink}{#1}}}
\newtcolorbox{experience}[1][]{popbase,colframe=popblue,colback=popblueL,before upper={\popboxhead{Expérience}{popblue}{#1}}}
% « Le savais-tu ? » : carte violette pleine (contexte, culture, Cameroun)
\newtcolorbox{savaistu}[1][]{popbase,colback=poppurple,colframe=popink,
  fontupper=\popsemi\fontsize{10.4}{14.2}\selectfont\color{white},
  before upper={\poppill{Le savais-tu\,?}{white}{poppurple}\ifx\relax#1\relax\else\hspace{6pt}{\popxbold\color{white}#1}\fi\par\vspace{7pt}}}
% Exercice résolu : énoncé en haut, \tcblower puis la solution sur fond crème
\newcounter{popexo}\newcounter{popexr}
\newtcolorbox{exoresolu}[1][]{popbase,colframe=poporange,colbacklower=popcream,
  segmentation style={popink,line width=1.2pt,dashed},
  before upper={\stepcounter{popexr}\popboxhead{Exercice résolu \thepopexr}{poporange}{#1}},
  before lower={\poppill{Solution}{popink}{popcream}\par\vspace{6pt}}}
% Exercice (non résolu en place) — la correction est dans la partie Corrigés
\newtcolorbox{exercice}[1][]{popbase,colframe=popink,
  before upper={\stepcounter{popexo}\popboxhead{Exercice \thepopexo}{popink}{#1}}}
% Corrigé
\newtcolorbox{corrige}[1][]{popbase,colframe=popgreen,colback=popgreenL,
  before upper={\popboxhead{Corrigé}{popgreen}{#1}}}
% Objectifs / prérequis / bilan
\newtcolorbox{objectifs}{popbase,colframe=popink,colback=poporangeL,
  before upper={\popboxhead{Objectifs de la leçon}{popink}{}}}
\newtcolorbox{prerequis}{popbase,colframe=popink,colback=popblueL,
  before upper={\popboxhead{Prérequis}{popblue}{}}}
\newtcolorbox{bilan}{popbase,colframe=popink,colback=white,boxrule=2.8pt,
  before upper={\popboxhead{Fiche bilan}{poporange}{L'essentiel à connaître pour l'examen}}}
% Figure encadrée avec légende
\newtcolorbox{popfigure}[1][]{enhanced,breakable=false,colback=white,colframe=popline,boxrule=1.4pt,arc=8pt,
  left=10pt,right=10pt,top=10pt,bottom=8pt,before skip=10pt,after skip=12pt,halign=center,
  lower separated=false,halign lower=center,
  fontlower=\popsemi\fontsize{9}{11.5}\selectfont\color{popmuted},
  #1}
\newcommand{\legende}[1]{\par\vspace{6pt}{\popsemi\fontsize{9}{11.5}\selectfont\color{popmuted}#1}}

% ============================================================
% Bloc de code source — \begin{lstlisting}[language=Php]...\end{lstlisting}
% (ou language=C, HTML, JavaScript, SQL ; option omise = pseudo-code brut).
% Même recette visuelle que les figures (\popfigure) : cadre encre épais,
% fond clair (popcream), légende possible juste après via \legende{...}
% comme pour une figure (listings ne sait pas arrondir ses coins : on garde
% un cadre rectangulaire, cohérent avec les autres tableaux du document).
% CONTRAT : les modèles ne doivent utiliser QUE \begin{lstlisting}[...] tel
% quel (pas de \begin{tcblisting}, pas de \verb pour un bloc multi-lignes).
% ============================================================
\lstdefinestyle{poplst}{
  backgroundcolor=\color{popcream},
  basicstyle=\popsemi\fontsize{8.6}{11}\selectfont\color{popink},
  keywordstyle=\bfseries\color{popblue},
  commentstyle=\itshape\color{popmuted},
  stringstyle=\color{popgreen},
  numberstyle=\tiny\color{popmuted},
  identifierstyle=\color{popink},
  frame=l,
  rulecolor=\color{popink},
  framerule=1.6pt,
  framesep=7pt,
  framexleftmargin=7pt,
  framexrightmargin=7pt,
  xleftmargin=2pt,
  xrightmargin=2pt,
  breaklines=true,
  breakatwhitespace=false,
  showstringspaces=false,
  showspaces=false,
  showtabs=false,
  tabsize=2,
  columns=flexible,
  keepspaces=true,
  numbers=none,
  aboveskip=10pt,
  belowskip=10pt,
  captionpos=b,
}
\lstset{style=poplst, language=PHP}
% La distribution TeX embarquée ne fournit pas le fichier de langue
% JavaScript de listings (lstlang*.dat) : on la redéfinit nous-mêmes
% pour que \begin{lstlisting}[language=JavaScript] fonctionne.
\lstdefinelanguage{JavaScript}{
  keywords={break,case,catch,class,const,continue,debugger,default,delete,do,
    else,export,extends,finally,for,function,if,import,in,instanceof,let,new,
    return,static,super,switch,this,throw,try,typeof,var,void,while,with,yield,
    async,await,of,get,set},
  keywordstyle=\bfseries\color{popblue},
  ndkeywords={true,false,null,undefined,NaN,Infinity},
  ndkeywordstyle=\bfseries\color{poporange},
  identifierstyle=\color{popink},
  sensitive=true,
  comment=[l]{//},
  morecomment=[s]{/*}{*/},
  string=[b]",
  morestring=[b]',
  morestring=[b]`,
}
% Même limitation pour CSS et les fichiers de configuration (apacheconf,
% nginx, ini...) : ils ne sont pas fournis. CSS et un style générique de
% type "config" (commentaires # ou ;, pas de mots-clés) couvrent les besoins.
\lstdefinelanguage{CSS}{
  keywords={color,background,background-color,margin,padding,border,
    border-radius,font,font-size,font-weight,font-family,width,height,
    display,position,top,left,right,bottom,float,clear,text-align,
    line-height,list-style,cursor,overflow,box-shadow,transition,
    flex,grid,align-items,justify-content,z-index},
  keywordstyle=\bfseries\color{popblue},
  identifierstyle=\color{popink},
  sensitive=false,
  comment=[s]{/*}{*/},
  string=[b]",
  morestring=[b]',
}
\lstdefinelanguage{apacheconf}{
  sensitive=false,
  morecomment=[l]{\#},
}
\lstdefinelanguage{Apache}{
  sensitive=false,
  morecomment=[l]{\#},
}
\lstdefinelanguage{text}{sensitive=false}
\lstdefinelanguage{powershell}{
  keywords={New-Object,New-SmbShare,Get-Acl,Set-Acl,Get-ChildItem,Set-Content,
    Get-Content,Write-Host,ForEach-Object,Where-Object,Select-Object,
    Test-Path,Remove-Item,Copy-Item,Move-Item,Start-Process,Stop-Process,
    If,Else,ElseIf,ForEach,While,Function,Return,Try,Catch},
  keywordstyle=\bfseries\color{popblue},
  identifierstyle=\color{popink},
  sensitive=false,
  comment=[l]{\#},
  string=[b]",
  morestring=[b]',
}
\lstdefinelanguage{ini}{
  sensitive=false,
  morecomment=[l]{;},
  morecomment=[l]{\#},
}

% ============================================================
% Signature OnBuch+
% ============================================================
\newcommand{\poplogoinline}[1]{{\poplogo\fontsize{#1}{#1}\selectfont\color{popink}OnBuch\color{poporange}+}}
\newcommand{\signatureonbuch}{%
  \begin{tcolorbox}[enhanced,colback=popink,colframe=popink,boxrule=2.2pt,arc=10pt,
      drop shadow=poporange,left=14pt,right=14pt,top=10pt,bottom=10pt,before skip=0pt,after skip=0pt]
    \tikz[baseline=-0.7ex]\node[circle,fill=popgreen,draw=popcream,line width=1.3pt,minimum size=22pt,inner sep=0]{\popcheck{white}};%
    \hspace{9pt}\begin{minipage}[c]{0.62\linewidth}
      {\popxbold\fontsize{12}{13}\selectfont\color{popcream}Signé\ \textbullet\ {\poplogo OnBuch\color{poporange}+}}\par\vspace{2pt}
      {\raggedright\popsemi\fontsize{8}{10.5}\selectfont\color{popline}Ludovic A.\\\MakeUppercase{Conforme au programme officiel MINESEC}\par}
    \end{minipage}\hfill
    {\poplogo\fontsize{22}{22}\selectfont\color{popcream}OnBuch\color{poporange}+}
  \end{tcolorbox}%
}

% ============================================================
% Page de garde
% #1 titre · #2 matière · #3 niveau · #4 module · #5 n° leçon · #6 durée ·
% #7 description / objectifs (texte ou itemize)
% ============================================================
\newcommand{\pagegarde}[7]{%
  \thispagestyle{empty}
  \newgeometry{a4paper,margin=1.7cm,top=1.6cm,bottom=1.4cm}%
  \begin{tikzpicture}[remember picture,overlay]
    \fill[poporange!18] ([xshift=-3.2cm,yshift=-3.6cm]current page.north east) circle (4.6cm);
    \fill[poppurple!14] ([xshift=2.4cm,yshift=7.6cm]current page.south west) circle (3.8cm);
  \end{tikzpicture}%
  {\poplogo\fontsize{30}{30}\selectfont\color{popink}OnBuch\color{poporange}+}\hfill
  \raisebox{0.6ex}{\poppill{Cours signé \textbullet\ Premium}{popink}{popcream}}\par
  \vspace{1pt}\popsquiggle{poppurple}\par
  \vspace{8pt}
  \begin{tcolorbox}[enhanced,colback=poporange,colframe=popink,boxrule=2.8pt,arc=13pt,
      drop shadow=popink,left=18pt,right=18pt,top=12pt,bottom=13pt,after skip=10pt]
    \poppill{#2 \textbullet\ #3}{popink}{popcream}\hspace{5pt}\poppill{#4}{white}{popink}\par\vspace{8pt}
    {\popdisplay\fontsize{14}{14}\selectfont\color{popink}LEÇON #5}\par\vspace{4pt}
    {\raggedright\hyphenpenalty=10000\exhyphenpenalty=10000\popdisplay\fontsize{25}{27}\selectfont\color{popcream}\MakeUppercase{#1}\par}
    \vspace{8pt}
    {\popxbold\fontsize{9.5}{9.5}\selectfont\color{popink}\MakeUppercase{Durée conseillée : #6}}
  \end{tcolorbox}
  \begin{tcolorbox}[enhanced,colback=white,colframe=popink,boxrule=2.2pt,arc=10pt,
      drop shadow=popink,left=16pt,right=16pt,top=10pt,bottom=10pt,after skip=8pt]
    \poppill{Dans cette leçon}{poppurple}{white}\par\vspace{5pt}
    \popsemi\fontsize{9.3}{12.6}\selectfont\color{popink}#7
  \end{tcolorbox}
  \noindent\begin{minipage}[t]{0.487\linewidth}
    \begin{tcolorbox}[enhanced,colback=popgreen,colframe=popink,boxrule=2.2pt,arc=10pt,
        drop shadow=popink,left=11pt,right=11pt,top=10pt,bottom=10pt,height=3.1cm,valign=center]
      \tcbox[colback=white,colframe=popink,boxrule=1.3pt,arc=5pt,boxsep=0pt,nobeforeafter,
        left=3pt,right=3pt,top=3pt,bottom=3pt,valign=center]{\qrcode[height=1.9cm]{https://play.google.com/store/apps/details?id=com.luvvix.onbuch&hl=fr}}%
      \hspace{8pt}\begin{minipage}[c]{0.5\linewidth}
        {\popdisplay\fontsize{11}{12.5}\selectfont\color{white}\MakeUppercase{Télécharge OnBuch}}\par\vspace{4pt}
        {\raggedright\popsemi\fontsize{8.4}{11}\selectfont\color{white}Gratuit \textbullet\ Google Play\\Tuteur IA, annales, concours, résultats\par}
      \end{minipage}
    \end{tcolorbox}
  \end{minipage}\hfill
  \begin{minipage}[t]{0.487\linewidth}
    \begin{tcolorbox}[enhanced,colback=poppurple,colframe=popink,boxrule=2.2pt,arc=10pt,
        drop shadow=popink,left=11pt,right=11pt,top=10pt,bottom=10pt,height=3.1cm,valign=center]
      \tikz[baseline=-0.7ex]\node[circle,fill=white,draw=popink,line width=1.3pt,minimum size=1.6cm,inner sep=0]{\popdisplay\fontsize{24}{24}\selectfont\color{poppurple}+};%
      \hspace{8pt}\begin{minipage}[c]{0.6\linewidth}
        {\popdisplay\fontsize{11}{12.5}\selectfont\color{white}\MakeUppercase{Passe à OnBuch+}}\par\vspace{4pt}
        {\raggedright\popsemi\fontsize{8.4}{11}\selectfont\color{white}Tous les cours signés, Tuteur IA illimité, fiches PDF — onglet OnBuch+ de l'app\par}
      \end{minipage}
    \end{tcolorbox}
  \end{minipage}
  \vspace{8pt}
  \popcontenu
  \vfill
  \signatureonbuch
  \restoregeometry
  \newpage
}

% Rangée « Ce cours contient » (page de garde)
\newcommand{\poptile}[3]{% #1 couleur #2 gros texte #3 légende
  \begin{tcolorbox}[enhanced,colback=#1,colframe=popink,boxrule=2pt,arc=9pt,shadow={2.5pt}{-2.5pt}{0pt}{popink},
      left=6pt,right=6pt,top=9pt,bottom=9pt,halign=center,valign=center,height=1.75cm,width=\linewidth,nobeforeafter]
    {\popdisplay\fontsize{12.5}{13}\selectfont\color{white}\MakeUppercase{#2}}\par\vspace{3pt}
    {\centering\popsemi\fontsize{7.8}{9.5}\selectfont\color{white}#3\par}
  \end{tcolorbox}}
\newcommand{\popcontenu}{%
  \par\noindent{\popxboldls\fontsize{7.5}{7.5}\selectfont\color{popmuted}CE COURS SIGNÉ CONTIENT}\par\vspace{7pt}
  \noindent\begin{minipage}[t]{0.235\linewidth}\poptile{popblue}{Cours}{complet, illustré, pas à pas}\end{minipage}\hfill
  \begin{minipage}[t]{0.235\linewidth}\poptile{poporange}{Méthodes}{et exercices résolus}\end{minipage}\hfill
  \begin{minipage}[t]{0.235\linewidth}\poptile{poppink}{Exercices}{type examen + corrigés}\end{minipage}\hfill
  \begin{minipage}[t]{0.235\linewidth}\poptile{popgreen}{Fiche bilan}{l'essentiel à retenir}\end{minipage}\par}

% Sommaire
\newtcolorbox{sommaire}{enhanced,colback=white,colframe=popink,boxrule=2.2pt,arc=10pt,
  drop shadow=popink,left=18pt,right=18pt,top=13pt,bottom=13pt,before skip=6pt,after skip=20pt,
  before upper={\poppill{Sommaire}{poppurple}{white}\par\vspace{9pt}\popsemi\fontsize{10.6}{14.5}\selectfont\color{popink}}}

% Signature de fin de cours — poussée tout en bas de la dernière page
\newcommand{\findecours}{%
  \par\vfill\nopagebreak
  \begin{center}\popsquiggle{poporange}\end{center}\vspace{4pt}
  \signatureonbuch
}
\AtBeginDocument{\def\llbracket{\mathopen{[\mkern-3.5mu[}}\def\rrbracket{\mathclose{]\mkern-3.5mu]}}} % intervalles d'entiers (glyphe ⟦ absent de la police)
% Glyphes absents de Fira Math : redéfinis à partir de glyphes disponibles
\AtBeginDocument{%
  \def\setminus{\mathbin{\backslash}}\let\smallsetminus\setminus
  \def\vdots{\mathord{\vbox{\baselineskip3.2pt\lineskiplimit0pt\kern4pt\hbox{.}\hbox{.}\hbox{.}}}}%
  \def\ddots{\mathinner{\mkern1mu\raise6.5pt\hbox{.}\mkern2mu\raise3.6pt\hbox{.}\mkern2mu\raise0.7pt\hbox{.}\mkern1mu}}%
  \def\triangle{\mathord{\scalebox{0.9}{$\symup{\Delta}$}}}%
  \def\nabla{\mathord{\raisebox{\depth}{\scalebox{1}[-1]{$\symup{\Delta}$}}}}%
  \def\checkmark{\popcheck{popdark}}%
  \def\star{\mathbin{\ast}}\def\bigstar{\mathord{\ast}}%
  \def\diamond{\mathbin{\scalebox{0.6}{$\Diamond$}}}%
  \def\bigcup{\mathop{\mathchoice{\raisebox{-0.3em}{\scalebox{1.7}{$\cup$}}}{\scalebox{1.25}{$\cup$}}{\cup}{\cup}}\displaylimits}%
  \def\bigcap{\mathop{\mathchoice{\raisebox{-0.3em}{\scalebox{1.7}{$\cap$}}}{\scalebox{1.25}{$\cap$}}{\cap}{\cap}}\displaylimits}%
  \def\lesseqqgtr{\mathrel{\lessgtr}}\def\bigoplus{\mathop{\oplus}\displaylimits}%
}
\providecommand{\ssi}{si et seulement si} % abréviation fréquente
\AtBeginDocument{\providecommand{\wideparen}{\overparen}} % arc de cercle

%% FIGFIX-BEGIN v3 — correctifs globaux des figures (docs/onbuch-generation/tools/figfix)
\usepackage{adjustbox}\usepackage{etoolbox}
% toute figure TikZ/pgfplots plus large que la ligne est réduite à la largeur disponible
\BeforeBeginEnvironment{tikzpicture}{\begin{adjustbox}{max width=\linewidth}}
\AfterEndEnvironment{tikzpicture}{\end{adjustbox}}
% légendes pgfplots lisibles par-dessus les courbes
\pgfplotsset{every axis/.append style={legend style={fill=white,fill opacity=0.92,text opacity=1,draw=black!15,inner sep=3pt}}}
% étiquettes d'arêtes (syntaxe edge["..."]) avec fond blanc
\tikzset{every edge quotes/.append style={fill=white,inner sep=1.2pt}}
% bulles de cartes mentales : texte à la ligne dans la bulle, bulle un peu plus grande
\tikzset{every concept/.append style={text width=2.0cm,align=center,minimum size=2.3cm,inner sep=2pt}}
%% FIGFIX-END
\begin{document}

\pagegarde{Présentation du langage de définition des algorithmes}{Informatique}{3ème}{Informatique – classe de 3ème}{N13}{1 séance (fiche de progression harmonisée)}{Cette leçon présente le langage de définition des algorithmes (pseudo-code) et la structure de base d'un algorithme : en-tête, déclarations, corps. L'élève apprend à rédiger des algorithmes clairs et non ambigus pour résoudre des situations concrètes de la vie courante.
\vspace{4pt}
\begin{multicols}{2}\small
\begin{itemize}
\item À la fin de cette leçon, je sais définir un algorithme et le langage de définition des algorithmes (pseudo-code)
\item À la fin de cette leçon, je sais décrire la structure de base d'un algorithme (en-tête, partie déclarative, corps)
\item À la fin de cette leçon, je sais identifier et nommer correctement les variables et constantes d'un problème
\item À la fin de cette leçon, je sais distinguer les types de données élémentaires (entier, réel, caractère, chaîne, booléen)
\item À la fin de cette leçon, je sais utiliser les instructions de base : lecture, écriture, affectation
\item À la fin de cette leçon, je sais rédiger en pseudo-code un algorithme simple résolvant une situation concrète de la vie courante
\end{itemize}\end{multicols}}

\begin{sommaire}
\begin{multicols}{2}
\begin{enumerate}
\item Rappels et notion de langage de définition des algorithmes
\item Structure de base d'un algorithme
\item Variables, constantes et types de données
\item Les instructions de base : lecture, écriture, affectation
\item Rédaction complète d'algorithmes sur des situations concrètes
\item Activité d'intégration
\item Méthodes et exercices résolus
\item Exercices
\item Corrigés des exercices
\item Fiche bilan
\end{enumerate}
\end{multicols}
\end{sommaire}

\begin{objectifs}
\begin{itemize}
\item À la fin de cette leçon, je sais définir un algorithme et le langage de définition des algorithmes (pseudo-code)
\item À la fin de cette leçon, je sais décrire la structure de base d'un algorithme (en-tête, partie déclarative, corps)
\item À la fin de cette leçon, je sais identifier et nommer correctement les variables et constantes d'un problème
\item À la fin de cette leçon, je sais distinguer les types de données élémentaires (entier, réel, caractère, chaîne, booléen)
\item À la fin de cette leçon, je sais utiliser les instructions de base : lecture, écriture, affectation
\item À la fin de cette leçon, je sais rédiger en pseudo-code un algorithme simple résolvant une situation concrète de la vie courante
\item À la fin de cette leçon, je sais justifier la démarche suivie et vérifier la correction d'un algorithme par un test d'exécution à la main
\end{itemize}
\end{objectifs}

\begin{prerequis}
\begin{itemize}
\item Notion d'algorithme vue en classe de 4ème (suite ordonnée d'instructions)
\item Notion de variable et de donnée en informatique
\item Opérations arithmétiques élémentaires (addition, soustraction, multiplication, division)
\item Résolution de problèmes simples de la vie courante (calcul de moyenne, de prix, de périmètre)
\item Notion d'ordinateur et de programme (unité 1)
\end{itemize}
\end{prerequis}

\newpage

\coursec{Rappels et notion de langage de définition des algorithmes}

Mama Ngo Bell veut expliquer à son neveu, qui vit à Bafoussam, comment préparer son célèbre poulet DG. Elle peut lui décrire la recette en français, en ewondo ou en pidgin : le résultat sera le même si les étapes sont claires et dans le bon ordre. « D'abord tu découpes le poulet, ensuite tu le fais frire, puis tu ajoutes les légumes... » Peu importe la langue utilisée : ce qui compte, c'est la \emph{suite d'étapes}. De même, en informatique, un algorithme peut être décrit en français courant, mais pour éviter toute ambiguïté et se rapprocher de l'ordinateur, on utilise un langage spécial appelé \cle{langage de définition des algorithmes} (ou \cle{pseudo-code}). Comment écrire correctement un algorithme dans ce langage normalisé ? C'est l'objet de cette leçon.

\courssub{Rappel : qu'est-ce qu'un algorithme ?}

Avant d'apprendre un nouveau langage, rappelons d'abord ce que nous cherchons à exprimer avec lui.

\begin{definition}[Algorithme]
Un \cle{algorithme} est une \textbf{suite finie et ordonnée d'instructions élémentaires} permettant de résoudre un problème ou d'accomplir une tâche précise.
\end{definition}

Décomposons cette définition, car chaque mot y compte :

\begin{itemize}
\item \textbf{Suite finie} : l'algorithme doit s'arrêter après un nombre limité d'étapes. Une recette qui ne se termine jamais ne sert à rien !
\item \textbf{Ordonnée} : les instructions s'exécutent dans un ordre précis. On ne peut pas « ajouter les légumes » avant d'avoir « découpé le poulet ».
\item \textbf{Instructions élémentaires} : chaque étape doit être simple, claire et réalisable sans réflexion supplémentaire.
\item \textbf{Résoudre un problème} : un algorithme répond toujours à un besoin précis (calculer une moyenne, trouver un itinéraire, effectuer un retrait d'argent).
\end{itemize}

\begin{savaistu}[D'où vient le mot « algorithme » ?]
Le mot \emph{algorithme} vient du nom du mathématicien perse \textbf{Al-Khawarizmi}, qui vécut au IX\textsuperscript{e} siècle à Bagdad. Dans ses ouvrages, il décrivait des méthodes de calcul étape par étape, par exemple pour résoudre des équations. Son nom, latinisé en \emph{Algoritmi}, est devenu le mot que nous utilisons aujourd'hui dans le monde entier. Plus de mille ans après lui, chaque téléphone mobile, chaque guichet automatique et chaque application exécute des millions d'algorithmes chaque jour !
\end{savaistu}

\courssub{Des algorithmes dans la vie courante}

Tu exécutes des algorithmes tous les jours sans t'en rendre compte. En voici trois exemples tirés de la vie quotidienne au Cameroun.

\textbf{Exemple 1 : la recette de cuisine.} Préparer un poulet DG, c'est suivre un algorithme : découper le poulet, assaisonner, faire frire, préparer les légumes, mélanger, laisser mijoter, servir. Chaque étape est élémentaire et l'ordre est essentiel.

\textbf{Exemple 2 : l'itinéraire.} Expliquer à un ami comment aller de la gare de Bafoussam au marché central : « Sors de la gare, tourne à droite, marche jusqu'au rond-point, prends la deuxième sortie, le marché est sur ta gauche. » C'est une suite ordonnée d'instructions.

\textbf{Exemple 3 : le retrait d'argent au guichet mobile (OM/MoMo).} Quand tu retires de l'argent avec Orange Money ou MTN Mobile Money, le téléphone exécute un algorithme :

\begin{lstlisting}
Composer le code du service (#150# ou *126#)
Choisir "Retrait d'argent"
Saisir le montant a retirer
Saisir le code secret
Confirmer l'operation
Recevoir le message de confirmation
\end{lstlisting}

Remarque que si tu saisis un mauvais code secret, l'opération est refusée : l'algorithme prévoit même les cas d'erreur !

\begin{aretenir}[L'essentiel]
Un algorithme n'est pas forcément lié à l'ordinateur : toute méthode précise, décrite étape par étape, pour accomplir une tâche est un algorithme. L'ordinateur n'est qu'une machine capable d'exécuter des algorithmes très rapidement et sans se tromper.
\end{aretenir}

\courssub{Les limites du langage naturel}

Si l'on peut décrire un algorithme en français courant, pourquoi inventer un autre langage ? Parce que le langage naturel (le français de tous les jours) présente de sérieux défauts lorsqu'il s'agit de donner des instructions à une machine.

\begin{itemize}
\item \textbf{L'ambiguïté} : une même phrase peut avoir plusieurs sens. Si l'on écrit « ajoute le sel et le poivre à la fin », faut-il ajouter le poivre à la fin seulement, ou les deux ? Un humain devine ; un ordinateur, non.
\item \textbf{La longueur} : décrire en phrases complètes un calcul simple prend beaucoup de place et de temps. « Tu prends la somme des deux notes, tu la divises par deux, et tu obtiens la moyenne » est bien plus long que « moyenne $\leftarrow$ (note1 + note2) / 2 ».
\item \textbf{L'imprécision} : des mots comme « un peu », « environ », « bientôt » n'ont pas de signification exacte pour une machine. L'ordinateur a besoin de valeurs et d'actions parfaitement définies.
\end{itemize}

\begin{exemplebox}[Une consigne ambiguë]
Le père de famille dit à son fils : « Va acheter du pain, et s'il y a des œufs, prends-en six. » Le fils revient avec six pains ! Pourquoi ? Parce que la phrase est ambiguë : « en » peut se rapporter aux œufs (le sens voulu) ou au pain (le sens compris). Un algorithme bien écrit lèverait cette ambiguïté : \texttt{SI il y a des œufs ALORS acheter 6 œufs}.
\end{exemplebox}

\courssub{Le langage de définition des algorithmes (pseudo-code)}

Pour décrire un algorithme de façon claire, courte et sans ambiguïté, les informaticiens ont créé un langage intermédiaire, à mi-chemin entre le français et les langages de programmation.

\begin{definition}[Langage de définition des algorithmes]
Le \cle{langage de définition des algorithmes} (LDA), appelé aussi \cle{pseudo-code}, est un \textbf{langage conventionnel, proche du français}, utilisé pour décrire un algorithme de manière \textbf{précise et non ambiguë}, indépendamment de tout langage de programmation.
\end{definition}

Le pseudo-code possède trois caractéristiques essentielles :

\begin{itemize}
\item \textbf{Des mots-clés réservés} : certains mots ont un sens fixe et ne peuvent pas être utilisés autrement. Par exemple : \texttt{Algorithme}, \texttt{Début}, \texttt{Fin}, \texttt{Lire}, \texttt{Écrire}, \texttt{Si}, \texttt{Alors}. Ces mots sont « réservés » : on ne peut pas s'en servir pour nommer autre chose.
\item \textbf{Une syntaxe précise} : chaque instruction obéit à une forme exacte. On écrit \texttt{Lire(note)} et non « lis un peu la note si tu veux ». Cette rigueur supprime toute ambiguïté.
\item \textbf{L'indépendance vis-à-vis des langages de programmation} : le pseudo-code n'appartient à aucun langage (ni C, ni JavaScript, ni Python). Un même algorithme écrit en pseudo-code pourra ensuite être traduit dans n'importe quel langage de programmation, exactement comme la recette de Mama Ngo Bell peut être racontée en français, en ewondo ou en pidgin.
\end{itemize}

\begin{exemplebox}[Premier algorithme en pseudo-code]
Voici comment on écrit, en pseudo-code, l'algorithme qui calcule la moyenne de deux notes d'un élève de 3\textsuperscript{e} :
\begin{lstlisting}
Algorithme Moyenne
Variable note1, note2, moyenne : reel
Debut
    Ecrire("Entrez la premiere note : ")
    Lire(note1)
    Ecrire("Entrez la deuxieme note : ")
    Lire(note2)
    moyenne <- (note1 + note2) / 2
    Ecrire("La moyenne est : ", moyenne)
Fin
\end{lstlisting}
On reconnaît les mots-clés réservés \texttt{Algorithme}, \texttt{Variable}, \texttt{Debut}, \texttt{Fin}, \texttt{Lire} et \texttt{Ecrire}. Nous étudierons chacun d'eux en détail dans les sections suivantes de cette leçon.
\end{exemplebox}

\courssub{De l'idée à la machine : la chaîne de traduction}

L'ordinateur ne comprend ni le français ni le pseudo-code : il ne comprend que le \cle{langage machine}, fait de 0 et de 1. Pour passer d'un problème concret à son exécution par l'ordinateur, on suit donc une chaîne en quatre étapes :

\begin{enumerate}
\item \textbf{Le problème} : on part d'une situation de la vie réelle (calculer les moyennes d'une classe, gérer les ventes d'une boutique).
\item \textbf{L'algorithme} : par \emph{analyse} du problème, on conçoit la suite d'étapes qui le résout, et on la rédige en pseudo-code.
\item \textbf{Le programme} : on \emph{traduit} l'algorithme dans un langage de programmation (C, JavaScript, Python...). C'est le travail du programmeur.
\item \textbf{Le résultat} : l'ordinateur \emph{exécute} le programme et produit le résultat attendu.
\end{enumerate}

\begin{popfigure}
\begin{center}
\begin{tikzpicture}[
    box/.style={draw=popblue, thick, fill=popblueL, rounded corners=3pt,
                minimum width=2.6cm, minimum height=1.1cm, align=center, font=\bfseries},
    arr/.style={-{Stealth[length=3mm]}, very thick, poporange},
    lbl/.style={font=\small\itshape, popdark}
]
\node[box] (pb) at (0,0) {Problème};
\node[box, align=center] (algo) at (4.2,0){Algorithme\\ \mdseries\small (pseudo-code)};
\node[box, align=center] (prog) at (8.4,0){Programme\\ \mdseries\small (langage)};
\node[box] (res) at (12.6,0) {Résultat};

\draw[arr] (pb) -- (algo) node[midway, above, lbl]{analyse};
\draw[arr] (algo) -- (prog) node[midway, above, lbl]{traduction};
\draw[arr] (prog) -- (res) node[midway, above, lbl]{exécution};
\end{tikzpicture}
\end{center}
\legende{La chaîne de traduction : du problème réel au résultat produit par l'ordinateur.}
\end{popfigure}

\begin{attention}[Ne confonds pas algorithme et programme !]
C'est l'erreur la plus fréquente en classe de 3\textsuperscript{e}. L'\textbf{algorithme} est la \emph{méthode de résolution}, écrite en pseudo-code, indépendante de toute machine. Le \textbf{programme} est la \emph{traduction} de cet algorithme dans un langage de programmation précis (C, JavaScript...), exécutable par l'ordinateur. Retiens l'image : l'algorithme est la recette écrite sur le cahier ; le programme est le plat cuisiné dans une cuisine précise, avec des ustensiles précis.
\end{attention}

\begin{exoresolu}[Identifier les étapes de la chaîne de traduction]
Énoncé. Le proviseur d'un collège de Douala veut un outil qui affiche, pour chaque élève, la mention obtenue au BEPC à partir de sa moyenne. Un informaticien rédige d'abord la méthode en pseudo-code, puis l'écrit en langage C, et l'ordinateur affiche finalement « Mention : Bien » pour une moyenne de 14. Pour chaque élément suivant, indique à quelle étape de la chaîne de traduction il correspond : (a) la moyenne de l'élève et sa mention ; (b) le texte en pseudo-code ; (c) le code en langage C ; (d) le besoin du proviseur.
\tcblower
Solution détaillée. On applique la chaîne : problème $\rightarrow$ algorithme $\rightarrow$ programme $\rightarrow$ résultat.
\begin{itemize}
\item (a) La moyenne et la mention affichée constituent le \textbf{résultat}, produit par l'exécution du programme.
\item (b) Le texte en pseudo-code est \textbf{l'algorithme}, obtenu par l'analyse du problème.
\item (c) Le code en langage C est \textbf{le programme}, obtenu par traduction de l'algorithme.
\item (d) Le besoin du proviseur est \textbf{le problème} de départ, la situation concrète à résoudre.
\end{itemize}
Conclusion : chaque étape de la chaîne a un rôle précis, et aucune ne peut être sautée : sans analyse, pas d'algorithme ; sans algorithme, pas de programme fiable.
\end{exoresolu}

\begin{exercice}[À toi de jouer]
1. Donne deux exemples d'algorithmes de la vie courante, différents de ceux du cours, en précisant les étapes dans l'ordre.
2. Pourquoi le français courant est-il mal adapté pour donner des instructions à un ordinateur ? Cite deux raisons.
3. Quelle est la différence entre un algorithme et un programme ?
4. Remets dans l'ordre les étapes de la chaîne de traduction : programme, problème, résultat, algorithme.
\end{exercice}

\begin{corrige}[Exercice — À toi de jouer]
1. Exemples possibles : la lessive (trier le linge, tremper, savonner, rincer, essorer, étendre) ; la recharge de crédit téléphonique (composer le code, saisir le numéro de la carte, valider, recevoir la confirmation).
2. Le français courant est \textbf{ambigu} (une phrase peut avoir plusieurs sens) et \textbf{imprécis} (mots vagues comme « un peu »), alors qu'un ordinateur exige des instructions exactes. Il est aussi trop long.
3. L'algorithme est la méthode de résolution écrite en pseudo-code ; le programme est sa traduction dans un langage de programmation, exécutable par l'ordinateur.
4. Ordre correct : problème $\rightarrow$ algorithme $\rightarrow$ programme $\rightarrow$ résultat.
\end{corrige}

\begin{aretenir}[Bilan de la section]
Un algorithme est une suite finie et ordonnée d'instructions élémentaires qui résout un problème. Pour le décrire sans ambiguïté, on utilise le \textbf{langage de définition des algorithmes} (pseudo-code), un langage conventionnel proche du français, doté de mots-clés réservés et d'une syntaxe précise, indépendant de tout langage de programmation. La chaîne de traduction mène du problème au résultat : analyse, pseudo-code, programme, exécution. Dans la prochaine section, nous découvrirons la structure de base d'un algorithme écrit en pseudo-code.
\end{aretenir}

\coursec{Structure de base d'un algorithme}

Dans la section précédente, nous avons vu qu'un algorithme est une suite finie et ordonnée d'instructions permettant de résoudre un problème, et que pour le rédiger sans ambiguïté, on utilise un \cle{langage de définition des algorithmes}, appelé aussi \cle{pseudo-code}. Mais une question se pose immédiatement : comment ce texte doit-il être organisé ? Peut-on écrire les instructions n'importe comment, dans le désordre ? Non. Tout comme une lettre administrative comporte toujours un en-tête, un corps et une signature, un algorithme suit une \cle{structure} précise et immuable. C'est cette structure que nous allons découvrir et maîtriser dans cette section.

\begin{definition}[Langage de définition des algorithmes (Pseudo-code)]
Un \cle{langage de définition des algorithmes} (ou \cle{pseudo-code}) est un langage normalisé, indépendant de toute machine, qui permet d'écrire des algorithmes avec une syntaxe précise, rigoureuse et compréhensible par tout informaticien. Il utilise des mots-clés en français (\texttt{Algorithme}, \texttt{Var}, \texttt{Debut}, \texttt{Fin}, \texttt{Si}, \texttt{TantQue}, \ldots) pour décrire la logique du traitement sans se soucier des détails syntaxiques d'un langage de programmation réel (C, Python, JavaScript\ldots).
\end{definition}

\courssub{Les trois parties d'un algorithme}

\textbf{Intuition : à quoi sert une structure ?}

Imagine que tu demandes à un ami de préparer du ndolé. Si tu lui donnes les étapes dans le désordre (« mets le sel, lave les feuilles, achète les arachides, sers le plat »), le résultat sera un désastre. De même, l'ordinateur (ou la personne qui lit ton algorithme) a besoin de savoir \emph{avant de commencer} : de quoi on parle (le nom du problème), quels ingrédients on va utiliser (les variables), et seulement ensuite ce qu'on fait avec (les instructions).

\begin{propriete}[Structure générale d'un algorithme]
Tout algorithme écrit en pseudo-code comporte \textbf{trois parties}, toujours dans cet ordre :
\begin{enumerate}
    \item \textbf{L'en-tête} : il donne le nom de l'algorithme, introduit par le mot-clé \texttt{Algorithme} ;
    \item \textbf{La partie déclarative} : elle annonce toutes les variables et constantes utilisées, avec leurs types, après les mots-clés \texttt{Var} et \texttt{Const} ;
    \item \textbf{Le corps} : il contient les instructions à exécuter, délimité par les mots-clés \texttt{Debut} et \texttt{Fin}.
\end{enumerate}
Cette propriété est admise : elle constitue la convention du langage de définition des algorithmes utilisé au collège. Nous la justifierons par l'exemple tout au long de la leçon.
\end{propriete}

\begin{popfigure}
\begin{center}
\begin{tikzpicture}[font=\small]
% Bloc en-tête
\draw[fill=popblueL, draw=popblue, thick, rounded corners=2pt] (0,4.2) rectangle (7,5.2);
\node[align=center] at (3.5,4.7) {\textbf{En-tête} : \texttt{Algorithme Somme}};
% Bloc déclarations
\draw[fill=poporangeL, draw=poporange, thick, rounded corners=2pt] (0,2.2) rectangle (7,4.0);
\node[align=center] at (3.5,3.4) {\textbf{Partie déclarative} : \texttt{Var} / \texttt{Const}};
\node[align=center] at (3.5,2.8) {annonce des variables et constantes avec leurs types};
% Bloc corps
\draw[fill=popgreenL, draw=popgreen, thick, rounded corners=2pt] (0,0) rectangle (7,2.0);
\node[align=center] at (3.5,1.5) {\textbf{Corps} : \texttt{Debut} \ldots\ \texttt{Fin}};
\node[align=center] at (3.5,0.9) {suite d'instructions (lire, écrire, calculer)};
\node[align=center] at (3.5,0.35) {exécutées dans l'ordre};
% Annotations latérales
\node[align=left, text=popblue, anchor=west] at (7.4,4.7) {Qui ? Le nom\\de l'algorithme};
\draw[-{Stealth}, popblue, thick] (7.1,4.7) -- (7.35,4.7);
\node[align=left, text=poporange, anchor=west] at (7.4,3.1) {Avec quoi ? Les\\objets manipulés};
\draw[-{Stealth}, poporange, thick] (7.1,3.1) -- (7.35,3.1);
\node[align=left, text=popgreen, anchor=west] at (7.4,1.0) {Quoi faire ? Les\\actions à exécuter};
\draw[-{Stealth}, popgreen, thick] (7.1,1.0) -- (7.35,1.0);
\end{tikzpicture}
\end{center}
\legende{Les trois parties d'un algorithme : l'en-tête (le nom), la partie déclarative (les objets utilisés) et le corps (les actions).}
\end{popfigure}

\courssub{L'en-tête : nommer son algorithme}

L'\cle{en-tête} est la première ligne de l'algorithme. Sa syntaxe est très simple :

\begin{propriete}[Syntaxe de l'en-tête]
L'en-tête s'écrit :
\begin{center}
\texttt{Algorithme} \emph{nom\_de\_l\_algorithme}
\end{center}
Le mot-clé \texttt{Algorithme} est obligatoire et s'écrit avec la majuscule. Il est suivi du \cle{nom} de l'algorithme, qui est un \cle{identificateur}.
\end{propriete}

Un \cle{identificateur} est un nom choisi par le programmeur pour désigner un algorithme, une variable ou une constante. On ne peut pas choisir n'importe quel nom : il faut respecter des règles strictes.

\begin{methode}[Bien choisir un identificateur]
Pour nommer un identificateur (algorithme, variable, constante), respecte les règles suivantes :
\begin{enumerate}
    \item Il \textbf{commence obligatoirement par une lettre} (jamais par un chiffre) ;
    \item Il ne contient \textbf{ni espace, ni accent, ni caractère spécial} (pas de \texttt{é}, \texttt{ç}, \texttt{-}, \texttt{?}, \texttt{!}\ldots) ; seul le tiret bas \texttt{\_} est toléré ;
    \item Il est \textbf{court et significatif} : il doit rappeler ce que l'objet représente.
\end{enumerate}
\textbf{Exemples corrects} : \texttt{Somme}, \texttt{moyenne}, \texttt{prixUnitaire}, \texttt{age}, \texttt{note1}.\\
\textbf{Exemples incorrects} : \texttt{1note} (commence par un chiffre), \texttt{prix unitaire} (espace), \texttt{moyenné} (accent), \texttt{nbre-eleves} (tiret).
\end{methode}

\begin{attention}[Erreurs fréquentes sur les identificateurs]
Un nom comme \texttt{x} ou \texttt{truc} est syntaxiquement correct, mais il ne veut rien dire : trois jours plus tard, tu ne sauras plus à quoi il sert ! Préfère \texttt{moyenneClasse} à \texttt{m}. À l'inverse, évite les noms trop longs comme \texttt{lamoyennedeselevesdelaclassedetroisieme}. L'équilibre : court \emph{et} clair.
\end{attention}

\courssub{La partie déclarative : annoncer les objets utilisés}

\textbf{Intuition.} Avant de cuisiner, on sort tous les ingrédients sur la table. Avant d'exécuter les instructions, l'algorithme doit \emph{annoncer} tous les objets qu'il va manipuler : c'est le rôle de la \cle{partie déclarative}.

\begin{definition}[Partie déclarative]
La partie déclarative est la zone de l'algorithme où l'on \textbf{déclare} (annonce) toutes les \cle{variables} et toutes les \cle{constantes} qui seront utilisées dans le corps, en précisant leur \cle{type} (entier, réel, caractère, chaîne, booléen). Elle est introduite par les mots-clés \texttt{Var} (pour les variables) et \texttt{Const} (pour les constantes).
\end{definition}

\begin{propriete}[Syntaxe des déclarations]
\begin{itemize}
    \item Variables : \texttt{Var} \emph{nom1, nom2, \ldots} \texttt{:} \emph{type}
    \item Constantes : \texttt{Const} \emph{nom} \texttt{=} \emph{valeur}
\end{itemize}
On peut regrouper sur une même ligne plusieurs variables de même type, séparées par des virgules. Une constante reçoit sa valeur dès sa déclaration et cette valeur ne changera jamais pendant l'exécution. Les types de base sont : \texttt{entier}, \texttt{reel}, \texttt{caractere}, \texttt{chaine}, \texttt{booleen}.
\end{propriete}

\begin{exemplebox}[Une partie déclarative]
\begin{lstlisting}
Var a, b, s : entier
Var moyenne : reel
Var nomEleve : chaine
Const TAUX = 18
\end{lstlisting}
Ici, on annonce trois variables entières (\texttt{a}, \texttt{b}, \texttt{s}), une variable réelle (\texttt{moyenne}), une chaîne de caractères (\texttt{nomEleve}) et une constante \texttt{TAUX} qui vaudra toujours 18 (par exemple un taux de TVA en pourcentage).
\end{exemplebox}

\courssub{Le corps : les instructions entre Debut et Fin}

Le \cle{corps} de l'algorithme est la partie « active » : c'est là que se trouvent les \cle{instructions} (lire des données, calculer, afficher des résultats). Ces instructions seront étudiées en détail dans la section 4 ; ici, retenons surtout leur cadre.

\begin{propriete}[Syntaxe du corps]
Le corps est \textbf{obligatoirement délimité} par les mots-clés \texttt{Debut} et \texttt{Fin} :
\begin{itemize}
    \item \texttt{Debut} marque le point de départ de l'exécution ;
    \item \texttt{Fin} marque l'arrêt de l'algorithme.
\end{itemize}
Les instructions situées entre ces deux mots-clés sont exécutées \textbf{dans l'ordre où elles sont écrites}, de haut en bas.
\end{propriete}

\begin{attention}[Les deux erreurs classiques du débutant]
\begin{itemize}
    \item \textbf{Utiliser une variable non déclarée} : si tu écris \texttt{s $\leftarrow$ a + b} sans avoir déclaré \texttt{s} dans la partie \texttt{Var}, l'algorithme est faux. Toute variable utilisée dans le corps doit avoir été annoncée avant.
    \item \textbf{Oublier \texttt{Debut} ou \texttt{Fin}} : sans ces mots-clés, on ne sait plus où commence et où s'arrête l'algorithme. Vérifie toujours que chaque \texttt{Debut} a bien son \texttt{Fin}.
\end{itemize}
\end{attention}

\courssub{Commentaires et indentation : rendre l'algorithme lisible}

Un algorithme n'est pas écrit seulement pour l'ordinateur : il est écrit pour être \textbf{lu par des humains} (toi, ton camarade, ton professeur). Deux outils rendent un algorithme lisible.

\begin{definition}[Commentaire]
Un \cle{commentaire} est un texte libre placé dans l'algorithme pour expliquer ce que fait une ligne ou un bloc d'instructions. Il est \textbf{ignoré lors de l'exécution}. On l'introduit par \texttt{//} : tout ce qui suit sur la ligne est un commentaire.
\end{definition}

\begin{definition}[Indentation]
L'\cle{indentation} est le décalage vers la droite (avec des espaces) des instructions situées à l'intérieur d'un bloc. Elle permet de voir d'un coup d'œil ce qui se trouve entre \texttt{Debut} et \texttt{Fin}.
\end{definition}

Compare les deux versions suivantes : elles font exactement la même chose, mais laquelle préfères-tu lire ?

\begin{lstlisting}
Algorithme Somme
Var a, b, s : entier
Debut
Lire(a, b)
s <- a + b
Ecrire(s)
Fin
\end{lstlisting}

\begin{lstlisting}
Algorithme Somme
// Calcule et affiche la somme de deux nombres
Var a, b, s : entier   // les deux nombres et leur somme
Debut
    Lire(a, b)          // saisie des deux valeurs
    s <- a + b          // calcul de la somme
    Ecrire(s)           // affichage du resultat
Fin
\end{lstlisting}

La seconde version, indentée et commentée, est incomparablement plus claire. Prends dès maintenant l'habitude d'indenter et de commenter : c'est un \cle{savoir-être} du bon programmeur.

\courssub{Exemple complet commenté : la somme de deux nombres}

\begin{exemplebox}[Algorithme Somme, complet et commenté]
\begin{lstlisting}
Algorithme Somme
// Role : lit deux nombres entiers et affiche leur somme
Var a, b, s : entier   // a et b : nombres saisis ; s : somme
Debut
    Lire(a, b)          // l'utilisateur entre les deux valeurs
    s <- a + b          // on range le resultat de a + b dans s
    Ecrire(s)           // on affiche le contenu de s
Fin
\end{lstlisting}
On retrouve bien les trois parties : l'en-tête (\texttt{Algorithme Somme}), la partie déclarative (\texttt{Var a, b, s : entier}) et le corps (entre \texttt{Debut} et \texttt{Fin}).
\end{exemplebox}

\begin{exoresolu}[Identifier les trois parties d'un algorithme]
\textbf{Énoncé.} Voici un algorithme écrit par un élève de 3ème :
\begin{lstlisting}
Algorithme Double
Var n, d : entier
Debut
    Lire(n)
    d <- 2 * n
    Ecrire(d)
Fin
\end{lstlisting}
1) Indique l'en-tête, la partie déclarative et le corps. 2) Que fait cet algorithme si l'utilisateur saisit la valeur 7 ?
\tcblower
\textbf{Solution.}
1) L'\textbf{en-tête} est la ligne \texttt{Algorithme Double}. La \textbf{partie déclarative} est la ligne \texttt{Var n, d : entier} : elle annonce deux variables entières. Le \textbf{corps} est l'ensemble des lignes comprises entre \texttt{Debut} et \texttt{Fin}, soit les trois instructions \texttt{Lire(n)}, \texttt{d <- 2 * n} et \texttt{Ecrire(d)}.

2) Traçons l'exécution pas à pas avec $n = 7$ :
\begin{itemize}
    \item \texttt{Lire(n)} : la variable $n$ reçoit la valeur $7$ ;
    \item \texttt{d <- 2 * n} : la variable $d$ reçoit $2 \times 7 = 14$ ;
    \item \texttt{Ecrire(d)} : l'algorithme affiche $14$.
\end{itemize}
Cet algorithme calcule et affiche donc le \textbf{double} du nombre saisi, ce qui justifie le nom bien choisi de l'identificateur \texttt{Double}.
\end{exoresolu}

\begin{exercice}[À toi de jouer]
Un élève veut écrire un algorithme qui calcule le périmètre d'un carré à partir de la longueur de son côté. Il écrit :
\begin{lstlisting}
Algorithme perimetre-carre
Debut
    Lire(c)
    p <- 4 * c
    Ecrire(p)
Fin
\end{lstlisting}
Trouve les \textbf{trois erreurs} commises (indice : regarde l'identificateur, la partie déclarative et les mots-clés), puis réécris l'algorithme correctement, avec un commentaire et une indentation soignée.
\end{exercice}

\begin{corrige}[Exercice 1]
Les trois erreurs sont :
\begin{enumerate}
    \item L'identificateur \texttt{perimetre-carre} est incorrect : il contient un tiret. Il faut écrire par exemple \texttt{PerimetreCarre} (pas d'accent sur Perimetre non plus).
    \item La \textbf{partie déclarative} est absente : les variables \texttt{c} et \texttt{p} ne sont pas déclarées.
    \item L'absence de déclaration rend le corps invalide ; on doit ajouter \texttt{Var c, p : reel} avant \texttt{Debut}.
\end{enumerate}
Version corrigée :
\begin{lstlisting}
Algorithme PerimetreCarre
// Calcule le perimetre d'un carre a partir de son cote
Var c, p : reel        // c : cote du carre ; p : perimetre
Debut
    Lire(c)             // saisie du cote
    p <- 4 * c          // calcul du perimetre
    Ecrire(p)           // affichage du resultat
Fin
\end{lstlisting}
\end{corrige}

\begin{aretenir}[L'essentiel de la section]
\begin{itemize}
    \item Tout algorithme comporte \textbf{trois parties} : l'\textbf{en-tête} (\texttt{Algorithme} + nom), la \textbf{partie déclarative} (\texttt{Var} / \texttt{Const} + types) et le \textbf{corps} (entre \texttt{Debut} et \texttt{Fin}).
    \item Un \textbf{identificateur} commence par une lettre, sans espace, ni accent, ni caractère spécial ; il doit être court et significatif.
    \item Toute variable utilisée dans le corps doit avoir été \textbf{déclarée} au préalable.
    \item Les \textbf{commentaires} (\texttt{//}) et l'\textbf{indentation} rendent l'algorithme lisible sans modifier son exécution.
\end{itemize}
\end{aretenir}

Maintenant que la charpente est bien en place, il reste à donner vie au corps de l'algorithme : comment lire une donnée, afficher un résultat, ranger une valeur dans une variable ? C'est l'objet des sections suivantes, qui détaillent les instructions de base.

\coursec{Variables, constantes et types de données}

Dans la section précédente, nous avons découvert la structure de base d'un algorithme : un en-tête qui donne son nom, une partie \textbf{déclarative} où l'on annonce les objets utilisés, puis un corps délimité par \texttt{Début} et \texttt{Fin} contenant les instructions. Mais que déclare-t-on exactement dans cette partie déclarative ? Pour le comprendre, posons-nous une question simple.

Imagine que le proviseur de ton collège te demande d'écrire un algorithme qui calcule la moyenne d'un élève à partir de ses notes. Pour faire ce calcul, l'algorithme doit \emph{retenir} quelque part les notes, la somme, puis la moyenne. Or un ordinateur ne retient rien « dans sa tête » : il stocke les informations dans sa \cle{mémoire}. Il faut donc réserver des emplacements dans cette mémoire, leur donner des noms pour pouvoir les retrouver, et préciser quelle sorte d'information chacun contiendra. C'est exactement le rôle des \cle{variables}, des \cle{constantes} et des \cle{types de données}, que nous allons étudier dans cette section.

\courssub{La notion de variable}

\textbf{Intuition : la boîte étiquetée.}

Pense à une boîte sur laquelle tu colles une étiquette, par exemple « age ». À l'intérieur de la boîte, tu places une fiche sur laquelle est écrit un nombre : 14. L'année suivante, tu peux ouvrir la boîte, retirer la fiche et en mettre une nouvelle portant le nombre 15. La boîte garde le même nom (« age »), mais son \emph{contenu} a changé. Une variable en algorithmique fonctionne exactement comme cette boîte étiquetée : c'est un emplacement de la mémoire de l'ordinateur, identifié par un nom, et dont le contenu peut être modifié au cours de l'exécution de l'algorithme.

\begin{definition}[Variable]
Une \cle{variable} est un emplacement de la mémoire de l'ordinateur, identifié par un \textbf{nom} (appelé \emph{identificateur}) et contenant une \textbf{valeur} qui peut \textbf{varier} (changer) au cours de l'exécution de l'algorithme.
\end{definition}

Une variable possède donc trois caractéristiques essentielles, qu'il faut toujours avoir en tête :

\begin{itemize}
\item son \textbf{nom} (ou identificateur) : c'est l'étiquette qui permet de désigner l'emplacement mémoire, par exemple \texttt{age}, \texttt{moyenne}, \texttt{nomEleve} ;
\item son \textbf{type} : c'est la nature des valeurs qu'elle peut contenir (nombre entier, nombre décimal, texte, etc.) ;
\item sa \textbf{valeur} : c'est le contenu actuel de l'emplacement mémoire, qui peut changer pendant l'exécution.
\end{itemize}

\begin{propriete}[Règles de choix d'un nom de variable]
Le nom d'une variable (identificateur) doit respecter les règles suivantes :
\begin{itemize}
\item il commence obligatoirement par une \textbf{lettre} ;
\item il ne contient ni espace, ni accent, ni caractère spécial (seuls les lettres, les chiffres et parfois le tiret bas \texttt{\_} sont admis) ;
\item il doit être \textbf{différent des mots réservés} du langage (\texttt{Algorithme}, \texttt{Var}, \texttt{Début}, \texttt{Fin}, etc.) ;
\item il doit être \textbf{explicite} : on préfère \texttt{prixHT} à \texttt{x}, car un bon nom rend l'algorithme plus facile à lire et à comprendre.
\end{itemize}
\end{propriete}

\begin{attention}[Noms de variables interdits ou déconseillés]
Les déclarations suivantes sont \textbf{incorrectes} : \texttt{2note} (commence par un chiffre), \texttt{prix TTC} (contient un espace), \texttt{prénom} (contient un accent), \texttt{Début} (mot réservé du langage). Choisis toujours des noms simples, sans espace ni accent, et qui décrivent clairement le contenu : \texttt{note2}, \texttt{prixTTC}, \texttt{prenom}.
\end{attention}

\courssub{La notion de constante}

\textbf{Intuition : ce qui ne change jamais.}

Dans certaines situations, une valeur reste \emph{toujours la même} pendant tout le déroulement du programme. Par exemple, au Cameroun, le taux de la Taxe sur la Valeur Ajoutée (TVA) est fixé à 19,25 \%. Si tu écris un algorithme qui calcule le prix TTC des articles vendus dans la boutique de ta famille, ce taux sera le même pour tous les calculs : il n'a aucune raison de changer pendant l'exécution. Plutôt que de réécrire le nombre 19,25 partout (avec le risque de se tromper une fois sur dix), on le range une seule fois dans une \cle{constante}.

\begin{definition}[Constante]
Une \cle{constante} est un emplacement de la mémoire de l'ordinateur, identifié par un nom, dont la valeur est fixée au départ et \textbf{ne peut pas changer} pendant l'exécution de l'algorithme.
\end{definition}

L'intérêt des constantes est double :

\begin{itemize}
\item \textbf{lisibilité} : écrire \texttt{TAUX} est plus parlant qu'écrire \texttt{19,25} au milieu d'un calcul ;
\item \textbf{facilité de modification} : si le taux de TVA change un jour, il suffit de modifier \emph{une seule ligne} (la déclaration de la constante) au lieu de chercher le nombre dans tout l'algorithme.
\end{itemize}

\begin{exemplebox}[Une constante en situation réelle]
Pour calculer le prix toutes taxes comprises (TTC) d'un article à partir de son prix hors taxes (HT), on déclare :
\begin{lstlisting}
Const TAUX = 19,25
Var prixHT, prixTTC : reel
\end{lstlisting}
Le calcul s'écrira ensuite : \texttt{prixTTC $\leftarrow$ prixHT + prixHT * TAUX / 100}. La valeur de \texttt{TAUX} reste fixe pendant toute l'exécution, tandis que \texttt{prixHT} et \texttt{prixTTC} peuvent changer d'un article à l'autre.
\end{exemplebox}

\begin{aretenir}[Variable ou constante ?]
Si la valeur d'une donnée \textbf{peut changer} pendant l'exécution (une note saisie, un total qui s'accumule, un âge), on utilise une \textbf{variable}. Si la valeur est \textbf{fixée une fois pour toutes} (taux de TVA, nombre de jours d'une semaine, valeur de $\pi$), on utilise une \textbf{constante}. Par convention, on écrit souvent les noms des constantes en \textbf{majuscules} (\texttt{TAUX}, \texttt{PI}, \texttt{NB\_JOURS}) pour les distinguer facilement des variables.
\end{aretenir}

\courssub{Syntaxe de déclaration}

Avant d'utiliser une variable ou une constante dans le corps d'un algorithme, il faut obligatoirement l'\textbf{annoncer} dans la partie déclarative : c'est ce qu'on appelle la \cle{déclaration}. Déclarer, c'est dire à l'ordinateur : « réserve-moi une boîte, voici son nom, et voici la nature de ce qu'elle contiendra ».

\begin{propriete}[Syntaxe de déclaration]
La déclaration d'une \textbf{variable} s'écrit :
\begin{center}
\texttt{Var nom\_de\_la\_variable : type}
\end{center}
On peut déclarer plusieurs variables de même type sur une seule ligne, en séparant leurs noms par des virgules :
\begin{center}
\texttt{Var note1, note2, note3 : reel}
\end{center}
La déclaration d'une \textbf{constante} s'écrit :
\begin{center}
\texttt{Const NOM\_DE\_LA\_CONSTANTE = valeur}
\end{center}
Le type d'une constante n'est pas précisé : il est déterminé automatiquement par la valeur qu'on lui donne.
\end{propriete}

\begin{exemplebox}[Partie déclarative complète]
Voici la partie déclarative d'un algorithme qui gère les informations d'un élève :
\begin{lstlisting}
Algorithme Bulletin
Const TAUX = 19,25
Const NB_TRIMESTRES = 3
Var nom : chaine
Var age : entier
Var moyenne, total : reel
Var admis : booleen
Debut
    (* instructions *)
Fin
\end{lstlisting}
Remarque l'ordre : on déclare d'abord les constantes (\texttt{Const}), puis les variables (\texttt{Var}), avant le mot \texttt{Debut}.
\end{exemplebox}

\courssub{Les types élémentaires de données}

Le \cle{type} d'une variable indique la \textbf{nature} des valeurs qu'elle peut contenir. C'est une information indispensable pour l'ordinateur : selon le type, il ne réserve pas la même place en mémoire et n'autorise pas les mêmes opérations (on peut additionner deux nombres, mais pas deux prénoms !). En classe de 3\textsuperscript{e}, nous utilisons cinq \textbf{types élémentaires}.

\begin{definition}[Types élémentaires]
Les cinq types élémentaires du langage algorithmique sont :
\begin{itemize}
\item \texttt{entier} : un nombre sans partie décimale, positif ou négatif. Exemples : $15$, $-3$, $2024$.
\item \texttt{reel} : un nombre qui peut avoir une partie décimale. Exemples : $3,14$, $12,5$, $-0,75$.
\item \texttt{caractere} : un seul caractère (lettre, chiffre ou symbole), écrit entre apostrophes. Exemples : \texttt{'A'}, \texttt{'b'}, \texttt{'7'}, \texttt{'?'}.
\item \texttt{chaine} : une suite de caractères (un texte), écrite entre guillemets. Exemples : \texttt{"Douala"}, \texttt{"Amina"}, \texttt{"3e C"}.
\item \texttt{booleen} : une valeur de vérité, qui ne peut être que \texttt{Vrai} ou \texttt{Faux}. Sert à représenter une réponse de type oui/non.
\end{itemize}
\end{definition}

Le tableau suivant récapitule ces types avec des exemples tirés de la vie scolaire.

\begin{center}
\begin{tabularx}{\linewidth}{|l|X|X|}
\hline
\rowcolor{popblueL} \textbf{Type} & \textbf{Exemples de valeurs} & \textbf{Exemple d'utilisation} \\
\hline
\texttt{entier} & 15 ; 2024 ; $-3$ & un âge, un nombre d'élèves, une année \\
\hline
\texttt{reel} & 3,14 ; 12,5 ; 19,25 & un prix, une moyenne, une taille en mètres \\
\hline
\texttt{caractere} & 'A' ; 'g' ; '7' & une seule lettre, la mention d'un bulletin \\
\hline
\texttt{chaine} & "Douala" ; "Amina" & un nom, une adresse, une classe \\
\hline
\texttt{booleen} & Vrai ; Faux & « l'élève est-il admis ? », « le stock est-il vide ? » \\
\hline
\end{tabularx}
\end{center}

\begin{attention}[Apostrophes ou guillemets ?]
Ne confonds pas le type \texttt{caractere} et le type \texttt{chaine} : un caractère seul s'écrit entre \textbf{apostrophes} (\texttt{'A'}), tandis qu'une chaîne s'écrit entre \textbf{guillemets} (\texttt{"Amina"}). Attention aussi : \texttt{'7'} est un \emph{caractère} (le symbole 7), \texttt{"7"} est une \emph{chaîne}, et $7$ est un \emph{entier} : ce sont trois données de types différents, avec lesquelles on ne peut pas faire les mêmes opérations.
\end{attention}

Pour bien visualiser ce qui se passe en mémoire après les déclarations et les premières affectations, représentons nos « boîtes étiquetées » :

\begin{popfigure}
\begin{center}
\begin{tikzpicture}[box/.style={draw=popblue, thick, rounded corners=2pt, minimum width=3.2cm, minimum height=1.1cm, fill=popblueL},
tag/.style={draw=poporange, thick, fill=poporangeL, rounded corners=2pt, font=\small\bfseries, minimum height=0.7cm},
typ/.style={font=\small\itshape, text=popdark}]
% Variable age
\node[tag] (t1) at (0,1.15) {age};
\node[box] (b1) at (0,0) {\Large 14};
\node[typ] at (0,-0.95) {type : entier};
% Variable nom
\node[tag] (t2) at (4.6,1.15) {nom};
\node[box] (b2) at (4.6,0) {\Large "Amina"};
\node[typ] at (4.6,-0.95) {type : chaine};
% Variable moyenne
\node[tag] (t3) at (9.2,1.15) {moyenne};
\node[box] (b3) at (9.2,0) {\Large 12,5};
\node[typ] at (9.2,-0.95) {type : reel};
\end{tikzpicture}
\end{center}
\legende{Représentation de trois variables en mémoire : chaque « boîte » possède un nom (étiquette orange), une valeur (contenu) et un type.}
\end{popfigure}

Chaque boîte correspond à un emplacement de la mémoire. L'étiquette orange est le nom de la variable, le contenu de la boîte est sa valeur actuelle, et le type précise la nature de ce contenu. Pendant l'exécution, on pourra remplacer le contenu de la boîte \texttt{age} par 15, celui de la boîte \texttt{moyenne} par 13,75, etc. En revanche, on ne pourra jamais mettre un texte dans la boîte \texttt{age} : son type \texttt{entier} l'interdit.

\courssub{Choisir le bon type pour une donnée}

Face à une situation concrète, la première question à se poser avant de déclarer une variable est : « \emph{quelles valeurs cette donnée peut-elle prendre ?} ». La réponse détermine le type.

\begin{methode}[Comment choisir le type d'une donnée]
Pour choisir correctement le type d'une variable, procède par étapes :
\begin{enumerate}
\item \textbf{Identifie la donnée} à mémoriser (par exemple : l'âge d'un élève).
\item \textbf{Demande-toi quelles valeurs elle peut prendre} :
\begin{itemize}
\item un nombre sans virgule (14, 15, 16\ldots) $\rightarrow$ type \texttt{entier} ;
\item un nombre qui peut avoir une virgule (12,5 ; 999,99\ldots) $\rightarrow$ type \texttt{reel} ;
\item une seule lettre ou un seul symbole ('A', 'F'\ldots) $\rightarrow$ type \texttt{caractere} ;
\item un texte de plusieurs caractères ("Amina", "Yaoundé"\ldots) $\rightarrow$ type \texttt{chaine} ;
\item une réponse de type oui/non $\rightarrow$ type \texttt{booleen}.
\end{itemize}
\item \textbf{Vérifie les cas limites} : la valeur peut-elle un jour devenir décimale ? Si oui, choisis \texttt{reel} même si les premiers exemples sont des entiers.
\item \textbf{Déclare la variable} avec le type retenu : \texttt{Var age : entier}.
\end{enumerate}
\end{methode}

Appliquons cette méthode à quelques données de la vie courante :

\begin{center}
\begin{tabularx}{\linewidth}{|X|l|X|}
\hline
\rowcolor{popblueL} \textbf{Donnée} & \textbf{Type choisi} & \textbf{Justification} \\
\hline
L'âge d'un élève & \texttt{entier} & on exprime l'âge en années entières : 14, 15\ldots \\
\hline
Le prix d'un cahier & \texttt{reel} & un prix peut être décimal : 350,50 FCFA \\
\hline
Le nom d'un élève & \texttt{chaine} & c'est un texte de plusieurs caractères \\
\hline
La moyenne générale & \texttt{reel} & une moyenne est rarement un nombre entier : 12,5 \\
\hline
Le sexe (M ou F) & \texttt{caractere} & une seule lettre : 'M' ou 'F' \\
\hline
L'élève est admis ? & \texttt{booleen} & la réponse est Vrai ou Faux \\
\hline
\end{tabularx}
\end{center}

\begin{attention}[Erreur fréquente : le prix ou la moyenne déclarés en entier]
Une erreur très classique consiste à déclarer un prix ou une moyenne comme \texttt{entier} parce que « dans mon exemple, c'est un nombre rond ». Si tu déclares \texttt{Var moyenne : entier} et que le calcul donne 12,5, la partie décimale sera \textbf{perdue} : l'algorithme affichera 12 au lieu de 12,5, ce qui peut changer la décision d'admission d'un élève ! Règle pratique : dès qu'une valeur \emph{peut} être décimale (prix, moyenne, taille, taux), déclare-la toujours en \texttt{reel}.
\end{attention}

\begin{exoresolu}[Déclarer les données d'une cantine scolaire]
Énoncé : la cantine du collège vend des repas. On veut écrire un algorithme qui, à partir du prix hors taxes d'un repas, calcule le prix TTC (taux de TVA : 19,25 \%) et indique si le repas est payé ou non. Écris la partie déclarative complète de cet algorithme en justifiant le type de chaque donnée.
\tcblower
Solution détaillée : identifions d'abord toutes les données du problème.
\begin{itemize}
\item Le \textbf{taux de TVA} : il est fixé à 19,25 et ne change jamais pendant l'exécution $\rightarrow$ \textbf{constante} \texttt{TAUX = 19,25}.
\item Le \textbf{prix hors taxes} du repas : c'est un nombre qui peut être décimal (par exemple 500 FCFA ou 750,50 FCFA) et qui varie d'un repas à l'autre $\rightarrow$ variable de type \texttt{reel}.
\item Le \textbf{prix TTC} : résultat d'un calcul avec un taux décimal, il est forcément susceptible d'être décimal $\rightarrow$ variable de type \texttt{reel}.
\item L'information « le repas est-il payé ? » : la réponse est oui ou non $\rightarrow$ variable de type \texttt{booleen}.
\end{itemize}
La partie déclarative s'écrit donc :
\begin{lstlisting}
Algorithme Cantine
Const TAUX = 19,25
Var prixHT, prixTTC : reel
Var estPaye : booleen
Debut
    (* instructions de calcul *)
Fin
\end{lstlisting}
On remarque que le taux est déclaré en constante (valeur fixe), que les deux prix sont déclarés en \texttt{reel} (et non en \texttt{entier}, pour ne pas perdre les centimes), et que l'état du paiement est un booléen puisqu'il ne vaut que Vrai ou Faux.
\end{exoresolu}

\begin{exercice}[À toi de jouer]
Le secrétariat de ton collège veut un algorithme pour gérer la fiche d'un élève. Pour chacune des données suivantes, indique s'il s'agit d'une variable ou d'une constante, donne-lui un nom correct et précise son type en justifiant : le nom de l'élève ; sa date de naissance (année seulement) ; sa moyenne du trimestre ; le nombre de classes du collège (fixé à 12) ; l'information « l'élève a payé la scolarité » ; la première lettre de son quartier.
\end{exercice}

\begin{corrige}[Exercice 1]
\begin{itemize}
\item Nom de l'élève : variable \texttt{nom} de type \texttt{chaine} (texte de plusieurs caractères, différent pour chaque élève).
\item Année de naissance : variable \texttt{anneeNaissance} de type \texttt{entier} (un nombre entier comme 2011).
\item Moyenne du trimestre : variable \texttt{moyenne} de type \texttt{reel} (peut être décimale : 12,5).
\item Nombre de classes : constante \texttt{NB\_CLASSES = 12} (valeur fixe, non précisée car le type est déduit de la valeur).
\item Scolarité payée : variable \texttt{scolaritePayee} de type \texttt{booleen} (Vrai ou Faux).
\item Première lettre du quartier : variable \texttt{initiale} de type \texttt{caractere} (une seule lettre, par exemple 'B' pour Bonabéri).
\end{itemize}
\end{corrige}

\begin{aretenir}[L'essentiel de la section]
\begin{itemize}
\item Une \textbf{variable} est un emplacement mémoire nommé dont la valeur \textbf{peut changer} pendant l'exécution ; une \textbf{constante} est un emplacement nommé dont la valeur \textbf{reste fixe}.
\item On déclare une variable par \texttt{Var nom : type} et une constante par \texttt{Const NOM = valeur}, dans la partie déclarative, avant \texttt{Debut}.
\item Les cinq types élémentaires sont : \texttt{entier}, \texttt{reel}, \texttt{caractere} (entre apostrophes), \texttt{chaine} (entre guillemets) et \texttt{booleen} (Vrai/Faux).
\item Pour choisir un type, on se demande quelles valeurs la donnée peut prendre, en pensant aux cas limites : toute valeur susceptible d'être décimale (prix, moyenne) se déclare en \texttt{reel}.
\end{itemize}
\end{aretenir}

Maintenant que nous savons réserver et nommer nos « boîtes mémoire » avec le bon type, il reste à apprendre à les utiliser : comment mettre une valeur dans une variable, comment demander une information à l'utilisateur et comment afficher un résultat. C'est précisément l'objet de la section suivante, consacrée aux instructions de base : la lecture, l'écriture et l'affectation.

\coursec{Structure de base d'un algorithme}

\courssub{Variables, constantes et types de données}

Avant d'écrire le moindre calcul, un algorithme doit \emph{nommer} et \emph{typer} les données qu'il va manipuler. C'est le rôle de la partie \textbf{déclarative} (avant le \texttt{Début}).

\begin{definition}[Variable]
Une \cle{variable} est une \textbf{zone mémoire nommée} dont la valeur \textbf{peut changer} au cours de l'exécution de l'algorithme. Elle est caractérisée par :
\begin{itemize}
\item un \textbf{nom} (identificateur) : commence par une lettre, sans espace, sans accent (ex. \texttt{total}, \texttt{nbBeignets}, \texttt{moyenneClasse}) ;
\item un \textbf{type} : l'ensemble des valeurs qu'elle peut contenir ;
\item une \textbf{valeur actuelle} : le contenu de la case mémoire à un instant donné.
\end{itemize}
\end{definition}

\begin{definition}[Constante]
Une \cle{constante} est une \textbf{zone mémoire nommée} dont la valeur \textbf{ne change jamais} pendant l'exécution. Elle est fixée dès la déclaration.
\begin{center}
\texttt{Constante <nom> = <valeur> : <Type>}
\end{center}
Exemple : \texttt{Constante PRIX\_BEIGNET = 100 : Entier}
\end{definition}

\begin{propriete}[Types de données de base (niveau 3ème)]
\begin{center}
\begin{tabularx}{\linewidth}{|l|X|l|}\hline
\rowcolor{popblueL} \textbf{Type} & \textbf{Ensemble des valeurs} & \textbf{Exemples littéraux} \\ \hline
\texttt{Entier} & $\mathbb{Z} = \{\dots, -2, -1, 0, 1, 2, \dots\}$ & \texttt{5}, \texttt{-3}, \texttt{0} \\ \hline
\texttt{Réel} & $\mathbb{R}$ (nombres à virgule) & \texttt{3,14}, \texttt{-0,5}, \texttt{12,0} \\ \hline
\texttt{Caractère} & Un seul symbole (lettre, chiffre, signe) & \texttt{'A'}, \texttt{'7'}, \texttt{'+'} \\ \hline
\texttt{Booléen} & \{\texttt{Vrai}, \texttt{Faux}\} & \texttt{Vrai}, \texttt{Faux} \\ \hline
\end{tabularx}
\end{center}
\emph{Note :} En algorithmique au BEPC, on écrit les réels avec une \textbf{virgule} (notation française), pas un point.
\end{propriete}

\begin{propriete}[Syntaxe de la déclaration]
\begin{itemize}
\item \textbf{Variables} : \texttt{Variable <liste noms> : <Type>} (séparés par des virgules).
\item \textbf{Constantes} : \texttt{Constante <nom> = <valeur> : <Type>} (une par ligne).
\end{itemize}
La zone de déclaration se place \textbf{avant} le mot-clé \texttt{Début}.
\end{propriete}

\begin{exemplebox}[Déclarations variées]
\begin{lstlisting}
Algorithme ExempleDeclarations
Constante TVA = 19,25 : Reel
Constante NOM_ETAB = "Lycee Bilingue" : Chaine
Variable age, note : Entier
Variable moyenne, prix : Reel
Variable reponse : Caractere
Variable admis : Booleen
Debut
    // ... corps de l'algorithme
Fin
\end{lstlisting}
\emph{Remarque :} Le type \texttt{Chaine} (chaîne de caractères) est souvent accepté pour les messages longs, bien que le type de base soit \texttt{Caractere}.
\end{exemplebox}

\begin{attention}[Règles de nommage (identificateurs)]
\begin{itemize}
\item Commencer par une \textbf{lettre} (pas de chiffre, pas de \texttt{\_} au début).
\item Pas d'espace, pas d'accent (\texttt{moyenne\_classe} autorisé, \texttt{moyenne classe} interdit, \texttt{âge} interdit).
\item Respecter la \textbf{casse} : \texttt{total} et \texttt{Total} sont deux variables différentes (sensible à la casse dans la plupart des langages, et recommandé en pseudo-code).
\item Éviter les mots réservés du langage (\texttt{Si}, \texttt{Alors}, \texttt{Variable}, \texttt{Entier}, \dots).
\end{itemize}
\end{attention}

\courssub{Les instructions de base : lecture, écriture, affectation}

Maintenant que nos données sont déclarées et typées, il faut les \emph{faire vivre}. Trois instructions fondamentales permettent l'interaction et le calcul : \cle{Lire} (entrée), \cle{Écrire} (sortie) et l'\cle{affectation} (calcul/mémorisation).

Prenons une situation concrète : la cantine du collège vend le beignet à \num{100} FCFA. On veut un algorithme qui demande le nombre de beignets achetés, calcule le prix total, puis l'affiche. Il faut donc : \textbf{lire} une donnée, \textbf{calculer} (affecter), puis \textbf{écrire} le résultat.

\courssub{L'instruction d'entrée : la lecture}

\courssub{L'instruction de sortie : l'écriture}

Commençons par les deux instructions qui permettent à l'algorithme de \emph{communiquer} avec l'utilisateur.

\begin{definition}[Lecture et écriture]
\begin{itemize}
\item L'instruction \cle{Lire} (ou instruction d'\textbf{entrée}) permet à l'utilisateur de \textbf{fournir une valeur} au clavier ; cette valeur est rangée dans une variable \emph{préalablement déclarée} et de type compatible.
\item L'instruction \cle{Écrire} (ou instruction de \textbf{sortie}) permet d'\textbf{afficher} à l'écran un message (texte entre guillemets doubles) ou la valeur d'une variable / d'une expression. On peut enchaîner plusieurs arguments séparés par des virgules.
\end{itemize}
\end{definition}

\begin{propriete}[Syntaxe exacte]
\begin{itemize}
\item \texttt{Lire(variable)} : l'algorithme s'arrête, attend que l'utilisateur tape une valeur valide pour le type de la variable, puis la range dedans.
\item \texttt{Écrire("message")} : affiche le texte tel quel (sans les guillemets).
\item \texttt{Écrire(variable)} ou \texttt{Écrire(expression)} : affiche la \emph{valeur} de la variable ou le résultat du calcul.
\item \textbf{Combinaison} : \texttt{Écrire("Le total est : ", total, " FCFA")}.
\end{itemize}
\end{propriete}

\begin{exemplebox}[Premier algorithme complet : la cantine]
\begin{lstlisting}
Algorithme Cantine
Constante PRIX_UNITAIRE = 100 : Entier
Variable nb, total : Entier
Debut
    Ecrire("Combien de beignets ? ")
    Lire(nb)                       // l'utilisateur tape par exemple 5
    total <- nb * PRIX_UNITAIRE    // calcul du prix total
    Ecrire("Le prix total est : ", total, " FCFA")
Fin
\end{lstlisting}
Si l'utilisateur tape \texttt{5}, l'écran affiche : \texttt{Le prix total est : 500 FCFA}.
\end{exemplebox}

\begin{attention}[Guillemets ou pas guillemets ?]
\texttt{Écrire("total")} affiche le mot \emph{total}, alors que \texttt{Écrire(total)} affiche la \emph{valeur} de la variable (par exemple 500). C'est une erreur très fréquente au BEPC : le texte entre guillemets est affiché tel quel, sans guillemets c'est le contenu de la variable qui est affiché.
\end{attention}

\courssub{L'instruction d'affectation}

Voici l'instruction la plus importante de toute l'algorithmique : c'est elle qui permet de \emph{calculer} et de \emph{mémoriser} les résultats.

\begin{definition}[Affectation]
L'\cle{affectation} (symbole $\leftarrow$, flèche vers la gauche, saisi souvent \texttt{<-} au clavier) consiste à \textbf{donner à une variable la valeur d'une expression}. Sa syntaxe est :
\begin{center}
\texttt{variable $\leftarrow$ expression}
\end{center}
On lit : « la variable \emph{reçoit} la valeur de l'expression » ou « la variable \emph{prend la valeur de} l'expression ».
\end{definition}

\begin{propriete}[Sens de l'affectation -- Règle d'or]
L'affectation s'exécute toujours en \textbf{deux étapes, dans cet ordre strict} :
\begin{enumerate}
\item On \textbf{évalue d'abord l'expression à droite} de la flèche (on effectue le calcul avec les valeurs \emph{actuelles} des variables) ;
\item Puis on \textbf{range le résultat dans la variable à gauche}, en \emph{écrasant} son ancienne valeur.
\end{enumerate}
La variable de gauche doit être de type \emph{compatible} avec le résultat de l'expression.
\end{propriete}

\begin{exemplebox}[L'ancienne valeur est écrasée]
\begin{lstlisting}
a <- 5         // a contient 5
a <- a + 3     // 1. on évalue a + 3 = 5 + 3 = 8
               // 2. on range 8 dans a (l'ancien 5 est perdu)
Ecrire(a)      // affiche 8
\end{lstlisting}
Après ces deux affectations, la variable \texttt{a} contient \textbf{8}. L'instruction \texttt{a <- a + 3} se lit : « le \emph{nouveau} \texttt{a} vaut l'\emph{ancien} \texttt{a} plus 3 ».
\end{exemplebox}

\begin{attention}[Deux erreurs classiques à ne jamais commettre]
\begin{itemize}
\item \textbf{Sens interdit (côté gauche)} : écrire \texttt{a + b <- s} est \emph{faux}. À gauche de la flèche, il doit y avoir \textbf{une seule variable} (un nom de case mémoire), jamais un calcul ni une constante. C'est la variable qui reçoit, pas l'expression !
\item \textbf{Confusion avec les mathématiques} : en maths, $a = a + 3$ est une équation impossible (aucun nombre n'est égal à lui-même plus 3). En algorithmique, \texttt{a <- a + 3} est parfaitement correct : ce n'est \emph{pas} une égalité, c'est un \emph{transfert} d'une nouvelle valeur. Le symbole $\leftarrow$ n'a \textbf{rien à voir} avec le signe $=$ des mathématiques.
\end{itemize}
\end{attention}

\courssub{Les opérateurs}

Les expressions à droite de la flèche (ou dans les conditions, ou dans \texttt{Écrire}) sont construites avec des \cle{opérateurs}.

\begin{propriete}[Opérateurs arithmétiques (sur Entiers et Réels)]
\begin{center}
\begin{tabularx}{\linewidth}{|l|X|l|}\hline
\rowcolor{popblueL} \textbf{Opérateur} & \textbf{Signification} & \textbf{Exemple (résultat)} \\ \hline
\texttt{+} & addition & \texttt{7 + 2} donne 9 \\ \hline
\texttt{-} & soustraction & \texttt{7 - 2} donne 5 \\ \hline
\texttt{*} & multiplication & \texttt{7 * 2} donne 14 \\ \hline
\texttt{/} & division réelle (résultat Réel) & \texttt{7 / 2} donne 3,5 \\ \hline
\texttt{div} & quotient de la division euclidienne (Entiers) & \texttt{7 div 2} donne 3 \\ \hline
\texttt{mod} & reste de la division euclidienne (Entiers) & \texttt{7 mod 2} donne 1 \\ \hline
\end{tabularx}
\end{center}
\emph{Priorités :} \texttt{*}, \texttt{/}, \texttt{div}, \texttt{mod} prioritaires sur \texttt{+}, \texttt{-}. Parenthèses \texttt{( )} pour forcer l'ordre.
\end{propriete}

\begin{propriete}[Opérateurs de comparaison (résultat : Booléen)]
Ils comparent deux valeurs de même type (Entier, Réel, Caractère) et renvoient \texttt{Vrai} ou \texttt{Faux}.
\begin{center}
\begin{tabularx}{\linewidth}{|c|X|}\hline
\rowcolor{popblueL} \textbf{Symbole} & \textbf{Sens} \\ \hline
\texttt{=} & égal à \\ \hline
$\neq$ (\texttt{<>} parfois toléré) & différent de \\ \hline
\texttt{<} & strictement inférieur à \\ \hline
\texttt{>} & strictement supérieur à \\ \hline
$\leq$ & inférieur ou égal à \\ \hline
$\geq$ & supérieur ou égal à \\ \hline
\end{tabularx}
\end{center}
Exemple : \texttt{age >= 15} vaut \texttt{Vrai} si \texttt{age} contient 17.
\end{propriete}

\begin{exemplebox}[Bons de cantine : \texttt{div} et \texttt{mod} en action]
Un élève a 750 FCFA. Combien de beignets à 100 FCFA peut-il acheter, et combien lui restera-t-il ?
\begin{lstlisting}
nb <- 750 div 100    // nb contient 7 (quotient entier)
reste <- 750 mod 100 // reste contient 50 (reste entier)
Ecrire("Nombre : ", nb, " Reste : ", reste)
\end{lstlisting}
Il achète 7 beignets et il lui reste 50 FCFA. Vérification : $7 \times 100 + 50 = 750$.
\end{exemplebox}

\courssub{Rédaction complète d'algorithmes et vérification par trace}

Comment être sûr qu'un algorithme est correct \emph{avant} de le faire exécuter par un ordinateur (ou de le rendre au BEPC) ? En jouant soi-même le rôle de la machine, à la main : c'est la \cle{trace d'exécution} (ou \cle{tableau de trace}).

\begin{methode}[Construire un tableau de trace]
Pour vérifier un algorithme, on construit un \textbf{tableau de trace} :
\begin{enumerate}
\item On crée une colonne pour le \textbf{numéro/texte de l'instruction} et une colonne \textbf{par variable} de l'algorithme (dans l'ordre de déclaration) ;
\item On exécute les instructions \textbf{dans l'ordre séquentiel}, une ligne par instruction exécutée ;
\item Après chaque instruction, on note la \textbf{nouvelle valeur} de chaque variable modifiée (une variable non modifiée garde sa valeur : on recopie ou on laisse vide) ;
\item À la fin, on vérifie que les valeurs obtenues correspondent au résultat attendu (calcul mental, logique métier).
\end{enumerate}
\end{methode}

\courssub{Exemple guidé : échanger le contenu de deux variables}

\textbf{Problème.} Deux variables \texttt{a} et \texttt{b} contiennent respectivement 5 et 9. On veut \emph{échanger} leurs contenus : à la fin, \texttt{a} doit contenir 9 et \texttt{b} doit contenir 5. C'est comme échanger le contenu de deux verres pleins : il faut un troisième verre vide !

\begin{attention}[Pourquoi pas simplement \texttt{a <- b} puis \texttt{b <- a} ?]
Si l'on écrit \texttt{a <- b} d'abord, la valeur 9 écrase le 5 contenu dans \texttt{a} : le 5 est \textbf{perdu à jamais}. Ensuite \texttt{b <- a} recopie 9 dans \texttt{b}. Résultat : les deux variables contiennent 9 ! Échanger deux variables \textbf{sans variable intermédiaire écrase l'une des valeurs}.
\end{attention}

\begin{methode}[L'échange avec une variable temporaire]
\begin{enumerate}
\item Sauvegarder la valeur de \texttt{a} dans une variable \texttt{temp} (ou \texttt{aux}, \texttt{tampon}) ;
\item Recopier la valeur de \texttt{b} dans \texttt{a} ;
\item Recopier la valeur sauvegardée (\texttt{temp}) dans \texttt{b}.
\end{enumerate}
\end{methode}

\begin{exemplebox}[Algorithme d'échange complet]
\begin{lstlisting}
Algorithme Echange
Variable a, b, temp : Entier
Debut
    a <- 5
    b <- 9
    temp <- a     // 1. Sauvegarde : temp vaut 5
    a <- b        // 2. a reçoit 9
    b <- temp     // 3. b reçoit l'ancien a (5)
    Ecrire("Apres echange : a = ", a, " b = ", b)
Fin
\end{lstlisting}
\end{exemplebox}

\begin{popfigure}
\begin{center}
\begin{tikzpicture}[x=2.8cm, y=0.65cm, font=\small]
% Grille
\draw[thick, popblue] (0,0) rectangle (4,-9);
\foreach \x in {1,2,3} \draw[popblue] (\x,0) -- (\x,-9);
\foreach \y in {-1,...,-8} \draw[popline] (0,\y) -- (4,\y);
% En-têtes
\fill[popblueL] (0,0) rectangle (4,-1);
\node[font=\bfseries, text=popdark] at (0.5,-0.5) {Instruction exécutée};
\node[font=\bfseries, text=popdark] at (1.5,-0.5) {a};
\node[font=\bfseries, text=popdark] at (2.5,-0.5) {b};
\node[font=\bfseries, text=popdark] at (3.5,-0.5) {temp};
% Lignes de trace
% Ligne 1 : Debut (valeurs initiales inconnues ?)
\node[align=center] at (0.5,-1.5) {Début\\(init.)};
\node at (1.5,-1.5) {?};
\node at (2.5,-1.5) {?};
\node at (3.5,-1.5) {?};
% Ligne 2 : a <- 5
\node at (0.5,-2.5) {\texttt{a <- 5}};
\node[popgreen, font=\bfseries] at (1.5,-2.5) {5};
\node at (2.5,-2.5) {?};
\node at (3.5,-2.5) {?};
% Ligne 3 : b <- 9
\node at (0.5,-3.5) {\texttt{b <- 9}};
\node at (1.5,-3.5) {5};
\node[popgreen, font=\bfseries] at (2.5,-3.5) {9};
\node at (3.5,-3.5) {?};
% Ligne 4 : temp <- a
\node at (0.5,-4.5) {\texttt{temp <- a}};
\node at (1.5,-4.5) {5};
\node at (2.5,-4.5) {9};
\node[popgreen, font=\bfseries] at (3.5,-4.5) {5};
% Ligne 5 : a <- b
\node at (0.5,-5.5) {\texttt{a <- b}};
\node[poporange, font=\bfseries] at (1.5,-5.5) {9};
\node at (2.5,-5.5) {9};
\node at (3.5,-5.5) {5};
% Ligne 6 : b <- temp
\node at (0.5,-6.5) {\texttt{b <- temp}};
\node at (1.5,-6.5) {9};
\node[popgreen, font=\bfseries] at (2.5,-6.5) {5};
\node at (3.5,-6.5) {5};
% Ligne 7 : Ecrire
\node at (0.5,-7.5) {\texttt{Ecrire(...)}};
\node at (1.5,-7.5) {9};
\node at (2.5,-7.5) {5};
\node at (3.5,-7.5) {5};
% Ligne 8 : Fin
\node[font=\itshape] at (0.5,-8.5) {Fin};
\node[font=\bfseries] at (1.5,-8.5) {9};
\node[font=\bfseries] at (2.5,-8.5) {5};
\node at (3.5,-8.5) {5};
\end{tikzpicture}
\end{center}
\legende{Tableau de trace de l'algorithme d'échange. Les valeurs modifiées sont surlignées (vert = initialisation/sauvegarde, orange = modification cible). À la fin, \texttt{a} et \texttt{b} ont bien échangé leurs contenus grâce à \texttt{temp}.}
\end{popfigure}

\begin{experience}[TP guidé : Trace manuelle d'un calcul de moyenne]
\textbf{Objectif :} S'entraîner à construire une trace complète sur papier.
\textbf{Matériel :} Feuille, stylo.
\textbf{Algorithme à tracer :}
\begin{lstlisting}
Algorithme MoyenneDeuxNotes
Variable n1, n2, somme, moy : Reel
Debut
    Ecrire("Note 1 : "); Lire(n1)
    Ecrire("Note 2 : "); Lire(n2)
    somme <- n1 + n2
    moy <- somme / 2
    Ecrire("Moyenne = ", moy)
Fin
\end{lstlisting}
\textbf{Manipulation :}
\begin{enumerate}
\item Dessiner le tableau (colonnes : Instruction, n1, n2, somme, moy).
\item Simuler l'exécution pour \texttt{n1 = 12} et \texttt{n2 = 16}.
\item Remplir le tableau ligne par ligne.
\item Vérifier que \texttt{moy} vaut bien 14 à la fin.
\end{enumerate}
\textbf{Interprétation :} La trace prouve que l'algorithme calcule correctement la moyenne arithmétique. C'est la méthode de validation officielle au BEPC.
\end{experience}

\begin{exoresolu}[Prix de trois cahiers et monnaie]
Énoncé. Un cahier coûte 350 FCFA. Écrire un algorithme qui demande le nombre de cahiers achetés, calcule le prix total, puis affiche le prix total et la monnaie à rendre si l'élève paie avec un billet de 2000 FCFA. Donner la trace pour 4 cahiers.
\tcblower
Solution détaillée.
\begin{lstlisting}
Algorithme AchatCahiers
Constante PRIX_CAHIER = 350 : Entier
Constante BILLET = 2000 : Entier
Variable nb, total, monnaie : Entier
Debut
    Ecrire("Nombre de cahiers ? ")
    Lire(nb)                    // L'utilisateur tape 4
    total <- nb * PRIX_CAHIER   // total = 4 * 350 = 1400
    monnaie <- BILLET - total   // monnaie = 2000 - 1400 = 600
    Ecrire("Total a payer : ", total, " FCFA")
    Ecrire("Monnaie a rendre : ", monnaie, " FCFA")
Fin
\end{lstlisting}
\textbf{Trace d'exécution (pour nb = 4) :}
\begin{center}
\begin{tabularx}{0.9\linewidth}{|l|X|X|X|}\hline
\rowcolor{popblueL} \textbf{Instruction} & \textbf{nb} & \textbf{total} & \textbf{monnaie} \\ \hline
Début & ? & ? & ? \\ \hline
\texttt{Lire(nb)} & 4 & ? & ? \\ \hline
\texttt{total <- nb * 350} & 4 & 1400 & ? \\ \hline
\texttt{monnaie <- 2000 - total} & 4 & 1400 & 600 \\ \hline
\texttt{Ecrire(...)} & 4 & 1400 & 600 \\ \hline
\end{tabularx}
\end{center}
L'écran affiche : \texttt{Total a payer : 1400 FCFA} puis \texttt{Monnaie a rendre : 600 FCFA}. Le résultat est correct : $4 \times 350 = 1400$ et $2000 - 1400 = 600$.
\end{exoresolu}

\begin{savaistu}[Le mot « algorithme » et Al-Khwarizmi]
Le mot \emph{algorithme} vient du nom du mathématicien perse \textbf{Al-Khwarizmi} (vers 780--850), qui a travaillé à la « Maison de la Sagesse » de Bagdad. Son livre \emph{Al-Kitab al-mukhtasar fi hisab al-jabr wa'l-muqabala} a donné le mot « algèbre » et a systématisé les méthodes de calcul étape par étape (les algorithmes) pour résoudre les équations. Aujourd'hui, chaque fois que vous suivez une recette de cuisine, un mode d'emploi de montage de meuble, ou que vous faites une division posée, vous exécutez un algorithme !
\end{savaistu}

\begin{aretenir}[L'essentiel de la leçon -- Structure de base]
\begin{itemize}
\item \textbf{Déclaration} : \texttt{Variable <noms> : <Type>} | \texttt{Constante <nom> = <val> : <Type>} (avant \texttt{Début}).
\item \textbf{Types de base} : \texttt{Entier}, \texttt{Reel}, \texttt{Caractere}, \texttt{Booleen}.
\item \texttt{Lire(variable)} : entrée clavier $\rightarrow$ variable.
\item \texttt{Écrire(...)} : sortie écran (texte "..." ou valeurs).
\item \texttt{variable <- expression} : \textbf{1. Calculer à droite}, \textbf{2. Ranger à gauche} (écrase l'ancien).
\item \textbf{Opérateurs} : \texttt{+ - * / div mod} (arithmétiques) ; \texttt{= <> < > <= >=} (comparatifs, renvoient Booléen).
\item \textbf{Trace d'exécution} : tableau (Instruction $\times$ Variables) pour valider la logique pas à pas.
\item \textbf{Échange de deux variables} : \textbf{obligatoirement} via une variable temporaire (\texttt{temp}).
\end{itemize}
\end{aretenir}

\begin{exercice}[À toi de jouer -- Entraînement BEPC]
\begin{enumerate}
\item \textbf{Calcul successif.} Que contient la variable \texttt{x} après la séquence : \texttt{x <- 10} ; \texttt{x <- x - 4} ; \texttt{x <- x * 2} ?
\item \textbf{Algorithme complet.} Écrire un algorithme qui lit le prix unitaire d'un stylo (entier) et la quantité achetée (entier), puis affiche le montant total à payer (avec libellé « Montant : ... FCFA »).
\item \textbf{Trace table.} Construire le tableau de trace complet de l'algorithme suivant pour les valeurs initiales \texttt{a = 12} et \texttt{b = 5} :
\begin{lstlisting}
q <- a div b
r <- a mod b
Ecrire(q, r)
\end{lstlisting}
\item \textbf{Erreur ou pas ?} Dire si les instructions suivantes sont syntaxiquement correctes (selon nos conventions). Si non, corriger.
\begin{itemize}
\item \texttt{5 <- x}
\item \texttt{total <- "100"}
\item \texttt{Lire("Entrez age : ", age)}
\item \texttt{Ecrire(age >= 18)}
\end{itemize}
\end{enumerate}
\end{exercice}

\begin{corrige}[Exercices -- Corrigé détaillé]
\begin{enumerate}
\item \texttt{x <- 10} $\Rightarrow$ \texttt{x} = 10. \quad \texttt{x <- x - 4} $\Rightarrow$ évalue 10-4=6, range 6 dans \texttt{x}. \quad \texttt{x <- x * 2} $\Rightarrow$ évalue 6*2=12, range 12 dans \texttt{x}. \textbf{Réponse finale : 12}.
\item \begin{lstlisting}
Algorithme AchatStylos
Variable prix, qte, montant : Entier
Debut
    Ecrire("Prix unitaire du stylo ? "); Lire(prix)
    Ecrire("Quantite ? "); Lire(qte)
    montant <- prix * qte
    Ecrire("Montant : ", montant, " FCFA")
Fin
\end{lstlisting}
\item \textbf{Trace :}
\begin{center}
\begin{tabularx}{0.7\linewidth}{|l|X|X|}\hline
\rowcolor{popblueL} \textbf{Instruction} & \textbf{q} & \textbf{r} \\ \hline
Début & ? & ? \\ \hline
\texttt{q <- a div b} (12 div 5) & 2 & ? \\ \hline
\texttt{r <- a mod b} (12 mod 5) & 2 & 2 \\ \hline
\texttt{Ecrire(q, r)} & 2 & 2 \\ \hline
\end{tabularx}
\end{center}
L'écran affiche \texttt{2 2} (car $12 = 5 \times 2 + 2$).
\item \textbf{Analyse :}
\begin{itemize}
\item \texttt{5 <- x} : \textbf{INCORRECT}. Côté gauche = littéral (constante). \emph{Correction :} \texttt{x <- 5}.
\item \texttt{total <- "100"} : \textbf{INCORRECT} si \texttt{total} est \texttt{Entier} (incompatibilité type). \emph{Correction :} \texttt{total <- 100} ou déclarer \texttt{total : Chaine}.
\item \texttt{Lire("Entrez age : ", age)} : \textbf{INCORRECT}. \texttt{Lire} ne prend qu'\textbf{un seul argument} (la variable). \emph{Correction :} \texttt{Ecrire("Entrez age : ")} ; \texttt{Lire(age)}.
\item \texttt{Ecrire(age >= 18)} : \textbf{CORRECT}. Affiche \texttt{Vrai} ou \texttt{Faux} selon la valeur de \texttt{age}.
\end{itemize}
\end{enumerate}
\end{corrige}

\coursec{Rappels et notion de langage de définition des algorithmes}

Avant d'écrire notre première ligne d'algorithme, posons-nous la question fondamentale : \emph{qu'est-ce qu'un algorithme ?} Imagine que tu donnes la recette du \emph{ndolé} à un ami : « Lave les feuilles, fais-les bouillir, ajoute l'arachide, le poisson, laisse mijoter… ». Cette suite d'étapes ordonnées, précises et finies pour obtenir un résultat garanti, \textbf{c'est un algorithme}. En informatique, l'ordinateur est ton « cuisinier » : il ne comprend que des ordres extrêmement précis, sans ambiguïté, exprimés dans un langage qu'il sait exécuter.

\begin{definition}[Algorithme]
Un \cle{algorithme} est une suite \textbf{finie}, \textbf{précise} et \textbf{ordonnée} d'instructions élémentaires permettant de résoudre un problème ou d'obtenir un résultat à partir de données d'entrée. Il doit être \textbf{général} (fonctionner pour toute donnée valide) et \textbf{efficace} (terminer en un temps raisonnable).
\end{definition}

\begin{propriete}[Caractéristiques d'un algorithme]
Tout algorithme valide respecte quatre propriétés essentielles :
\begin{enumerate}
\item \textbf{Finitude :} Il se termine après un nombre fini d'étapes (pas de boucle infinie).
\item \textbf{Précision (non-ambiguïté) :} Chaque instruction a un sens unique, sans place à l'interprétation.
\item \textbf{Généralité :} Il traite une \emph{classe} de problèmes (ex. : « calculer la moyenne de \emph{deux} notes » fonctionne pour 10 et 12, comme pour 14,5 et 8,25).
\item \textbf{Entrées/Sorties définies :} Les données initiales (entrées) et les résultats attendus (sorties) sont explicitement identifiés.
\end{enumerate}
\end{propriete}

\courssub{Pourquoi un langage de définition (pseudo-code) ?}

Tu connais le français (langage naturel) et tu entendras parler de Python, C, JavaScript (langages de programmation). Le langage naturel est trop flou : « ajoute un peu de sel » — combien ? Le langage machine (binaire) ou même le C sont trop rigides pour \emph{réfléchir} à la logique. Le \cle{pseudo-code} (ou langage algorithmique) est le langage de \emph{conception} : il utilise des mots-clés en français (ou anglais) structurés comme du code, mais sans la syntaxe stricte d'un compilateur. C'est le « brouillon propre » du programmeur.

\begin{exemplebox}[Même problème, trois langages]
\textbf{Problème :} Demander un nombre et afficher son double.
\begin{center}
\begin{tabularx}{\linewidth}{|l|X|}
\hline
\rowcolor{popblueL} \textbf{Français (naturel)} & « Écris "Donne un nombre", lis le nombre, calcule le double, écris le résultat. » \\
\hline
\rowcolor{popgreenL} \textbf{Pseudo-code (cours 3ème)} & \texttt{Algorithme Double}\newline\texttt{Var x, d : entier}\newline\texttt{Debut}\newline\texttt{~~Ecrire("Donne un nombre : ")}\newline\texttt{~~Lire(x)}\newline\texttt{~~d <- x * 2}\newline\texttt{~~Ecrire("Le double est ", d)}\newline\texttt{Fin} \\
\hline
\rowcolor{poporangeL} \textbf{Python (exécution)} & \texttt{x = int(input("Donne un nombre : "))}\newline\texttt{d = x * 2}\newline\texttt{print("Le double est", d)} \\
\hline
\end{tabularx}
\end{center}
\legende{Le pseudo-code est le pont entre la logique humaine et le code machine. Il est universel : on le comprend sans connaître Python ni C.}
\end{exemplebox}

\begin{aretenir}[Le trio gagnant]
\begin{itemize}
\item \textbf{Langage naturel} : pour \emph{formuler} le besoin (cahier des charges).
\item \textbf{Pseudo-code} : pour \emph{concevoir} la logique (algorithme), indépendamment de la machine.
\item \textbf{Langage de programmation (C, Python...)} : pour \emph{implémenter} et \emph{exécuter} sur l'ordinateur.
\end{itemize}
En 3ème, nous maîtrisons le \textbf{pseudo-code} : c'est la compétence examinée au BEPC.
\end{aretenir}

\coursec{Structure de base d'un algorithme}

Tout algorithme écrit en pseudo-code dans ce cours respecte une \textbf{architecture rigide en trois blocs}. Respecter cette structure, c'est garantir la lisibilité et la correction.

\begin{propriete}[Squelette standard d'un algorithme]
\begin{lstlisting}
Algorithme <NomSignificatif>        // 1. En-tête : nom unique, sans espace, explicite
Const                                 // 2. Bloc constantes (optionnel)
    <NOM_CONST> = <valeur> : <type>   //    Noms en MAJUSCULES par convention
Var                                   // 3. Bloc variables (obligatoire si données)
    <nom_var> : <type>                //    Noms en minuscules (camelCase autorisé)
    <nom_var2>, <nom_var3> : <type>   //    Plusieurs variables de même type sur une ligne
Debut                                 // 4. Début du corps principal
    <Instruction_1>                   //    Instructions indentées (2 espaces)
    <Instruction_2>
    ...
Fin                                   // 5. Fin obligatoire
\end{lstlisting}
\end{propriete}

\begin{methode}[Règles de rédaction du pseudo-code]
Pour que ton algorithme soit lu et corrigé facilement (par ton prof, par toi-même plus tard), applique ces règles :
\begin{enumerate}
\item \textbf{Nommage :} \texttt{Algorithme CalculMoyenne} (PascalCase), \texttt{Const TVA = 19,25} (MAJUSCULES), \texttt{Var note, moyenne : reel} (minuscules).
\item \textbf{Indentation :} Décale de 2 espaces (ou une tabulation) tout ce qui est entre \texttt{Debut} et \texttt{Fin}. C'est \emph{visuel} : on voit la structure.
\item \textbf{Commentaires :} Utilise \texttt{//} pour expliquer \emph{pourquoi}, pas \emph{quoi} (le code montre le quoi). Ex : \texttt{m <- (n1+n2)/2 // moyenne arithmétique}.
\item \textbf{Une instruction par ligne :} Ne mets pas \texttt{Lire(x) Ecrire(x)} sur la même ligne.
\item \textbf{Fin de ligne :} En pseudo-code, \textbf{pas de point-virgule ;} à la fin des lignes (contrairement au C).
\end{enumerate}
\end{methode}

\begin{attention}[Pièges de structure fréquents]
\begin{itemize}
\item Oublier \texttt{Debut} ou \texttt{Fin} : l'algorithme est invalide.
\item Mélanger déclarations et instructions : \texttt{Var} et \texttt{Const} \textbf{uniquement} avant \texttt{Debut}.
\item Écrire du code en dehors de \texttt{Debut ... Fin} : tout traitement doit être à l'intérieur.
\item Mettre des accents dans les noms de variables : \texttt{var âge} $\rightarrow$ \textbf{interdit}. Utilise \texttt{age} ou \texttt{ageEleve}.
\end{itemize}
\end{attention}

\coursec{Variables, constantes et types de données}

Un algorithme manipule des \emph{valeurs}. Pour les stocker en mémoire le temps du calcul, on utilise des \textbf{variables} (contenu changeant) ou des \textbf{constantes} (contenu figé). Chaque « boîte » a un \textbf{type} qui détermine quelles valeurs elle peut contenir et quelles opérations sont autorisées.

\begin{definition}[Variable]
Une \cle{variable} est une zone de mémoire identifiée par un \textbf{nom} (identificateur), qui contient \textbf{une seule valeur à la fois} d'un \textbf{type} donné, et dont la valeur \textbf{peut être modifiée} pendant l'exécution par une affectation ou une lecture.
\end{definition}

\begin{definition}[Constante]
Une \cle{constante} est une zone de mémoire identifiée par un nom, qui contient une valeur d'un type donné, \textbf{fixée à la déclaration} et \textbf{interdite de modification} pendant toute l'exécution. Elle sert à nommer une valeur fixe (taux, seuil, $\pi$) pour la lisibilité et la maintenance.
\end{definition}

\begin{propriete}[Types de données de base (niveau 3ème)]
\begin{center}
\begin{tabularx}{\linewidth}{|l|X|X|X|}
\hline
\rowcolor{popblueL} \textbf{Type} & \textbf{Mot-clé pseudo-code} & \textbf{Valeurs admises} & \textbf{Exemples} \\
\hline
Entier & \texttt{entier} & Nombres sans virgule (positifs, négatifs, zéro) & \texttt{-5, 0, 42, 10000} \\
\hline
Réel & \texttt{reel} & Nombres à virgule (notation décimale) & \texttt{3,14 ; -12,5 ; 0,0 ; 19,25} \\
\hline
Chaîne de caractères & \texttt{chaine} & Suite de caractères entre guillemets & \texttt{"Yaoundé" ; "Note : " ; "123"} \\
\hline
Booléen & \texttt{booleen} & Deux valeurs logiques seulement & \texttt{Vrai, Faux} \\
\hline
\end{tabularx}
\end{center}
\end{propriete}

\begin{exemplebox}[Déclarations variées]
\begin{lstlisting}
Algorithme ExempleTypes
Const PI = 3,14159       // reel (implicite par la valeur)
Const NOM_ETAB = "Lycée Bilingue"  // chaine
Var age, nbEleves : entier
Var moyenne, taille : reel
Var nomEleve, mention : chaine
Var estAdmis : booleen
Debut
    // ... corps de l'algorithme
Fin
\end{lstlisting}
\end{exemplebox}

\begin{attention}[Règles d'identificateurs (noms de variables/constants)]
\begin{enumerate}
\item Commence par une \textbf{lettre} (a-z, A-Z). Pas de chiffre au début.
\item Autorisé : lettres, chiffres, \texttt{\_} (underscore). \textbf{Interdit :} espaces, accents (\texttt{é, à, ç}), caractères spéciaux (\texttt{@, \#, \$}), mots-clés réservés (\texttt{Var, Const, Si, Pour, Lire...}).
\item \textbf{La casse compte :} \texttt{moyenne} et \texttt{Moyenne} sont deux variables différentes. \textbf{Convention :} variables en \texttt{minuscules} (ou \texttt{camelCase}), constantes en \texttt{MAJUSCULES}.
\item Choisis des noms \textbf{significatifs} : \texttt{mt} $\rightarrow$ mauvais ; \texttt{montantTTC} $\rightarrow$ excellent.
\end{enumerate}
\end{attention}

\begin{aretenir}[Variable vs Constante : le bon réflexe]
\begin{itemize}
\item \textbf{Constante :} « Le prix du repas est 500 FCFA », « La TVA est 19,25\% », « $\pi \approx 3,14$ ». $\rightarrow$ \texttt{Const PRIX = 500}.
\item \textbf{Variable :} « Saisir la note de l'élève », « Calculer la moyenne », « Compter les élèves présents ». $\rightarrow$ \texttt{Var note, moyenne, compteur : ...}.
\end{itemize}
Utiliser une constante pour une valeur fixe évite les « nombres magiques » dans le code et permet de changer la valeur en un seul endroit si la loi change (ex. TVA passant à 20\%).
\end{aretenir}

\coursec{Les instructions de base : lecture, écriture, affectation}

Maintenant que nous savons nommer et typer nos données, voyons les trois instructions fondamentales qui font \emph{vivre} l'algorithme : entrer des données, les transformer, sortir des résultats.

\courssub{L'instruction de lecture : \texttt{Lire}}

\begin{propriete}[Syntaxe \texttt{Lire}]
\texttt{Lire(<variable>)} \quad ou \quad \texttt{Lire(<var1>, <var2>, ...)}
\begin{itemize}
\item Effectue une \textbf{pause} : l'algorithme attend que l'utilisateur tape une valeur au clavier et valide (Entrée).
\item La valeur saisie est \textbf{convertie} selon le type de la variable (entier $\rightarrow$ entier, réel $\rightarrow$ réel, etc.) et stockée dans la variable.
\item La variable \textbf{doit être déclarée} avant dans \texttt{Var}.
\item \textbf{Message utilisateur :} \texttt{Lire} n'affiche rien. On place \textbf{toujours} un \texttt{Ecrire} avant pour guider l'utilisateur.
\end{itemize}
\end{propriete}

\begin{exemplebox}[Utilisation de \texttt{Lire}]
\begin{lstlisting}
Ecrire("Entrez votre age : ")
Lire(age)              // age est de type entier
Ecrire("Entrez votre taille (m) : ")
Lire(taille)           // taille est de type reel
Ecrire("Entrez votre nom : ")
Lire(nom)              // nom est de type chaine
\end{lstlisting}
\end{exemplebox}

\courssub{L'instruction d'écriture : \texttt{Écrire}}

\begin{propriete}[Syntaxe \texttt{Écrire}]
\texttt{Écrire(<expression\_1>, <expression\_2>, ...)}
\begin{itemize}
\item Affiche à l'écran la suite des expressions, séparées par des virgules dans l'instruction.
\item Les \textbf{chaînes de caractères} (messages) sont entre \textbf{guillemets doubles} \texttt{"..."}.
\item Les \textbf{variables} s'écrivent \textbf{sans guillemets} : leur \emph{valeur} est affichée.
\item On peut mélanger texte et variables : \texttt{Écrire("La moyenne est ", m, " sur 20")}.
\item Pas de retour à la ligne automatique standard (comportement dépend de l'interpréteur) ; on gère l'affichage par des \texttt{Écrire} successives.
\end{itemize}
\end{propriete}

\begin{attention}[Guillemets : le piège classique]
\begin{lstlisting}
// ERREUR : affiche littéralement "moyenne"
Ecrire("moyenne")       

// CORRECT : affiche la VALEUR de la variable moyenne
Ecrire(moyenne)         

// CORRECT : message + valeur
Ecrire("Moyenne = ", moyenne)
\end{lstlisting}
\textbf{Règle d'or :} \texttt{"texte"} = texte fixe ; \texttt{variable} = valeur variable.
\end{attention}

\courssub{L'instruction d'affectation : $\leftarrow$}

C'est le cœur du \emph{traitement}. Elle calcule une expression et stocke le résultat dans une variable.

\begin{propriete}[Syntaxe et sémantique de l'affectation]
\texttt{<variable> $\leftarrow$ <expression>}
\begin{enumerate}
\item \textbf{Évaluer} l'expression à droite (calculer sa valeur actuelle en utilisant les valeurs \emph{actuelles} des variables).
\item \textbf{Stocker} ce résultat dans la variable à gauche (écraser l'ancienne valeur).
\item Le type de l'expression doit être \textbf{compatible} avec le type de la variable (entier dans entier, réel dans réel, etc.).
\end{enumerate}
\textbf{Attention :} $\leftarrow$ n'est \textbf{PAS} une égalité mathématique. \texttt{x $\leftarrow$ x + 1} signifie « prend la valeur actuelle de x, ajoute 1, remets dans x » (incrémentation). En maths, $x = x+1$ est faux.
\end{propriete}

\begin{propriete}[Opérateurs arithmétiques et priorités]
\begin{center}
\begin{tabularx}{\linewidth}{|l|l|X|}
\hline
\rowcolor{popblueL} \textbf{Opérateur} & \textbf{Sens} & \textbf{Priorité} \\
\hline
\texttt{+ , -} & Addition, Soustraction & 1 (la plus faible) \\
\hline
\texttt{* , /} & Multiplication, Division & 2 \\
\hline
\texttt{\^{}} ou \texttt{**} & Puissance & 3 (la plus forte) \\
\hline
\texttt{( )} & Parenthèses (forcent l'ordre) & Priorité absolue \\
\hline
\end{tabularx}
\end{center}
\textbf{Associativité :} Pour une même priorité, on évalue de \textbf{gauche à droite}. Ex : \texttt{10 - 3 - 2} = \texttt{(10-3)-2} = 5.
\end{propriete}

\begin{exemplebox}[Affectations et priorités]
\begin{lstlisting}
Var a, b, c, res : reel
Debut
    a <- 10
    b <- 3
    c <- 2
    res <- a + b * c      // 1. b*c = 6, 2. a+6 = 16  (res = 16)
    res <- (a + b) * c    // 1. (a+b) = 13, 2. 13*2 = 26 (res = 26)
    res <- a / b * c      // 1. a/b = 3.333..., 2. *2 = 6.666... (gauche à droite)
    a <- a + 1            // Incrémentation : a devient 11
Fin
\end{lstlisting}
\end{exemplebox}

\begin{experience}[TP Papier : Trace d'exécution manuelle]
\textbf{Objectif :} Comprendre l'état mémoire pas à pas.
\textbf{Matériel :} Feuille, stylo.
\textbf{Algorithme :}
\begin{lstlisting}
Algorithme TraceDemo
Var x, y, z : entier
Debut
    x <- 5
    y <- 2
    z <- x + y
    x <- z * 2
    y <- x - y
    Ecrire(x, y, z)
Fin
\end{lstlisting}
\textbf{Manipulation :} Complète le tableau de trace ci-dessous (colonnes = variables, lignes = instructions).
\begin{center}
\begin{tabularx}{\linewidth}{|l|X|X|X|}
\hline
\rowcolor{popblueL} \textbf{Instruction} & \textbf{x} & \textbf{y} & \textbf{z} \\
\hline
x $\leftarrow$ 5 & & & \\
\hline
y $\leftarrow$ 2 & & & \\
\hline
z $\leftarrow$ x + y & & & \\
\hline
x $\leftarrow$ z * 2 & & & \\
\hline
y $\leftarrow$ x - y & & & \\
\hline
Écrire(x, y, z) & \multicolumn{3}{X|}{Affiche : ...} \\
\hline
\end{tabularx}
\end{center}
\textbf{Interprétation :} Vérifie que la dernière ligne affiche \texttt{14 12 7}. Si ce n'est pas le cas, reprends chaque ligne : l'affectation change \textbf{une seule} variable à la fois.
\end{experience}

\coursec{Rédaction complète d'algorithmes sur des situations concrètes}

Dans les sections précédentes, tu as découvert les trois instructions de base du langage algorithmique : la \cle{lecture} (\texttt{Lire}), l'\cle{écriture} (\texttt{Écrire}) et l'\cle{affectation} ($\leftarrow$). Chacune prise isolément est simple. Mais un algorithme utile n'est jamais une seule instruction : c'est une \emph{suite ordonnée} d'instructions qui résout un problème concret, du début à la fin. Le moment est venu d'assembler ces briques pour rédiger des algorithmes complets, comme le ferait un programmeur face à un vrai besoin : calculer une moyenne de bulletin, établir une facture à la boutique, mesurer un champ. C'est toute la puissance de l'algorithmique : passer d'une situation de la vie courante à une méthode de calcul automatisable.

\courssub{La démarche de résolution en cinq étapes}

Face à un problème, le débutant a souvent envie d'écrire immédiatement des instructions. C'est une erreur : on construit un algorithme comme on construit une maison, en suivant un plan. Avant toute rédaction, il faut \emph{comprendre} le problème et le \emph{découper}.

\begin{methode}[Résoudre un problème par un algorithme : les 5 étapes]
Pour passer d'une situation concrète à un algorithme complet, on suit toujours la même démarche :
\begin{enumerate}
\item \textbf{Analyser le problème} : lire attentivement l'énoncé, identifier les \cle{données} (ce que l'utilisateur fournit, les \emph{entrées}) et les \cle{résultats} attendus (ce que l'algorithme doit produire, les \emph{sorties}). Expliciter la ou les formules de calcul.
\item \textbf{Déclarer} : choisir un nom significatif et un type adapté pour chaque variable et chaque constante nécessaire (données d'entrée, résultats de sortie, valeurs fixes du problème).
\item \textbf{Saisir} : écrire les instructions \texttt{Lire} qui permettent d'entrer les données, précédées chacune d'un \texttt{Écrire} explicite (message de prompt).
\item \textbf{Traiter} : écrire les calculs (affectations) qui transforment les données en résultats, en utilisant les formules dégagées à l'étape 1. Respecter les priorités opératoires (parentheses !).
\item \textbf{Afficher} : écrire les instructions \texttt{Écrire} qui présentent les résultats à l'utilisateur, avec un message clair et l'unité si pertinent.
\end{enumerate}
Ensuite, on \textbf{vérifie} l'algorithme par une \cle{trace d'exécution} : on choisit des valeurs concrètes, on « joue » l'algorithme à la main et on compare le résultat obtenu au calcul manuel.
\end{methode}

\begin{definition}[Trace d'exécution]
La \cle{trace d'exécution} (ou test d'exécution) d'un algorithme est le tableau qui montre, ligne par ligne, la valeur prise par chaque variable au cours du déroulement, pour des données d'entrée choisies. Elle permet de vérifier que l'algorithme produit bien le résultat attendu, sans ordinateur.
\end{definition}

\courssub{Situation 1 : la moyenne trimestrielle d'un élève}

\textbf{Problème.} Au collège, le professeur principal veut calculer rapidement la moyenne trimestrielle d'un élève en informatique, obtenue à partir de deux notes : la note du devoir et la note de la composition. La formule est :
\[ m = \frac{n_1 + n_2}{2} \]

\textbf{Analyse (étape 1).}
\begin{itemize}
\item \emph{Entrées} : deux notes $n_1$ et $n_2$ (des nombres pouvant avoir des virgules, comme 12,5). Type \texttt{reel}.
\item \emph{Sortie} : la moyenne $m$. Type \texttt{reel}.
\item \emph{Traitement} : additionner les deux notes, puis diviser par 2.
\end{itemize}

\textbf{Déclaration (étape 2).} Trois variables de type \cle{reel} : \texttt{n1}, \texttt{n2} et \texttt{m}.

\begin{exemplebox}[Algorithme MoyenneTrimestre]
\begin{lstlisting}
Algorithme MoyenneTrimestre
Var n1, n2, m : reel        // les deux notes et la moyenne
Debut
    Ecrire("Entrez la premiere note : ")
    Lire(n1)                 // saisie de la note 1
    Ecrire("Entrez la deuxieme note : ")
    Lire(n2)                 // saisie de la note 2
    m <- (n1 + n2) / 2       // traitement : formule de la moyenne
    Ecrire("La moyenne est : ", m)
Fin
\end{lstlisting}
\end{exemplebox}

\textbf{Vérification par une trace.} Testons avec $n_1 = 12$ et $n_2 = 16$ :

\begin{center}
\begin{tabularx}{\linewidth}{|l|X|X|X|}
\hline
\rowcolor{popblueL} \textbf{Instruction} & \textbf{n1} & \textbf{n2} & \textbf{m} \\
\hline
Lire(n1) & 12 & ? & ? \\
\hline
Lire(n2) & 12 & 16 & ? \\
\hline
m $\leftarrow$ (n1 + n2)/2 & 12 & 16 & 14 \\
\hline
Écrire(...) & \multicolumn{3}{X|}{affiche : « La moyenne est : 14 »} \\
\hline
\end{tabularx}
\end{center}

Le calcul manuel donne $(12 + 16)/2 = 28/2 = 14$ : l'algorithme est correct.

\begin{attention}[Le piège des parenthèses]
L'erreur la plus fréquente ici est d'écrire \texttt{m <- n1 + n2 / 2} \emph{sans parenthèses}. Or la division est prioritaire sur l'addition : l'ordinateur calcule d'abord \texttt{n2 / 2}, puis ajoute \texttt{n1}. Avec $n_1 = 12$ et $n_2 = 16$, on obtiendrait $12 + 8 = 20$ au lieu de $14$ ! La formule correcte exige les parenthèses : \texttt{(n1 + n2) / 2}. Retiens : \textbf{en algorithmique comme en mathématiques, les priorités opératoires s'appliquent}.
\end{attention}

La figure suivante montre l'\cle{organigramme} de cet algorithme : c'est une représentation graphique du déroulement, où l'on voit clairement la logique \emph{entrée $\rightarrow$ traitement $\rightarrow$ sortie}, en parallèle du pseudo-code.

\begin{popfigure}
\begin{center}
\begin{tikzpicture}[
  font=\small,
  startstop/.style={rounded rectangle, rounded rectangle arc length=120, draw=popdark, fill=popgoldL, minimum width=2.6cm, minimum height=0.8cm, align=center},
  io/.style={trapezium, trapezium left angle=70, trapezium right angle=110, draw=popdark, fill=popblueL, minimum width=3.2cm, minimum height=0.8cm, align=center},
  process/.style={rectangle, draw=popdark, fill=popgreenL, minimum width=3.6cm, minimum height=0.8cm, align=center},
  arrow/.style={-{Stealth[length=2.5mm]}, thick, popdark}
]
\node[startstop] (debut) {D\'ebut};
\node[io, below=0.55cm of debut] (lire) {Lire n1, n2};
\node[process, below=0.55cm of lire] (calc) {m $\leftarrow$ (n1 + n2) / 2};
\node[io, below=0.55cm of calc] (ecrire) {Écrire m};
\node[startstop, below=0.55cm of ecrire] (fin) {Fin};
\draw[arrow] (debut) -- (lire);
\draw[arrow] (lire) -- (calc);
\draw[arrow] (calc) -- (ecrire);
\draw[arrow] (ecrire) -- (fin);
\node[right=0.35cm of lire, text=popmuted, align=left] {\footnotesize entrées};
\node[right=0.35cm of calc, text=popmuted, align=left] {\footnotesize traitement};
\node[right=0.35cm of ecrire, text=popmuted, align=left] {\footnotesize sortie};
\end{tikzpicture}
\end{center}
\legende{Organigramme de l'algorithme MoyenneTrimestre : chaque étape du pseudo-code correspond à un symbole (trapèze pour les entrées/sorties, rectangle pour le traitement).}
\end{popfigure}

\courssub{Situation 2 : le montant TTC d'un achat à la boutique}

\textbf{Problème.} La boutique du quartier affiche ses prix hors taxe (HT). La gérante veut un outil qui calcule le montant \emph{toutes taxes comprises} (TTC) en appliquant la TVA camerounaise de 19,25\,\%. La formule est :
\[ \text{TTC} = \text{HT} + \text{HT} \times \frac{19,25}{100} \]

\textbf{Analyse.}
\begin{itemize}
\item \emph{Entrée} : le montant hors taxe \texttt{ht} (en FCFA). Type \texttt{reel}.
\item \emph{Sortie} : le montant TTC \texttt{ttc}. Type \texttt{reel}.
\item \emph{Donnée fixe} : le taux de TVA, 19,25, qui ne change jamais pendant l'exécution : c'est une \cle{constante}.
\end{itemize}

C'est l'occasion de justifier un choix important : la TVA est déclarée comme \emph{constante} car sa valeur est fixée par la loi et ne doit pas être modifiée par l'utilisateur ; les montants \texttt{ht} et \texttt{ttc} sont des \emph{variables réelles} car ils changent à chaque client et peuvent comporter des décimales.

\begin{exemplebox}[Algorithme AchatTTC]
\begin{lstlisting}
Algorithme AchatTTC
Const TVA = 19,25            // taux de TVA en pourcentage (fixe)
Var ht, ttc : reel           // montant hors taxe et montant TTC
Debut
    Ecrire("Entrez le montant hors taxe (FCFA) : ")
    Lire(ht)                          // saisie du prix HT
    ttc <- ht + ht * TVA / 100        // calcul du prix TTC
    Ecrire("Le montant TTC est : ", ttc, " FCFA")
Fin
\end{lstlisting}
\end{exemplebox}

\textbf{Vérification.} Avec $\texttt{ht} = 10\,000$ FCFA :
\[ \text{TTC} = 10\,000 + 10\,000 \times \frac{19,25}{100} = 10\,000 + 1\,925 = 11\,925 \ \text{FCFA} \]
L'algorithme affiche bien 11925 FCFA : le test confirme le calcul manuel.

\begin{aretenir}[Variable ou constante ?]
On déclare en \cle{constante} toute valeur fixe du problème (un taux, un prix unitaire, le nombre $\pi$) : elle est définie une fois pour toutes et ne peut pas être modifiée pendant l'exécution. On déclare en \cle{variable} tout ce qui est saisi par l'utilisateur ou calculé par le programme. Ce choix rend l'algorithme plus sûr et plus lisible.
\end{aretenir}

\courssub{Situation 3 : périmètre et aire d'un champ rectangulaire}

\textbf{Problème.} Un agriculteur de Ngaoundéré possède un champ rectangulaire. Il veut connaître le \cle{périmètre} (pour acheter la longueur de fil de clôture) et l'\cle{aire} (pour estimer la quantité de semences). Pour un rectangle de longueur $L$ et de largeur $l$ :
\[ P = 2 \times (L + l) \qquad \text{et} \qquad A = L \times l \]

\textbf{Analyse.}
\begin{itemize}
\item \emph{Entrées} : la longueur \texttt{L} et la largeur \texttt{l}, en mètres (réels).
\item \emph{Sorties} : le périmètre \texttt{P} (en mètres) et l'aire \texttt{A} (en mètres carrés).
\item \emph{Traitements} : deux calculs indépendants, effectués l'un après l'autre.
\end{itemize}

\begin{exemplebox}[Algorithme ChampRectangulaire]
\begin{lstlisting}
Algorithme ChampRectangulaire
Var L, l, P, A : reel       // dimensions, perimetre et aire
Debut
    Ecrire("Longueur du champ (m) : ")
    Lire(L)
    Ecrire("Largeur du champ (m) : ")
    Lire(l)
    P <- 2 * (L + l)         // perimetre du rectangle
    A <- L * l               // aire du rectangle
    Ecrire("Le perimetre est : ", P, " m")
    Ecrire("L'aire est : ", A, " m2")
Fin
\end{lstlisting}
\end{exemplebox}

\textbf{Vérification.} Avec $L = 50$ m et $l = 30$ m : $P = 2 \times (50 + 30) = 160$ m de clôture, et $A = 50 \times 30 = 1\,500$ m$^2$. La trace d'exécution donne les mêmes valeurs.

\courssub{Exercice résolu et synthèse}

\begin{exoresolu}[La cantine du lycée]
La cantine d'un lycée de Yaoundé vend le repas à 500 FCFA. Le surveillant veut un algorithme qui demande le nombre d'élèves inscrits et affiche la recette totale attendue.
\tcblower
\textbf{Étape 1 — Analyse.} Entrée : le nombre d'élèves \texttt{n} (un entier). Sortie : la recette \texttt{r} (en FCFA). Donnée fixe : le prix du repas, 500 FCFA, déclaré en constante. Traitement : $r = n \times 500$.

\textbf{Étapes 2 à 5 — Algorithme.}
\begin{lstlisting}
Algorithme RecetteCantine
Const PRIX = 500             // prix d'un repas en FCFA (fixe)
Var n, r : entier            // nombre d'eleves et recette
Debut
    Ecrire("Nombre d'eleves inscrits : ")
    Lire(n)
    r <- n * PRIX            // calcul de la recette
    Ecrire("La recette attendue est : ", r, " FCFA")
Fin
\end{lstlisting}
\textbf{Vérification.} Avec $n = 240$ élèves : $r = 240 \times 500 = 120\,000$ FCFA. Le calcul manuel confirme le résultat. La démarche est justifiée : \texttt{n} est un entier car on ne compte pas une fraction d'élève ; \texttt{PRIX} est une constante car le tarif ne change pas pendant l'exécution.
\end{exoresolu}

\begin{savaistu}[Le pseudo-code est universel]
Les algorithmes que tu viens d'écrire en pseudo-code français peuvent être traduits presque mot à mot dans n'importe quel langage de programmation : Python, C, Java, JavaScript\ldots{} Seule la syntaxe change ; la \emph{logique} (saisir, calculer, afficher) reste identique. C'est pourquoi les programmeurs du monde entier commencent toujours par réfléchir en algorithmes avant de coder : le pseudo-code est leur langue commune.
\end{savaistu}

\begin{exercice}[À toi de jouer]
\begin{enumerate}
\item Un taxi-moto (benskin) facture la course 100 FCFA de prise en charge plus 50 FCFA par kilomètre parcouru. Rédige l'algorithme complet (analyse, déclarations, instructions) qui calcule le prix d'une course, puis vérifie-le avec une distance de 8 km.
\item Rédige l'algorithme qui calcule la moyenne générale d'un élève à partir de \emph{trois} notes, et teste-le avec 10, 14 et 16.
\end{enumerate}
\end{exercice}

\begin{corrige}[Exercices 1 et 2]
\textbf{Exercice 1.}
\textbf{Analyse.} Entrée : distance \texttt{d} (reel) ; sortie : prix \texttt{p} (reel). Constantes : prise en charge 100, prix au km 50. Traitement : $p = 100 + 50 \times d$.
\begin{lstlisting}
Algorithme PrixCourse
Const PRISE = 100
Const PRIXKM = 50
Var d, p : reel
Debut
    Ecrire("Distance parcourue (km) : ")
    Lire(d)
    p <- PRISE + PRIXKM * d
    Ecrire("Le prix de la course est : ", p, " FCFA")
Fin
\end{lstlisting}
\textbf{Test} avec $d = 8$ : $p = 100 + 50 \times 8 = 500$ FCFA. Correct.

\textbf{Exercice 2.}
\textbf{Analyse.} Entrées : trois notes \texttt{n1, n2, n3} (reel). Sortie : moyenne \texttt{m} (reel). Traitement : $m = (n1 + n2 + n3) / 3$.
\begin{lstlisting}
Algorithme MoyenneGenerale
Var n1, n2, n3, m : reel
Debut
    Ecrire("Note 1 : "); Lire(n1)
    Ecrire("Note 2 : "); Lire(n2)
    Ecrire("Note 3 : "); Lire(n3)
    m <- (n1 + n2 + n3) / 3
    Ecrire("La moyenne generale est : ", m)
Fin
\end{lstlisting}
\textbf{Test} : $(10 + 14 + 16)/3 = 40/3 \approx 13,33$. Correct.
\end{corrige}

\begin{aretenir}[L'essentiel de la section]
Tout algorithme de cette leçon suit le même squelette : \textbf{déclarer} les variables et constantes, \textbf{saisir} les données avec \texttt{Lire}, \textbf{traiter} par affectation, \textbf{afficher} les résultats avec \texttt{Écrire}. On vérifie toujours par une trace d'exécution comparée au calcul manuel. Ces algorithmes sont dits \cle{séquentiels} : les instructions s'exécutent dans l'ordre, de haut en bas, sans choix ni répétition. Dans l'unité suivante, tu les enrichiras avec les \emph{structures conditionnelles} (\texttt{Si...Alors}) et les \emph{boucles} (\texttt{Pour}, \texttt{Tant que}), qui permettront de traiter des situations bien plus riches.
\end{aretenir}

\newpage

\coursec{Activité d'intégration}

\coursec{Leçon 13 : Présentation du langage de définition des algorithmes}

\courssub{Objectifs de l'activité}

À l'issue de cette activité d'intégration, tu seras capable de :

\begin{itemize}
\item analyser une situation concrète de la vie courante pour y repérer les \cle{entrées}, les \cle{traitements} et les \cle{sorties} ;
\item choisir correctement les \cle{noms} et les \cle{types} des variables et des constantes d'un algorithme ;
\item rédiger un algorithme complet en \cle{pseudo-code}, en respectant la structure de base (en-tête, déclarations, corps) ;
\item vérifier le bon fonctionnement d'un algorithme à l'aide d'un \cle{tableau de trace} ;
\item présenter et \cle{justifier} ta démarche devant la classe, en rédigeant une conclusion claire.
\end{itemize}

\courssub{Situation}

Au marché Mokolo, à Yaoundé, le jeune Armand tient un petit kiosque de vente de crédit téléphonique (MTN, Orange, Nexttel). Comme tous les vendeurs de son quartier, Armand n'achète pas le crédit au prix où il le revend : il l'achète à un \cle{prix de gros} auprès d'un distributeur, puis il le revend à ses clients en ajoutant une \cle{marge fixe de 5 \%} sur le prix d'achat.

Chaque jour, Armand effectue des dizaines de ventes de montants différents : 500 FCFA, 1 000 FCFA, 2 500 FCFA, etc. À chaque fois, il doit calculer mentalement deux nombres :

\begin{itemize}
\item le \cle{prix de vente} à demander au client : prix d'achat augmenté de 5 \% ;
\item son \cle{bénéfice} sur la vente : la différence entre le prix de vente et le prix d'achat.
\end{itemize}

Un soir, en rentrant chez lui, Armand se trompe dans ses comptes et croit avoir gagné plus que la réalité. Il se dit alors : « Si seulement j'avais une petite procédure automatique qui, à partir du prix d'achat saisi, me donne immédiatement le prix de vente et mon bénéfice ! »

Armand sait que tu apprends l'algorithmique en classe de 3ème. Il te demande donc de l'aider. Pour tester ta solution, il te propose une vente réelle de sa journée : un achat de crédit dont le \textbf{prix d'achat est de 2 000 FCFA}.

\textbf{Problème posé :} rédiger un algorithme complet, en pseudo-code, qui lit le prix d'achat d'un crédit, calcule le prix de vente (marge de 5 \%) et le bénéfice, puis affiche ces deux résultats. Vérifier ensuite cet algorithme par un tableau de trace avec un achat de 2 000 FCFA.

\courssub{Consignes}

Travail en groupes de 3 ou 4 élèves. Chaque groupe répond par écrit aux questions suivantes, puis un rapporteur présente la solution du groupe devant la classe.

\begin{enumerate}
\item \textbf{Analyse du problème.} Pour cette situation, identifiez et notez sur votre feuille :
\begin{enumerate}
\item les \cle{données d'entrée} (ce que l'utilisateur fournit à l'algorithme) ;
\item les \cle{résultats de sortie} (ce que l'algorithme doit afficher) ;
\item les \cle{traitements} (les calculs à effectuer pour passer des entrées aux sorties). Écrivez les formules mathématiques utilisées.
\end{enumerate}
\item \textbf{Choix des variables et des constantes.} Proposez un nom significatif pour chaque information manipulée, en précisant s'il s'agit d'une \cle{variable} ou d'une \cle{constante}, ainsi que son \cle{type} (entier, réel, chaîne de caractères...). Justifiez chaque choix : pourquoi la marge est-elle une constante et non une variable ?
\item \textbf{Rédaction de l'algorithme.} Rédigez l'algorithme complet en pseudo-code, en respectant les trois parties de la structure de base : l'\cle{en-tête} (nom de l'algorithme), la partie \cle{déclarations} (constantes et variables avec leurs types) et le \cle{corps} (instructions entre Début et Fin). Utilisez les instructions de lecture, d'affectation et d'écriture vues dans la leçon.
\item \textbf{Vérification par une trace d'exécution.} Construisez le tableau de trace de votre algorithme pour un prix d'achat de 2 000 FCFA : un tableau dont les colonnes sont les noms des variables et dont chaque ligne correspond à une instruction exécutée. Vérifiez que les valeurs affichées à la fin sont bien celles attendues (faites le calcul à la main pour comparer).
\item \textbf{Présentation et justification.} Le rapporteur présente la solution du groupe au tableau et justifie chaque choix (noms, types, ordre des instructions). La classe valide la solution ou propose des corrections. Que se passerait-il si l'on calculait le bénéfice \emph{avant} le prix de vente ?
\end{enumerate}

\courssub{Ressources à mobiliser}

\begin{aretenir}[Ce qu'il faut avoir en tête pour réussir]
\begin{itemize}
\item \textbf{Structure de base d'un algorithme :} un en-tête (\texttt{Algorithme nom\_de\_l\_algorithme}), une partie déclarations (\texttt{Constante}, \texttt{Variable} avec le type de chacune), et un corps délimité par \texttt{Début} et \texttt{Fin}.
\item \textbf{Instructions de base :} \texttt{Lire(variable)} pour saisir une entrée ; \texttt{variable $\leftarrow$ expression} pour l'affectation ; \texttt{Écrire(...)} pour afficher un résultat ou un message.
\item \textbf{Variable et constante :} une \cle{variable} peut changer de valeur pendant l'exécution ; une \cle{constante} garde toujours la même valeur, fixée dans les déclarations.
\item \textbf{Types de données :} \texttt{Entier}, \texttt{Réel}, \texttt{Caractère}, \texttt{Chaîne}, \texttt{Booléen}. Un prix peut comporter des centimes : on choisit le type \texttt{Réel}.
\item \textbf{Calcul d'un pourcentage :} augmenter une valeur de 5 \% revient à la multiplier par $1 + \dfrac{5}{100} = 1{,}05$.
\item \textbf{Tableau de trace :} tableau donnant, instruction après instruction, la valeur de chaque variable ; il permet de vérifier l'algorithme « à la main ».
\end{itemize}
\end{aretenir}

\courssub{Démarche et corrigé commenté}

\begin{exoresolu}[L'algorithme du vendeur de crédit téléphonique]
\textbf{Énoncé.} Armand revend le crédit téléphonique avec une marge fixe de 5 \% sur le prix d'achat. Rédiger un algorithme complet qui lit le prix d'achat d'un crédit, calcule puis affiche le prix de vente et le bénéfice. Vérifier l'algorithme par un tableau de trace pour un achat de 2 000 FCFA.

\tcblower

\textbf{Étape 1 : analyse du problème (entrées, traitements, sorties).}

\begin{itemize}
\item \textbf{Entrée :} le prix d'achat du crédit, fourni par le vendeur (en FCFA).
\item \textbf{Traitements :} deux calculs.
\begin{itemize}
\item Prix de vente = prix d'achat augmenté de 5 \%, c'est-à-dire :
\[ \text{prix de vente} = \text{prix d'achat} \times 1{,}05 \]
\item Bénéfice = différence entre le prix de vente et le prix d'achat :
\[ \text{bénéfice} = \text{prix de vente} - \text{prix d'achat} \]
\end{itemize}
\item \textbf{Sorties :} le prix de vente et le bénéfice, affichés à l'écran avec des messages clairs.
\end{itemize}

\textbf{Étape 2 : choix des variables et de la constante.}

\begin{center}
\begin{tabularx}{\linewidth}{|l|l|X|}
\hline
\rowcolor{popblueL} \textbf{Nom} & \textbf{Nature et type} & \textbf{Rôle (justification)} \\
\hline
\texttt{TAUX} & Constante Réel = 1,05 & Coefficient de la marge : il est fixé à 5 \% une fois pour toutes, il ne change jamais pendant l'exécution : c'est donc une \cle{constante}. \\
\hline
\texttt{prixAchat} & Variable Réel & Prix d'achat saisi par le vendeur : il change à chaque vente. Type Réel car un prix peut avoir des centimes. \\
\hline
\texttt{prixVente} & Variable Réel & Résultat du calcul du prix de vente. \\
\hline
\texttt{benefice} & Variable Réel & Résultat du calcul du bénéfice. \\
\hline
\end{tabularx}
\end{center}

Les noms choisis sont \cle{significatifs} : ils décrivent clairement le contenu, sans espace ni accent, et commencent par une lettre.

\textbf{Étape 3 : rédaction de l'algorithme complet.}

\begin{lstlisting}
Algorithme vente_credit
// Calcule le prix de vente et le benefice d'une vente de credit

Constante
    TAUX : Reel <- 1,05        // marge fixe de 5 %

Variable
    prixAchat : Reel           // prix d'achat saisi (FCFA)
    prixVente : Reel           // prix de vente calcule (FCFA)
    benefice  : Reel           // benefice realise (FCFA)

Debut
    // 1. Lecture de l'entree
    Ecrire("Entrez le prix d'achat du credit (FCFA) : ")
    Lire(prixAchat)

    // 2. Traitements
    prixVente <- prixAchat * TAUX
    benefice  <- prixVente - prixAchat

    // 3. Affichage des sorties
    Ecrire("Prix de vente : ", prixVente, " FCFA")
    Ecrire("Benefice : ", benefice, " FCFA")
Fin
\end{lstlisting}

\textbf{Justification de l'ordre des instructions :} on doit d'abord \texttt{Lire} le prix d'achat avant tout calcul (sinon \texttt{prixAchat} n'a aucune valeur) ; on doit calculer \texttt{prixVente} \emph{avant} \texttt{benefice}, car le bénéfice utilise la valeur de \texttt{prixVente}. Inverser ces deux affectations provoquerait une erreur : le bénéfice serait calculé avec un prix de vente encore inconnu.

\textbf{Étape 4 : vérification par tableau de trace (achat de 2 000 FCFA).}

\begin{center}
\begin{tabularx}{\linewidth}{|X|l|l|l|}
\hline
\rowcolor{popblueL} \textbf{Instruction exécutée} & \textbf{prixAchat} & \textbf{prixVente} & \textbf{benefice} \\
\hline
Début & indéfini & indéfini & indéfini \\
\hline
Lire(prixAchat) & 2 000 & indéfini & indéfini \\
\hline
prixVente $\leftarrow$ prixAchat $\times$ TAUX & 2 000 & 2 100 & indéfini \\
\hline
benefice $\leftarrow$ prixVente $-$ prixAchat & 2 000 & 2 100 & 100 \\
\hline
Ecrire(...) & 2 000 & 2 100 & 100 \\
\hline
\end{tabularx}
\end{center}

\textbf{Vérification à la main :} $2\,000 \times 1{,}05 = 2\,100$ FCFA et $2\,100 - 2\,000 = 100$ FCFA. L'algorithme affiche donc :

\begin{lstlisting}
Prix de vente : 2100 FCFA
Benefice : 100 FCFA
\end{lstlisting}

Les résultats de la trace correspondent exactement au calcul manuel : \textbf{l'algorithme est correct}.

\textbf{Conclusion :} pour un achat de crédit à 2 000 FCFA, Armand doit vendre à 2 100 FCFA et gagne 100 FCFA. L'algorithme fonctionne pour n'importe quel prix d'achat saisi : il suffit de relancer l'exécution.
\end{exoresolu}

\begin{attention}[Erreurs fréquentes à éviter]
\begin{itemize}
\item Oublier de déclarer une variable (ou la constante \texttt{TAUX}) avant de l'utiliser.
\item Écrire la marge sous la forme \texttt{5\%} dans un calcul : en pseudo-code, on écrit le coefficient $1{,}05$ ou l'expression \texttt{prixAchat + prixAchat * 5 / 100}.
\item Calculer \texttt{benefice} avant \texttt{prixVente} : l'ordinateur exécute les instructions \cle{dans l'ordre} ; une variable non encore calculée n'a pas de valeur utilisable.
\item Confondre \texttt{Lire} et \texttt{Écrire} : \texttt{Lire} reçoit une donnée de l'utilisateur, \texttt{Écrire} affiche un résultat.
\end{itemize}
\end{attention}

\courssub{Bilan de l'activité}

Cette activité t'a permis de mobiliser, sur une situation commerciale authentique du quotidien camerounais, l'ensemble des savoirs de la leçon :

\begin{itemize}
\item l'\cle{analyse d'un problème} en termes d'entrées, de traitements et de sorties, première étape de toute résolution algorithmique ;
\item la \cle{structure de base d'un algorithme} : en-tête, déclarations, corps entre \texttt{Début} et \texttt{Fin} ;
\item la distinction entre \cle{variable} et \cle{constante}, et le choix justifié des \cle{types de données} ;
\item les trois \cle{instructions de base} : lecture, affectation et écriture, dont l'ordre conditionne la justesse du programme ;
\item la \cle{vérification par un tableau de trace}, méthode simple et puissante pour tester un algorithme « à la main » avant de le faire exécuter par une machine.
\end{itemize}

Tu as également constaté qu'un même raisonnement mathématique (augmenter de 5 \%) devient, une fois traduit en algorithme, un outil réutilisable à l'infini : Armand peut désormais traiter toutes ses ventes de la journée sans risque d'erreur de calcul. C'est exactement cette démarche — observer un besoin réel, le modéliser, rédiger un algorithme, le vérifier puis justifier sa démarche — qui est attendue de toi au BEPC.

\begin{exercice}[Pour aller plus loin]
Reprends l'algorithme de la vente de crédit et modifie-le pour qu'il demande en plus le \textbf{nombre de ventes identiques} réalisées dans la journée, puis qu'il affiche le \textbf{bénéfice total} de la journée. Vérifie ta solution par un tableau de trace avec 3 ventes d'un crédit acheté 1 000 FCFA chacune.
\end{exercice}

\begin{corrige}[Exercice Pour aller plus loin]
\textbf{Énoncé.} Modifier l'algorithme précédent pour prendre en compte le nombre de ventes identiques et afficher le bénéfice total. Tracer pour 3 ventes à 1 000 FCFA l'achat.

\textbf{Solution.}

\textbf{Analyse :}
\begin{itemize}
\item Entrées : \texttt{prixAchat} (Réel), \texttt{nbVentes} (Entier).
\item Traitements : Calcul du prix de vente, du bénéfice unitaire, puis du bénéfice total (\texttt{beneficeTotal $\leftarrow$ beneficeUnitaire * nbVentes}).
\item Sorties : \texttt{prixVente}, \texttt{beneficeUnitaire}, \texttt{beneficeTotal}.
\end{itemize}

\textbf{Algorithme modifié :}

\begin{lstlisting}
Algorithme vente_credit_journee
// Calcule le benefice total d'une journee pour des ventes identiques

Constante
    TAUX : Reel <- 1,05

Variable
    prixAchat      : Reel
    nbVentes       : Entier
    prixVente      : Reel
    beneficeUnitaire : Reel
    beneficeTotal  : Reel

Debut
    Ecrire("Entrez le prix d'achat du credit (FCFA) : ")
    Lire(prixAchat)
    Ecrire("Entrez le nombre de ventes identiques : ")
    Lire(nbVentes)

    prixVente       <- prixAchat * TAUX
    beneficeUnitaire <- prixVente - prixAchat
    beneficeTotal   <- beneficeUnitaire * nbVentes

    Ecrire("Prix de vente unitaire : ", prixVente, " FCFA")
    Ecrire("Benefice unitaire : ", beneficeUnitaire, " FCFA")
    Ecrire("Benefice total de la journee : ", beneficeTotal, " FCFA")
Fin
\end{lstlisting}

\textbf{Tableau de trace (prixAchat = 1 000, nbVentes = 3) :}

\begin{center}
\begin{tabularx}{\linewidth}{|X|l|l|l|l|l|}
\hline
\rowcolor{popblueL} \textbf{Instruction} & \textbf{prixAchat} & \textbf{nbVentes} & \textbf{prixVente} & \textbf{beneficeUnitaire} & \textbf{beneficeTotal} \\
\hline
Début & indéfini & indéfini & indéfini & indéfini & indéfini \\
\hline
Lire(prixAchat) & 1 000 & indéfini & indéfini & indéfini & indéfini \\
\hline
Lire(nbVentes) & 1 000 & 3 & indéfini & indéfini & indéfini \\
\hline
prixVente $\leftarrow$ prixAchat $\times$ TAUX & 1 000 & 3 & 1 050 & indéfini & indéfini \\
\hline
beneficeUnitaire $\leftarrow$ prixVente $-$ prixAchat & 1 000 & 3 & 1 050 & 50 & indéfini \\
\hline
beneficeTotal $\leftarrow$ beneficeUnitaire $\times$ nbVentes & 1 000 & 3 & 1 050 & 50 & 150 \\
\hline
Ecrire(...) & 1 000 & 3 & 1 050 & 50 & 150 \\
\hline
\end{tabularx}
\end{center}

\textbf{Vérification :} $1\,000 \times 1{,}05 = 1\,050$ ; $1\,050 - 1\,000 = 50$ ; $50 \times 3 = 150$. L'algorithme affiche un bénéfice total de 150 FCFA. \textbf{Correct.}
\end{corrige}

\newpage

\coursec{Méthodes et exercices résolus}

\coursec{Présentation du langage de définition des algorithmes}

\courssub{Méthode 1 : Identifier et rédiger la structure de base d'un algorithme}

\begin{methode}[Rédiger la structure de base d'un algorithme]
Pour écrire un algorithme correctement structuré, suis toujours les étapes suivantes :
\begin{enumerate}
  \item \textbf{Nommer l'algorithme} : commence par le mot-clé \texttt{Algorithme} suivi d'un nom simple, sans espace ni accent (ex. : \texttt{Algorithme Calcul\_Moyenne}).
  \item \textbf{Déclarer les constantes et les variables} : sous les mots-clés \texttt{Constantes} puis \texttt{Variables}, liste chaque objet avec son type (\cle{Entier}, \cle{Réel}, \cle{Chaîne}, \cle{Caractère}, \cle{Booléen}). Une constante reçoit sa valeur immédiatement (ex. \texttt{TVA = 19,25}).
  \item \textbf{Ouvrir le corps} par le mot-clé \texttt{Début}.
  \item \textbf{Écrire les instructions} dans l'ordre logique impératif : lecture des données (\texttt{Lire}), traitements (affectations), affichage des résultats (\texttt{Écrire}).
  \item \textbf{Fermer l'algorithme} par le mot-clé \texttt{Fin}.
\end{enumerate}
\textbf{Réflexes à retenir :} chaque variable utilisée doit avoir été déclarée ; les données doivent être lues \emph{avant} d'être utilisées ; le résultat doit être affiché \emph{avant} le \texttt{Fin}.
\end{methode}

\begin{exoresolu}[Structure d'un algorithme de calcul de moyenne]
\textbf{Énoncé.} Le professeur d'informatique du collège de Bafoussam veut un algorithme qui lit les deux notes d'un élève (sur 20) et affiche sa moyenne. Rédige cet algorithme en respectant la structure de base.
\tcblower
\textbf{Solution détaillée.}

\textbf{Étape 1 : analyse du problème.}
\begin{itemize}
  \item Données d'entrée : deux notes, qui peuvent avoir des décimales $\Rightarrow$ type \cle{Réel}.
  \item Traitement : $\text{moyenne} = \dfrac{\text{note1} + \text{note2}}{2}$.
  \item Sortie : la moyenne, de type Réel.
\end{itemize}

\textbf{Étape 2 : rédaction en respectant la structure de base.}
\begin{lstlisting}
Algorithme Calcul_Moyenne
Variables
    note1, note2, moyenne : Reel
Debut
    Ecrire("Entrez la premiere note : ")
    Lire(note1)
    Ecrire("Entrez la deuxieme note : ")
    Lire(note2)
    moyenne <- (note1 + note2) / 2
    Ecrire("La moyenne est : ", moyenne)
Fin
\end{lstlisting}

\textbf{Étape 3 : vérification.} Si \texttt{note1} $= 12$ et \texttt{note2} $= 15$, alors $\text{moyenne} = \dfrac{12+15}{2} = 13,5$. L'algorithme affiche : « La moyenne est : 13,5 ».

\textbf{Conclusion.} L'algorithme respecte les quatre parties de la structure de base : l'en-tête (\texttt{Algorithme}), la déclaration des variables, le corps délimité par \texttt{Debut} et \texttt{Fin}, avec les instructions dans l'ordre lecture $\rightarrow$ traitement $\rightarrow$ écriture.
\end{exoresolu}

\courssub{Méthode 2 : Choisir le bon type de donnée}

\begin{methode}[Déclarer variables et constantes avec le bon type]
\begin{enumerate}
  \item \textbf{Repère la nature de la donnée} dans l'énoncé : nombre entier, nombre à virgule, texte, seul caractère, oui/non.
  \item \textbf{Associe le type} selon le tableau de correspondance :
\end{enumerate}
\begin{center}
\begin{tabularx}{\linewidth}{|l|X|X|}
\hline
\rowcolor{popblueL} \textbf{Type} & \textbf{Nature de la donnée} & \textbf{Exemples} \\
\hline
Entier & Nombre sans virgule & âge, effectif d'une classe, quantité \\
\hline
Réel & Nombre à virgule & prix, moyenne, taille, température \\
\hline
Chaîne & Suite de caractères (texte) & nom, adresse, nom d'une ville \\
\hline
Caractère & Un seul caractère & une lettre, 'O' ou 'N' \\
\hline
Booléen & Vrai ou Faux uniquement & réponse à une question oui/non \\
\hline
\end{tabularx}
\end{center}
\begin{enumerate}
  \setcounter{enumi}{2}
  \item \textbf{Distingue variable et constante} : une \cle{constante} ne change jamais pendant l'exécution (ex. : la TVA de 19,25 \%), une \cle{variable} peut changer de valeur.
  \item \textbf{Déclare} les constantes avec leur valeur dès le départ, les variables avec leur type seulement.
\end{enumerate}
\end{methode}

\begin{exoresolu}[Choix des types pour une boutique]
\textbf{Énoncé.} Mme Fotso tient une boutique à Douala. Elle veut un algorithme qui gère les informations suivantes : le nom d'un article, son prix unitaire, la quantité en stock, et le fait que l'article soit ou non en promotion. Pour chaque information, indique le type de donnée approprié en justifiant, puis écris la partie déclaration de l'algorithme.
\tcblower
\textbf{Solution détaillée.}

\textbf{Analyse de chaque information :}
\begin{itemize}
  \item \textbf{Nom de l'article} : c'est un texte composé de plusieurs caractères (ex. : « Savon ») $\Rightarrow$ type \cle{Chaîne}.
  \item \textbf{Prix unitaire} : c'est un nombre qui peut avoir des décimales (ex. : 250,50 FCFA) $\Rightarrow$ type \cle{Réel}.
  \item \textbf{Quantité en stock} : on compte des articles entiers (on ne vend pas un demi-savon en stock) $\Rightarrow$ type \cle{Entier}.
  \item \textbf{En promotion} : la réponse est seulement Vrai ou Faux $\Rightarrow$ type \cle{Booléen}.
\end{itemize}

\textbf{Partie déclaration de l'algorithme :}
\begin{lstlisting}
Algorithme Gestion_Boutique
Variables
    nomArticle : Chaine
    prixUnitaire : Reel
    quantiteStock : Entier
    enPromotion : Booleen
Debut
    ...
Fin
\end{lstlisting}

\textbf{Conclusion.} Le choix du type dépend de la nature de la donnée : Chaîne pour le texte, Réel pour les nombres à virgule, Entier pour les nombres entiers, Booléen pour les valeurs vrai/faux. Un mauvais choix de type (par exemple Entier pour un prix) entraînerait la perte des décimales.
\end{exoresolu}

\courssub{Méthode 3 : Utiliser les instructions de lecture et d'écriture}

\begin{methode}[Communiquer avec l'utilisateur : Lire et Écrire]
\begin{enumerate}
  \item \textbf{Avant chaque \texttt{Lire}}, affiche un message avec \texttt{Ecrire} pour indiquer à l'utilisateur ce qu'il doit saisir : c'est la règle de convivialité.
  \item \textbf{Syntaxe de lecture} : \texttt{Lire(nom\_variable)} — la valeur tapée au clavier est rangée dans la variable, qui doit avoir été déclarée.
  \item \textbf{Syntaxe d'écriture} : \texttt{Ecrire("message")} pour un texte fixe, \texttt{Ecrire(variable)} pour une valeur, ou les deux combinés : \texttt{Ecrire("Le total est : ", total)}.
  \item \textbf{Ordre obligatoire} : on ne peut afficher ou utiliser une variable qu'\emph{après} lui avoir donné une valeur (par \texttt{Lire} ou par affectation).
  \item \textbf{Vérifie} que le type de la valeur saisie correspond au type déclaré de la variable.
\end{enumerate}
\end{methode}

\begin{exoresolu}[Lecture et écriture : facture d'un cybercafé]
\textbf{Énoncé.} Un cybercafé de Yaoundé facture la connexion Internet à \SI{300}{F} l'heure. Écris un algorithme qui demande le nom du client et le nombre d'heures de connexion, puis affiche le nom du client et le montant à payer.
\tcblower
\textbf{Solution détaillée.}

\textbf{Étape 1 : identification des données.}
\begin{itemize}
  \item Entrées : nom du client (Chaîne), nombre d'heures (Entier).
  \item Constante : le tarif horaire, 300 F, qui ne change pas $\Rightarrow$ constante Entier.
  \item Sortie : montant à payer (Entier, car $300 \times$ heures donne un entier).
\end{itemize}

\textbf{Étape 2 : rédaction.}
\begin{lstlisting}
Algorithme Facture_Cybercafe
Constantes
    TARIF = 300
Variables
    nomClient : Chaine
    heures, montant : Entier
Debut
    Ecrire("Entrez le nom du client : ")
    Lire(nomClient)
    Ecrire("Entrez le nombre d'heures : ")
    Lire(heures)
    montant <- TARIF * heures
    Ecrire("Client : ", nomClient)
    Ecrire("Montant a payer : ", montant, " FCFA")
Fin
\end{lstlisting}

\textbf{Étape 3 : test.} Pour le client « Amina » qui reste 3 heures : $\text{montant} = 300 \times 3 = 900$. L'écran affiche : « Client : Amina » puis « Montant a payer : 900 FCFA ».

\textbf{Conclusion.} Chaque \texttt{Lire} est précédé d'un \texttt{Ecrire} qui guide l'utilisateur, et le résultat est affiché avec un message clair accompagné de l'unité (FCFA).
\end{exoresolu}

\courssub{Méthode 4 : Maîtriser l'affectation et l'ordre des instructions}

\begin{methode}[Utiliser correctement l'affectation]
\begin{enumerate}
  \item \textbf{Syntaxe} : \texttt{variable <- expression}. La flèche va toujours \emph{vers la variable} : on calcule d'abord l'expression à droite, puis on range le résultat dans la variable à gauche.
  \item \textbf{Sens unique} : \texttt{a <- b} copie la valeur de \texttt{b} dans \texttt{a} ; ce n'est pas une égalité mathématique, \texttt{b} ne change pas.
  \item \textbf{Écrasement} : une nouvelle affectation remplace l'ancienne valeur de la variable, qui est définitivement perdue.
  \item \textbf{Incrémentation} : \texttt{compteur <- compteur + 1} signifie « la nouvelle valeur de compteur est l'ancienne valeur plus 1 ».
  \item \textbf{Trace d'exécution} : pour vérifier un algorithme, construis un tableau donnant la valeur de chaque variable après chaque instruction.
\end{enumerate}
\end{methode}

\begin{exoresolu}[Trace d'exécution et échange de deux variables]
\textbf{Énoncé.} On considère l'algorithme suivant :
\begin{lstlisting}
Algorithme Mystere
Variables
    a, b : Entier
Debut
    a <- 5
    b <- 8
    a <- a + b
    b <- a - b
    a <- a - b
    Ecrire(a, b)
Fin
\end{lstlisting}
\begin{enumerate}
  \item Construis la table d'exécution de cet algorithme.
  \item Que réalise cet algorithme ?
\end{enumerate}
\tcblower
\textbf{Solution détaillée.}

\textbf{1. Table d'exécution.} On note la valeur de chaque variable après chaque instruction :

\begin{center}
\begin{tabularx}{\linewidth}{|l|X|X|}
\hline
\rowcolor{popblueL} \textbf{Instruction} & \textbf{Valeur de a} & \textbf{Valeur de b} \\
\hline
\texttt{a <- 5} & 5 & indéfini \\
\hline
\texttt{b <- 8} & 5 & 8 \\
\hline
\texttt{a <- a + b} & $5+8=13$ & 8 \\
\hline
\texttt{b <- a - b} & 13 & $13-8=5$ \\
\hline
\texttt{a <- a - b} & $13-5=8$ & 5 \\
\hline
\end{tabularx}
\end{center}

L'algorithme affiche donc : \texttt{8 5}.

\textbf{2. Rôle de l'algorithme.} Au départ, $a = 5$ et $b = 8$ ; à la fin, $a = 8$ et $b = 5$. Les valeurs de \texttt{a} et \texttt{b} ont été \textbf{échangées}, sans utiliser de variable intermédiaire.

\textbf{Conclusion.} La table d'exécution est l'outil indispensable pour comprendre et vérifier un algorithme : elle montre pas à pas comment l'affectation modifie le contenu des variables. Ici, elle prouve que la suite des trois affectations réalise bien un échange.
\end{exoresolu}

\courssub{Méthode 5 : Rédiger un algorithme complet à partir d'une situation concrète}

\begin{methode}[De la situation de vie à l'algorithme complet]
Pour résoudre un problème concret par un algorithme, applique la démarche en 5 étapes :
\begin{enumerate}
  \item \textbf{Comprendre le problème} : lis l'énoncé et reformule ce qui est demandé.
  \item \textbf{Identifier les données} : distingue les \emph{entrées} (à lire), les \emph{constantes} (valeurs fixes de l'énoncé) et les \emph{sorties} (résultats à afficher).
  \item \textbf{Établir les formules de traitement} : écris les calculs qui transforment les entrées en sorties.
  \item \textbf{Rédiger l'algorithme} en respectant la structure de base : en-tête, déclarations, \texttt{Debut} ... \texttt{Fin}.
  \item \textbf{Tester} avec des valeurs simples : fais une trace d'exécution et vérifie que le résultat affiché est correct.
\end{enumerate}
\textbf{Justification de la démarche :} chaque étape doit pouvoir être expliquée — au BEPC, la démarche compte autant que le résultat final.
\end{methode}

\begin{exoresolu}[Algorithme complet : bulletin d'un élève]
\textbf{Énoncé.} Au collège de Buea, la moyenne d'un élève en informatique est calculée à partir de trois notes sur 20 : deux évaluations et un devoir qui compte double. Écris un algorithme complet qui lit le nom de l'élève et ses trois notes, calcule sa moyenne et affiche le nom de l'élève avec sa moyenne. Justifie ta démarche, puis teste ton algorithme avec les notes 10, 14 et 12.
\tcblower
\textbf{Solution détaillée.}

\textbf{Étape 1 : compréhension.} On veut la moyenne pondérée de trois notes, le devoir comptant double (coefficient 2).

\textbf{Étape 2 : identification des données.}
\begin{itemize}
  \item Entrées : nom de l'élève (Chaîne), note1, note2 (évaluations), noteDevoir (Réels).
  \item Sortie : moyenne (Réel).
\end{itemize}

\textbf{Étape 3 : formule de traitement.} Le devoir compte double, donc :
\[ \text{moyenne} = \frac{\text{note1} + \text{note2} + 2 \times \text{noteDevoir}}{4} \]
Le diviseur est 4 car le total des coefficients est $1 + 1 + 2 = 4$.

\textbf{Étape 4 : rédaction de l'algorithme.}
\begin{lstlisting}
Algorithme Bulletin_Eleve
Variables
    nom : Chaine
    note1, note2, noteDevoir, moyenne : Reel
Debut
    Ecrire("Entrez le nom de l'eleve : ")
    Lire(nom)
    Ecrire("Entrez la note de la premiere evaluation : ")
    Lire(note1)
    Ecrire("Entrez la note de la deuxieme evaluation : ")
    Lire(note2)
    Ecrire("Entrez la note du devoir : ")
    Lire(noteDevoir)
    moyenne <- (note1 + note2 + 2 * noteDevoir) / 4
    Ecrire("Eleve : ", nom)
    Ecrire("Moyenne : ", moyenne, " / 20")
Fin
\end{lstlisting}

\textbf{Étape 5 : test avec les notes 10, 14 et 12.}
\begin{align*}
\text{moyenne} &= \frac{10 + 14 + 2 \times 12}{4} \\
&= \frac{10 + 14 + 24}{4} \\
&= \frac{48}{4} \\
&= 12
\end{align*}
L'algorithme affiche : « Moyenne : 12 / 20 », ce qui est correct.

\textbf{Conclusion.} En suivant la démarche en 5 étapes (comprendre, identifier, modéliser, rédiger, tester), on obtient un algorithme complet, correct et vérifiable, qui répond exactement à la situation posée.
\end{exoresolu}

\courssub{Méthode 6 : Vérifier et corriger un algorithme}

\begin{methode}[Relire et corriger un algorithme]
Avant de rendre ta copie (ou d'exécuter ton algorithme), effectue ces contrôles systématiques :
\begin{enumerate}
  \item \textbf{Structure} : l'algorithme a-t-il un nom ? Les mots-clés \texttt{Debut} et \texttt{Fin} sont-ils présents ?
  \item \textbf{Déclarations} : chaque variable utilisée est-elle déclarée, avec un type cohérent ?
  \item \textbf{Initialisation} : chaque variable a-t-elle reçu une valeur (par \texttt{Lire} ou affectation) avant d'être utilisée dans un calcul ou un affichage ?
  \item \textbf{Ordre des instructions} : les lectures précèdent-elles les traitements, et les traitements précèdent-ils les affichages ?
  \item \textbf{Test} : refais une trace d'exécution avec des valeurs simples et vérifie le résultat à la main.
\end{enumerate}
\end{methode}

\begin{exoresolu}[Corriger un algorithme mal rédigé]
\textbf{Énoncé.} Un élève a écrit l'algorithme suivant pour calculer le périmètre d'un rectangle :
\begin{lstlisting}
Algorithme Perimetre
Debut
    perimetre <- 2 * (longueur + largeur)
    Ecrire("Le perimetre est : ", perimetre)
    Lire(longueur, largeur)
Fin
\end{lstlisting}
Releve toutes les erreurs commises, puis réécris l'algorithme corrigé.
\tcblower
\textbf{Solution détaillée.}

\textbf{Recherche des erreurs :}
\begin{enumerate}
  \item \textbf{Aucune variable n'est déclarée} : \texttt{longueur}, \texttt{largeur} et \texttt{perimetre} sont utilisées sans avoir été déclarées avec leur type.
  \item \textbf{Ordre des instructions incorrect} : le calcul \texttt{perimetre <- 2 * (longueur + largeur)} est effectué \emph{avant} la lecture des valeurs. Au moment du calcul, \texttt{longueur} et \texttt{largeur} n'ont aucune valeur : le calcul est impossible.
  \item \textbf{Absence de messages de saisie} : l'utilisateur ne sait pas ce qu'il doit entrer.
\end{enumerate}

\textbf{Algorithme corrigé :}
\begin{lstlisting}
Algorithme Perimetre_Rectangle
Variables
    longueur, largeur, perimetre : Reel
Debut
    Ecrire("Entrez la longueur : ")
    Lire(longueur)
    Ecrire("Entrez la largeur : ")
    Lire(largeur)
    perimetre <- 2 * (longueur + largeur)
    Ecrire("Le perimetre est : ", perimetre)
Fin
\end{lstlisting}

\textbf{Vérification.} Pour un terrain de longueur $20$ m et de largeur $12$ m : $\text{périmètre} = 2 \times (20 + 12) = 2 \times 32 = 64$ m. L'algorithme affiche « Le perimetre est : 64 », ce qui est exact.

\textbf{Conclusion.} La règle d'or est : \emph{déclarer} toutes les variables, \emph{lire} les données avant de les utiliser, \emph{calculer} ensuite, et \emph{afficher} en dernier.
\end{exoresolu}

\begin{attention}[Les 5 erreurs qui coûtent des points au BEPC]
\begin{enumerate}
  \item \textbf{Oublier de déclarer les variables} (ou déclarer un mauvais type) : toute variable utilisée doit figurer dans la partie \texttt{Variables} avec un type adapté. Mettre un prix en Entier fait perdre les décimales.
  \item \textbf{Utiliser une variable avant de lui avoir donné une valeur} : calculer avec \texttt{longueur} avant de l'avoir lue est une erreur grave. Respecte l'ordre : lecture $\rightarrow$ traitement $\rightarrow$ écriture.
  \item \textbf{Inverser le sens de l'affectation} : on écrit \texttt{moyenne <- (note1 + note2) / 2} et jamais \texttt{(note1 + note2) / 2 <- moyenne}. La variable recevant le résultat est toujours à \emph{gauche} de la flèche.
  \item \textbf{Oublier les mots-clés de structure} : un algorithme sans \texttt{Algorithme} (nom), \texttt{Debut} ou \texttt{Fin} est incomplet et sanctionné, même si les calculs sont justes.
  \item \textbf{Afficher un résultat brut sans message} : écrire seulement \texttt{Ecrire(moyenne)} rend l'affichage incompréhensible. Accompagne toujours le résultat d'un texte clair et de l'unité : \texttt{Ecrire("La moyenne est : ", moyenne, " / 20")}.
\end{enumerate}
\end{attention}

\newpage

\coursec{Exercices}

\coursec{Rappels et notion de langage de définition des algorithmes}

\begin{savaistu}[Origine du mot « algorithme »]
Le terme \cle{algorithme} vient du nom du mathématicien perse \textbf{Al-Khwârizmî} (vers 780--850), qui a écrit un traité sur le calcul indien et l'algèbre. Son nom, latinisé en \emph{Algoritmi}, a donné le mot « algorithme » : une méthode de calcul systématique, étape par étape, qui aboutit à un résultat garanti.
\end{savaistu}

\begin{definition}[Algorithme]
Un \textbf{algorithme} est une suite \textbf{finie} et \textbf{ordonnée} d'instructions \textbf{précises} et \textbf{non ambiguës}, permettant de résoudre un problème ou d'obtenir un résultat à partir de données d'entrée.
\end{definition}

\begin{aretenir}[Caractéristiques d'un algorithme]
Pour être valide, un algorithme doit respecter cinq propriétés fondamentales :
\begin{enumerate}
    \item \textbf{Finitude} : il se termine après un nombre fini d'étapes.
    \item \textbf{Précision (déterminisme)} : chaque instruction est claire et ne laisse place à aucune interprétation.
    \item \textbf{Entrées} : zéro ou plusieurs données initiales (fournies par l'utilisateur ou lues dans un fichier).
    \item \textbf{Sorties} : une ou plusieurs données résultantes (affichées, stockées ou transmises).
    \item \textbf{Efficacité} : chaque instruction doit être réalisable en un temps fini avec les moyens disponibles (pas d'opérations impossibles comme « diviser par zéro » ou « deviner un nombre »).
\end{enumerate}
\end{aretenir}

\begin{propriete}[Langage de définition d'algorithmes (Pseudo-code)]
Au lieu d'écrire directement dans un langage de programmation (C, Python, JavaScript\ldots), on utilise un \textbf{langage de définition d'algorithmes} (souvent appelé \emph{pseudo-code}). Ce langage :
\begin{itemize}
    \item utilise des mots-clés en \textbf{français} (ou en anglais selon les conventions) : \texttt{Algorithme}, \texttt{Variable}, \texttt{Début}, \texttt{Fin}, \texttt{Si}, \texttt{TantQue}, \texttt{Lire}, \texttt{Écrire}\ldots
    \item est \textbf{indépendant de la machine} et du compilateur ;
    \item permet de se concentrer sur la \textbf{logique} et la \textbf{structure} du programme sans se soucier de la syntaxe stricte d'un langage réel (points-virgules, accolades, includes\ldots).
\end{itemize}
C'est le langage que nous utiliserons tout au long de l'année pour concevoir nos programmes avant de les coder en C ou en HTML/JavaScript.
\end{propriete}

\coursec{Structure de base d'un algorithme}

\begin{propriete}[Les trois parties obligatoires]
Tout algorithme écrit en pseudo-code se compose de trois grandes parties, dans cet ordre strict :
\begin{enumerate}
    \item \textbf{L'en-tête} : commence par le mot-clé \texttt{Algorithme} suivi du \textbf{nom} de l'algorithme (sans espace, commençant par une majuscule, ex. : \texttt{CalculMoyenne}).
    \item \textbf{La partie déclarative} : introduite par le mot-clé \texttt{Variables} (ou \texttt{Variable}), elle liste toutes les variables et constantes utilisées avec leur \textbf{type} (Entier, Réel, Chaîne, Booléen). \emph{Optionnel si aucune variable n'est utilisée.}
    \item \textbf{Le corps de l'algorithme} : délimité par \texttt{Début} et \texttt{Fin}, il contient la suite des \textbf{instructions} exécutées séquentiellement.
\end{enumerate}
\end{propriete}

\begin{popfigure}
\begin{tikzpicture}[
    node distance=1.2cm and 2cm,
    box/.style={rectangle, draw=popdark, thick, rounded corners, fill=popcream, minimum height=1.2cm, text width=4.5cm, align=center, font=\small},
    arrow/.style={-Stealth, thick, popdark}
]
\node[box, align=center] (entete){\textbf{En-tête}\\\texttt{Algorithme MonAlgo}};
\node[box, below=of entete, align=center] (decl){\textbf{Déclarations}\\\texttt{Variables}\\ \texttt{a, b : Entier}\\ \texttt{c : Réel}};
\node[box, below=of decl, align=center] (corps){\textbf{Corps}\\\texttt{Début}\\ \texttt{\quad Instructions\ldots}\\ \texttt{Fin}};
\draw[arrow] (entete) -- (decl);
\draw[arrow] (decl) -- (corps);
\end{tikzpicture}
\legende{Structure générale d'un algorithme en pseudo-code}
\end{popfigure}

\begin{exemplebox}[Squelette minimal]
Voici la structure la plus simple possible (aucun traitement, aucune variable) :
\begin{lstlisting}
Algorithme Vide
Début
    (* Aucune instruction *)
Fin
\end{lstlisting}
\end{exemplebox}

\begin{attention}[Erreurs fréquentes sur la structure]
\begin{itemize}
    \item Oublier \texttt{Début} ou \texttt{Fin}.
    \item Écrire des instructions \textbf{en dehors} du bloc \texttt{Début ... Fin} (par exemple après \texttt{Fin}).
    \item Mettre les déclarations \textbf{après} \texttt{Début} : les déclarations précèdent toujours \texttt{Début}.
    \item Utiliser des accents ou des espaces dans le nom de l'algorithme (ex. \texttt{Algorithme Calcul Moyenne} $\rightarrow$ \textbf{incorrect} ; écrire \texttt{CalculMoyenne}).
\end{itemize}
\end{attention}

\coursec{Variables, constantes et types de données}

\begin{definition}[Variable]
Une \textbf{variable} est une \textbf{zone de mémoire} identifiée par un \textbf{nom} (identificateur), qui contient une \textbf{valeur unique} susceptible de \textbf{changer} au cours de l'exécution de l'algorithme.
\end{definition}

\begin{definition}[Constante]
Une \textbf{constante} est une zone de mémoire identifiée par un nom, dont la valeur est \textbf{fixée une seule fois} à la déclaration et \textbf{ne peut jamais être modifiée} pendant l'exécution.
\end{definition}

\begin{propriete}[Types de données de base (niveau 3ème)]
En algorithmique au collège, on utilise quatre types fondamentaux :
\begin{center}
\begin{tabularx}{\linewidth}{|l|X|l|}
\hline
\rowcolor{popblueL} \textbf{Type} & \textbf{Description / Valeurs admises} & \textbf{Exemple de littéral} \\
\hline
\texttt{Entier} & Nombres entiers relatifs ($\ldots, -2, -1, 0, 1, 2, \ldots$) & $42$, $-7$, $0$ \\
\hline
\texttt{Réel} & Nombres à virgule (notation décimale avec \textbf{point} ou virgule selon convention, ici virgule française) & $3,14$, $-0,5$, $12,0$ \\
\hline
\texttt{Chaîne} & Suite de caractères (lettres, chiffres, symboles) entre guillemets & \texttt{"Bonjour"}, \texttt{"3ème C"} \\
\hline
\texttt{Booléen} & Valeur logique : \texttt{Vrai} ou \texttt{Faux} & \texttt{Vrai}, \texttt{Faux} \\
\hline
\end{tabularx}
\end{center}
\end{propriete}

\begin{propriete}[Syntaxe de déclaration]
\begin{lstlisting}
Variables
    <nom1>, <nom2> : <Type> ;
    <nom3> : <Type> ;
Constantes
    <NOM_CONST> = <valeur> : <Type> ;
\end{lstlisting}
\begin{itemize}
    \item Les noms de variables commencent par une \textbf{minuscule} (ex. \texttt{note}, \texttt{prixUnitaire}).
    \item Les noms de constantes s'écrivent en \textbf{MAJUSCULES} (ex. \texttt{TAUX\_TVA}, \texttt{ANNEE\_COURANTE}).
    \item On sépare les variables de même type par des virgules ; on termine chaque ligne par un point-virgule.
\end{itemize}
\end{propriete}

\begin{exemplebox}[Déclarations variées]
\begin{lstlisting}
Variables
    age, nbEleves : Entier ;
    moyenne, prix : Réel ;
    nomEleve, ville : Chaîne ;
    estAdmis : Booléen ;
Constantes
    ANNEE_COURANTE = 2026 : Entier ;
    PI = 3,14159 : Réel ;
    NOM_ETAB = "Lycee Bilingue" : Chaîne ;
\end{lstlisting}
\end{exemplebox}

\begin{attention}[Règles de nommage (identificateurs)]
\begin{itemize}
    \item Commencer par une \textbf{lettre} (pas de chiffre au début).
    \item Pas d'espaces, pas d'accents, pas de caractères spéciaux (\%, \$, \#, @\ldots). Le trait de soulignement \texttt{\_} est autorisé.
    \item Respecter la casse : \texttt{maVariable} $\neq$ \texttt{mavariable}.
    \item Choisir des noms \textbf{significatifs} : \texttt{noteMaths} est mieux que \texttt{n} ou \texttt{x}.
\end{itemize}
\end{attention}

\coursec{Les instructions de base : lecture, écriture, affectation}

\begin{propriete}[Instruction d'affectation : $\leftarrow$]
L'affectation stocke une valeur dans une variable.
\begin{center}
\texttt{<Variable> $\leftarrow$ <Expression>}
\end{center}
\begin{itemize}
    \item La flèche $\leftarrow$ (saisie \texttt{<-} au clavier) signifie « \textbf{reçoit la valeur de} ».
    \item L'\texttt{<Expression>} peut être une valeur littérale, une variable, ou un calcul ($a + b * 2$, \texttt{Racine(16)}\ldots).
    \item \textbf{Sens unique} : la variable \textbf{cible} est \textbf{à gauche}, la valeur source \textbf{à droite}.
\end{itemize}
\end{propriete}

\begin{propriete}[Instruction de lecture : \texttt{Lire}]
Permet de saisir une valeur au clavier pendant l'exécution et de la stocker dans une variable.
\begin{center}
\texttt{Lire(<Variable>)}
\end{center}
\begin{itemize}
    \item La variable doit être \textbf{déclarée} et de \textbf{type compatible} avec la saisie attendue.
    \item L'exécution \textbf{suspend} l'algorithme jusqu'à la validation (touche Entrée).
\end{itemize}
\end{propriete}

\begin{propriete}[Instruction d'écriture : \texttt{Écrire}]
Affiche une information à l'écran (sortie standard).
\begin{center}
\texttt{Écrire(<Expression>)} \quad ou \quad \texttt{Écrire(<Message>, <Variable>)}
\end{center}
\begin{itemize}
    \item Les chaînes de caractères (messages) sont entre \textbf{guillemets doubles} : \texttt{"Total : "}.
    \item On peut enchaîner plusieurs arguments séparés par des virgules (selon les dialectes) ou faire plusieurs \texttt{Écrire} successifs. Nous utiliserons plusieurs \texttt{Écrire} pour la clarté.
    \item \texttt{Écrire} n'affiche \textbf{pas} de retour à la ligne automatique à la fin (sauf convention locale) ; on gère l'affichage ligne par ligne.
\end{itemize}
\end{propriete}

\begin{methode}[Rédiger un algorithme simple (3 étapes)]
\begin{enumerate}
    \item \textbf{Analyser le problème} : identifier les \textbf{entrées} (données connues ou demandées), les \textbf{sorties} (résultats attendus) et les \textbf{traitements} (formules, calculs).
    \item \textbf{Concevoir} : choisir les variables (noms, types), les constantes, et ordonner les instructions logiquement (Lire $\rightarrow$ Calculer $\rightarrow$ Écrire).
    \item \textbf{Coder en pseudo-code} : respecter la structure \texttt{Algorithme / Variables / Début / Fin} et la syntaxe des trois instructions de base.
\end{enumerate}
\end{methode}

\begin{exemplebox}[Algorithme complet : Calcul du prix TTC]
Problème : Saisir un prix HT, appliquer un taux de TVA constant (19,25\%), afficher le prix TTC.
\begin{lstlisting}
Algorithme PrixTTC
Variables
    prixHT, prixTTC : Réel ;
Constantes
    TAUX_TVA = 0,1925 : Réel ;
Début
    Écrire("Entrez le prix HT (FCFA) : ")
    Lire(prixHT)
    prixTTC <- prixHT * (1 + TAUX_TVA)
    Écrire("Le prix TTC est : ")
    Écrire(prixTTC)
    Écrire(" FCFA")
Fin
\end{lstlisting}
\end{exemplebox}

\coursec{Rédaction complète d'algorithmes sur des situations concrètes}

\begin{exoresolu}[Calcul de la moyenne de trois notes]
\textbf{Énoncé :} Écrire un algorithme \texttt{Moyenne3Notes} qui lit trois notes sur 20, calcule leur moyenne et l'affiche avec un message.

\textbf{Analyse :}
\begin{itemize}
    \item Entrées : trois notes (réelles, ex. 12,5 ; 14 ; 9,75).
    \item Sortie : la moyenne (réelle).
    \item Traitement : somme des notes divisée par 3.
\end{itemize}

\textbf{Algorithme :}
\begin{lstlisting}
Algorithme Moyenne3Notes
Variables
    note1, note2, note3, somme, moyenne : Réel ;
Début
    Écrire("Entrez la note 1 : ")
    Lire(note1)
    Écrire("Entrez la note 2 : ")
    Lire(note2)
    Écrire("Entrez la note 3 : ")
    Lire(note3)
    somme <- note1 + note2 + note3
    moyenne <- somme / 3
    Écrire("La moyenne des trois notes est : ")
    Écrire(moyenne)
    Écrire(" / 20")
Fin
\end{lstlisting}

\textbf{Test (trace) :} Si notes = 12, 14, 10 $\rightarrow$ somme = 36 $\rightarrow$ moyenne = 12. Affichage : \texttt{La moyenne des trois notes est : 12 / 20}.
\tcblower
\textbf{Remarques pédagogiques :}
\begin{itemize}
    \item Le type \texttt{Réel} est indispensable car les notes peuvent avoir des demi-points.
    \item La division par 3 (entier) est une division réelle car \texttt{somme} est réel.
    \item L'affichage est découpé pour être lisible.
\end{itemize}
\end{exoresolu}

\courssub{Je vérifie mes connaissances}

\begin{exercice}[QCM : à chaque question, une seule bonne réponse]
Pour chaque question, recopie la lettre de la bonne réponse.
\begin{enumerate}
\item Un algorithme commence toujours par le mot-clé :
\begin{enumerate}
\item[a)] \texttt{Fin}
\item[b)] \texttt{Début}
\item[c)] \texttt{Variable}
\item[d)] \texttt{Écrire}
\end{enumerate}
\item L'instruction qui permet à l'utilisateur d'entrer une valeur est :
\begin{enumerate}
\item[a)] \texttt{Écrire}
\item[b)] \texttt{Afficher}
\item[c)] \texttt{Lire}
\item[d)] \texttt{Calculer}
\end{enumerate}
\item Le symbole $\leftarrow$ (saisi \texttt{<-}) signifie :
\begin{enumerate}
\item[a)] une comparaison
\item[b)] une affectation
\item[c)] une lecture
\item[d)] une fin d'algorithme
\end{enumerate}
\item La valeur d'une constante au cours de l'exécution d'un algorithme :
\begin{enumerate}
\item[a)] change à chaque étape
\item[b)] ne change jamais
\item[c)] change seulement à la fin
\item[d)] est toujours nulle
\end{enumerate}
\end{enumerate}
\end{exercice}

\begin{exercice}[Vrai ou faux ? Justifie ta réponse]
Réponds par \textbf{vrai} ou \textbf{faux}, puis justifie en une phrase.
\begin{enumerate}
\item Une variable doit être déclarée avant d'être utilisée dans un algorithme.
\item L'instruction \texttt{Écrire("Bonjour")} permet de saisir le prénom de l'utilisateur.
\item Une variable de type \texttt{Entier} peut contenir la valeur $12,75$.
\item Dans l'instruction \texttt{A $\leftarrow$ B}, c'est la variable \texttt{A} qui reçoit la valeur de \texttt{B}.
\end{enumerate}
\end{exercice}

\begin{exercice}[Définitions]
Définis en une ou deux phrases chacun des termes suivants, puis donne un exemple pour chacun :
\begin{enumerate}
\item Algorithme ;
\item Variable ;
\item Constante ;
\item Affectation.
\end{enumerate}
\end{exercice}

\begin{exercice}[Rôles des instructions de base]
Recopie et complète le tableau ci-dessous en associant à chaque instruction son rôle, puis donne pour chacune un exemple d'utilisation dans un algorithme de la vie courante (cantine scolaire, marché, transport en moto-taxi\ldots).

\begin{center}
\begin{tabularx}{\linewidth}{|l|X|X|}
\hline
\rowcolor{popblueL} \textbf{Instruction} & \textbf{Rôle} & \textbf{Exemple d'utilisation} \\
\hline
\texttt{Lire} & & \\
\hline
\texttt{Écrire} & & \\
\hline
$\leftarrow$ (affectation) & & \\
\hline
\end{tabularx}
\end{center}
\end{exercice}

\courssub{Je m'entraîne}

\begin{exercice}[Au marché de Mokolo]
Écris en pseudo-code un algorithme nommé \texttt{TotalAchat} qui :
\begin{enumerate}
\item lit le prix unitaire (en FCFA) et la quantité d'un article acheté au marché ;
\item calcule le montant total à payer ;
\item affiche le montant total avec un message clair.
\end{enumerate}
Tu déclareras toutes les variables utilisées en précisant leur type.
\end{exercice}

\begin{exercice}[Année de naissance]
Écris en pseudo-code un algorithme nommé \texttt{Naissance} qui lit l'âge d'un élève de 3ème et affiche son année de naissance, sachant que l'année courante est 2026. Tu utiliseras une constante pour l'année courante et tu justifieras ce choix.
\end{exercice}

\begin{exercice}[Périmètre et aire d'un carré]
Le terrain de sport de ton collège a la forme d'un carré.
\begin{enumerate}
\item Écris en pseudo-code un algorithme qui lit la longueur du côté du terrain (en mètres), puis calcule et affiche son périmètre et son aire.
\item Justifie le choix du type (\texttt{Entier} ou \texttt{Réel}) de chaque variable utilisée.
\end{enumerate}
\end{exercice}

\begin{exercice}[Tableau de trace]
On considère l'algorithme suivant :

\begin{lstlisting}
Algorithme Mystere
Variable a, b, c : Entier
Début
    a <- 5
    b <- 3
    c <- a + b
    a <- c - b
    Écrire(a)
    Écrire(c)
Fin
\end{lstlisting}

\begin{enumerate}
\item Recopie et complète le tableau de trace ci-dessous en donnant la valeur de chaque variable après l'exécution de chaque ligne.

\begin{center}
\begin{tabularx}{\linewidth}{|l|X|X|X|}
\hline
\rowcolor{popblueL} \textbf{Instruction exécutée} & \textbf{a} & \textbf{b} & \textbf{c} \\
\hline
\texttt{a <- 5} & & & \\
\hline
\texttt{b <- 3} & & & \\
\hline
\texttt{c <- a + b} & & & \\
\hline
\texttt{a <- c - b} & & & \\
\hline
\end{tabularx}
\end{center}

\item Déduis-en les valeurs affichées par l'algorithme à la fin de son exécution.
\end{enumerate}
\end{exercice}

\courssub{Je me prépare au BEPC}

\begin{exercice}[Type BEPC : la chasse aux erreurs]
Marius, élève en classe de 3ème au collège bilingue de Bafoussam, a rédigé l'algorithme ci-dessous pour calculer la moyenne de deux notes d'un camarade. Malheureusement, son algorithme contient \textbf{trois erreurs} : une variable non déclarée, un mauvais type de variable et une affectation inversée (sens de la flèche).

\begin{lstlisting}
Algorithme MoyenneNotes
Variable note1, note2 : Entier
Variable moyenne : Entier
Début
    Écrire("Entrez la première note : ")
    Lire(note1)
    Écrire("Entrez la deuxième note : ")
    Lire(note2)
    note1 + note2 <- somme
    moyenne <- somme / 2
    Écrire("La moyenne est : ")
    Écrire(moyenne)
Fin
\end{lstlisting}

\begin{enumerate}
\item Recopie la structure de base d'un algorithme (les trois grandes parties) en les nommant.
\item Identifie les trois erreurs commises par Marius en précisant la ligne concernée et la nature de chaque erreur.
\item Réécris l'algorithme complet corrigé, en choisissant des types de variables adaptés (une moyenne peut être un nombre à virgule, par exemple $12,5$).
\item Donne la valeur affichée par l'algorithme corrigé si l'utilisateur entre les notes $13$ et $17$.
\end{enumerate}
\end{exercice}

\begin{exercice}[Type BEPC : la boutique de Mme Etoundi]
Mme Etoundi tient une boutique d'articles scolaires à Yaoundé. Pour la rentrée, elle souhaite disposer d'un programme simple qui l'aide à calculer la facture d'un client. Une règle de sa boutique est fixe : le prix d'un paquet de 10 stylos est de \num{1000} FCFA, et ce prix ne change jamais au cours de l'année.

\begin{enumerate}
\item Explique pourquoi le prix du paquet de stylos doit être représenté par une \textbf{constante} et non par une variable, puis donne la ligne de déclaration correspondante en pseudo-code.
\item Écris en pseudo-code un algorithme nommé \texttt{Facture} qui :
\begin{itemize}
\item lit le nombre de paquets de stylos achetés par le client ;
\item lit le montant des autres articles achetés (cahiers, règles\ldots) ;
\item calcule le montant total de la facture ;
\item affiche le montant total à payer avec un message poli du type : \og Merci, vous devez payer \ldots{} FCFA \fg.
\end{itemize}
\item Mme Etoundi vend $4$ paquets de stylos et pour \num{2500} FCFA d'autres articles à un client. Fais le tableau de trace de ton algorithme et donne le message exact affiché à l'écran.
\item Rédige une conclusion de deux phrases expliquant l'intérêt d'un tel algorithme pour la gestion quotidienne de la boutique.
\end{enumerate}
\end{exercice}

\begin{exercice}[Type BEPC : le transport scolaire]
Le collège de Nkongsamba organise le transport de ses élèves par cars. Chaque car peut transporter au maximum $50$ élèves. Le proviseur demande à la classe de 3ème de concevoir un algorithme qui, à partir du nombre d'élèves inscrits à une sortie pédagogique, affiche le nombre de cars complets et le nombre d'élèves restant dans le dernier car. On rappelle que :
\begin{itemize}
\item \texttt{div} désigne la division entière : $137 \textbf{ div } 50 = 2$ ;
\item \texttt{mod} désigne le reste de la division entière : $137 \textbf{ mod } 50 = 37$.
\end{itemize}

\begin{enumerate}
\item Donne la définition d'un algorithme et cite les trois instructions de base étudiées en classe.
\item Identifie les données d'entrée et les données de sortie du problème posé, puis propose un nom et un type pour chaque variable (le nombre de cars et le nombre d'élèves restants sont des nombres entiers).
\item Écris en pseudo-code l'algorithme \texttt{Transport} qui lit le nombre d'élèves inscrits, calcule le nombre de cars complets et le nombre d'élèves restants, puis affiche les deux résultats avec des messages explicites.
\item Vérifie ton algorithme en faisant son tableau de trace pour $137$ élèves inscrits, puis pour $100$ élèves inscrits. Que remarques-tu dans le second cas ?
\item Justifie en une phrase le choix du type \texttt{Entier} pour toutes les variables de cet algorithme.
\end{enumerate}
\end{exercice}

\newpage

\coursec{Corrigés des exercices}

\begin{corrige}[Exercice 1 — QCM : à chaque question, une seule bonne réponse]
\begin{enumerate}
\item \textbf{Réponse b)} : \texttt{Début}. Le corps de l'algorithme, qui contient les instructions, commence toujours par le mot-clé \texttt{Début} et se termine par \texttt{Fin}. (\texttt{Algorithme} ouvre l'en-tête, mais le \emph{corps} commence par \texttt{Début}.)
\item \textbf{Réponse c)} : \texttt{Lire}. C'est l'instruction de lecture (saisie au clavier) : \texttt{Lire(variable)}.
\item \textbf{Réponse b)} : une affectation. \texttt{A <- B} signifie « \texttt{A} reçoit la valeur de \texttt{B} ».
\item \textbf{Réponse b)} : ne change jamais. Par définition, la valeur d'une constante est fixée une seule fois à la déclaration et ne peut pas être modifiée pendant l'exécution.
\end{enumerate}
\end{corrige}

\begin{corrige}[Exercice 2 — Vrai ou faux ? Justifie ta réponse]
\begin{enumerate}
\item \textbf{Vrai.} Toute variable doit être déclarée (avec son type) dans la partie déclarative avant le mot-clé \texttt{Début} ; sinon l'algorithme ne sait pas quel type de valeur elle contient.
\item \textbf{Faux.} \texttt{Écrire("Bonjour")} est une instruction d'\emph{affichage} : elle affiche un message à l'écran. Pour saisir le prénom de l'utilisateur, il faut utiliser \texttt{Lire(prenom)}.
\item \textbf{Faux.} Une variable de type \texttt{Entier} ne peut contenir que des nombres entiers ($\ldots, -1, 0, 1, 2, \ldots$). La valeur $12,75$ possède une partie décimale : il faut un type \texttt{Réel}.
\item \textbf{Vrai.} Dans une affectation, la variable cible est toujours à \textbf{gauche} de la flèche et la valeur source à \textbf{droite} : \texttt{A <- B} place la valeur de \texttt{B} dans \texttt{A}.
\end{enumerate}
\end{corrige}

\begin{corrige}[Exercice 3 — Définitions]
\begin{enumerate}
\item \textbf{Algorithme} : suite finie et ordonnée d'instructions précises et non ambiguës permettant de résoudre un problème. \emph{Exemple :} l'algorithme \texttt{Moyenne3Notes} qui calcule la moyenne de trois notes.
\item \textbf{Variable} : zone de mémoire identifiée par un nom, contenant une valeur qui peut changer au cours de l'exécution. \emph{Exemple :} \texttt{note1 : Réel}.
\item \textbf{Constante} : zone de mémoire identifiée par un nom, dont la valeur est fixée à la déclaration et ne change jamais pendant l'exécution. \emph{Exemple :} \texttt{ANNEE\_COURANTE = 2026 : Entier}.
\item \textbf{Affectation} : instruction qui stocke une valeur dans une variable, notée \texttt{Variable <- Expression}. \emph{Exemple :} \texttt{total <- prix * quantite}.
\end{enumerate}
\end{corrige}

\begin{corrige}[Exercice 4 — Rôles des instructions de base]
\begin{center}
\begin{tabularx}{\linewidth}{|l|X|X|}
\hline
\rowcolor{popblueL} \textbf{Instruction} & \textbf{Rôle} & \textbf{Exemple d'utilisation} \\
\hline
\texttt{Lire} & Saisir une valeur au clavier et la stocker dans une variable & \texttt{Lire(prix)} : le commerçant entre le prix d'un article au marché \\
\hline
\texttt{Écrire} & Afficher un message ou un résultat à l'écran & \texttt{Écrire("Total à payer : ")} à la cantine scolaire \\
\hline
$\leftarrow$ (affectation) & Stocker une valeur ou le résultat d'un calcul dans une variable & \texttt{total <- prix * quantite} pour calculer le montant d'un achat \\
\hline
\end{tabularx}
\end{center}
\end{corrige}

\begin{corrige}[Exercice 5 — Au marché de Mokolo]
\textbf{Analyse :} entrées : prix unitaire et quantité ; sortie : montant total ; traitement : \texttt{total <- prixUnitaire * quantite}.
\begin{lstlisting}
Algorithme TotalAchat
Variables
    prixUnitaire : Reel ;
    quantite : Entier ;
    total : Reel ;
Début
    Écrire("Entrez le prix unitaire (FCFA) : ")
    Lire(prixUnitaire)
    Écrire("Entrez la quantité : ")
    Lire(quantite)
    total <- prixUnitaire * quantite
    Écrire("Le montant total à payer est : ")
    Écrire(total)
    Écrire(" FCFA")
Fin
\end{lstlisting}
\textbf{Justification des types :} \texttt{prixUnitaire} et \texttt{total} sont de type \texttt{Réel} car un prix peut avoir des centimes ; \texttt{quantite} est un \texttt{Entier} car on compte des articles entiers.
\end{corrige}

\begin{corrige}[Exercice 6 — Année de naissance]
\begin{lstlisting}
Algorithme Naissance
Variables
    age, anneeNaissance : Entier ;
Constantes
    ANNEE_COURANTE = 2026 : Entier ;
Début
    Écrire("Entrez votre âge : ")
    Lire(age)
    anneeNaissance <- ANNEE_COURANTE - age
    Écrire("Vous êtes né(e) en : ")
    Écrire(anneeNaissance)
Fin
\end{lstlisting}
\textbf{Justification :} l'année courante est déclarée en \textbf{constante} car elle est fixe pendant toute l'exécution (2026) et ne doit pas être modifiée par l'algorithme ; l'écrire en constante évite toute modification accidentelle et rend l'algorithme plus lisible.
\end{corrige}

\begin{corrige}[Exercice 7 — Périmètre et aire d'un carré]
\begin{enumerate}
\item Algorithme :
\begin{lstlisting}
Algorithme TerrainCarre
Variables
    cote, perimetre, aire : Reel ;
Début
    Écrire("Entrez la longueur du côté (m) : ")
    Lire(cote)
    perimetre <- 4 * cote
    aire <- cote * cote
    Écrire("Le périmètre est : ")
    Écrire(perimetre)
    Écrire(" m")
    Écrire("L'aire est : ")
    Écrire(aire)
    Écrire(" m²")
Fin
\end{lstlisting}
\item \textbf{Justification des types :} le type \texttt{Réel} est choisi pour les trois variables car la longueur du côté peut être un nombre à virgule (par exemple $25,5$ m), et le périmètre et l'aire, calculés à partir d'elle, peuvent donc aussi être décimaux. Si l'on était sûr que le côté est toujours entier, on pourrait utiliser \texttt{Entier}.
\end{enumerate}
\end{corrige}

\begin{corrige}[Exercice 8 — Tableau de trace]
\begin{enumerate}
\item Tableau de trace complété :
\begin{center}
\begin{tabularx}{\linewidth}{|l|X|X|X|}
\hline
\rowcolor{popblueL} \textbf{Instruction exécutée} & \textbf{a} & \textbf{b} & \textbf{c} \\
\hline
\texttt{a <- 5} & 5 & --- & --- \\
\hline
\texttt{b <- 3} & 5 & 3 & --- \\
\hline
\texttt{c <- a + b} & 5 & 3 & 8 \\
\hline
\texttt{a <- c - b} & 5 & 3 & 8 \\
\hline
\end{tabularx}
\end{center}
Détail du dernier calcul : \texttt{a <- c - b} donne $a = 8 - 3 = 5$. La nouvelle valeur de \texttt{a} est donc $5$ (elle remplace l'ancienne, qui était déjà $5$).
\item \textbf{Valeurs affichées :} \texttt{Écrire(a)} affiche $5$, puis \texttt{Écrire(c)} affiche $8$.
\end{enumerate}
\end{corrige}

\begin{corrige}[Exercice 9 — Type BEPC : la chasse aux erreurs]
\begin{enumerate}
\item Les trois grandes parties d'un algorithme : \textbf{l'en-tête} (\texttt{Algorithme NomAlgo}), \textbf{la partie déclarative} (\texttt{Variables} / \texttt{Constantes}) et \textbf{le corps} (entre \texttt{Début} et \texttt{Fin}).
\item Les trois erreurs :
\begin{itemize}
\item Ligne \texttt{note1 + note2 <- somme} : \textbf{affectation inversée} — la variable cible doit être à gauche ; on doit écrire \texttt{somme <- note1 + note2}.
\item Ligne \texttt{moyenne <- somme / 2} : la variable \texttt{somme} est \textbf{non déclarée} dans la partie déclarative.
\item Ligne \texttt{Variable moyenne : Entier} : \textbf{mauvais type} — une moyenne peut être un nombre à virgule (ex. $12,5$), il faut le type \texttt{Réel}.
\end{itemize}
\item Algorithme corrigé :
\begin{lstlisting}
Algorithme MoyenneNotes
Variables
    note1, note2 : Reel ;
    somme, moyenne : Reel ;
Début
    Écrire("Entrez la première note : ")
    Lire(note1)
    Écrire("Entrez la deuxième note : ")
    Lire(note2)
    somme <- note1 + note2
    moyenne <- somme / 2
    Écrire("La moyenne est : ")
    Écrire(moyenne)
Fin
\end{lstlisting}
\item Avec les notes $13$ et $17$ : somme $= 13 + 17 = 30$ ; moyenne $= 30 / 2 = 15$. L'algorithme affiche : \texttt{La moyenne est : 15}.
\end{enumerate}
\textbf{Barème indicatif (10 pts) :} structure (2 pts) ; trois erreurs identifiées avec nature (3 pts) ; algorithme corrigé complet (3 pts) ; valeur affichée justifiée (2 pts).
\textbf{Points de vigilance :} citer explicitement la règle « la cible de l'affectation est à gauche de la flèche » ; déclarer \texttt{somme} ; justifier le type \texttt{Réel} par l'existence possible de demi-points.
\end{corrige}

\begin{corrige}[Exercice 10 — Type BEPC : la boutique de Mme Etoundi]
\begin{enumerate}
\item Le prix du paquet de stylos (1000 FCFA) ne change jamais au cours de l'année : c'est une valeur fixe, donc une \textbf{constante}. La déclarer en constante garantit qu'aucune instruction ne pourra la modifier par erreur.
\begin{lstlisting}
Constantes
    PRIX_PAQUET = 1000 : Entier ;
\end{lstlisting}
\item Algorithme :
\begin{lstlisting}
Algorithme Facture
Variables
    nbPaquets : Entier ;
    autresArticles, total : Reel ;
Constantes
    PRIX_PAQUET = 1000 : Entier ;
Début
    Écrire("Nombre de paquets de stylos : ")
    Lire(nbPaquets)
    Écrire("Montant des autres articles (FCFA) : ")
    Lire(autresArticles)
    total <- nbPaquets * PRIX_PAQUET + autresArticles
    Écrire("Merci, vous devez payer ")
    Écrire(total)
    Écrire(" FCFA")
Fin
\end{lstlisting}
\item Tableau de trace pour $4$ paquets et 2500 FCFA d'autres articles :
\begin{center}
\begin{tabularx}{\linewidth}{|l|X|X|X|}
\hline
\rowcolor{popblueL} \textbf{Instruction} & \textbf{nbPaquets} & \textbf{autresArticles} & \textbf{total} \\
\hline
\texttt{Lire(nbPaquets)} & 4 & --- & --- \\
\hline
\texttt{Lire(autresArticles)} & 4 & 2500 & --- \\
\hline
\texttt{total <- ...} & 4 & 2500 & 6500 \\
\hline
\end{tabularx}
\end{center}
Calcul : $4 \times 1000 + 2500 = 4000 + 2500 = 6500$. Message affiché : \og Merci, vous devez payer 6500 FCFA \fg.
\item \textbf{Conclusion :} cet algorithme permet à Mme Etoundi de calculer rapidement et sans erreur la facture de chaque client. Il illustre comment l'informatique automatise les tâches répétitives de la vie quotidienne et fiabilise la gestion d'une boutique.
\end{enumerate}
\textbf{Barème indicatif (10 pts) :} justification de la constante + déclaration (2 pts) ; algorithme complet et correct (4 pts) ; trace et message exact (2 pts) ; conclusion rédigée (2 pts).
\textbf{Points de vigilance :} la constante doit être déclarée \emph{avant} \texttt{Début} ; le message affiché doit respecter exactement la formulation demandée ; le calcul doit être détaillé.
\end{corrige}

\begin{corrige}[Exercice 11 — Type BEPC : le transport scolaire]
\begin{enumerate}
\item Un \textbf{algorithme} est une suite finie et ordonnée d'instructions précises et non ambiguës permettant de résoudre un problème. Les trois instructions de base sont : \texttt{Lire} (saisie), \texttt{Écrire} (affichage) et l'\textbf{affectation} \texttt{<-} (calcul et stockage).
\item \textbf{Entrée :} le nombre d'élèves inscrits : \texttt{nbEleves : Entier}. \textbf{Sorties :} le nombre de cars complets : \texttt{nbCars : Entier} ; le nombre d'élèves restants : \texttt{restants : Entier}.
\item Algorithme :
\begin{lstlisting}
Algorithme Transport
Variables
    nbEleves, nbCars, restants : Entier ;
Constantes
    CAPACITE = 50 : Entier ;
Début
    Écrire("Entrez le nombre d'élèves inscrits : ")
    Lire(nbEleves)
    nbCars <- nbEleves div CAPACITE
    restants <- nbEleves mod CAPACITE
    Écrire("Nombre de cars complets : ")
    Écrire(nbCars)
    Écrire("Élèves dans le dernier car : ")
    Écrire(restants)
Fin
\end{lstlisting}
\item \textbf{Trace pour 137 élèves :} $137 \textbf{ div } 50 = 2$ (car $2 \times 50 = 100 \leq 137$) et $137 \textbf{ mod } 50 = 137 - 100 = 37$. Affichage : 2 cars complets, 37 élèves restants. Vérification : $2 \times 50 + 37 = 137$. \checkmark

\textbf{Trace pour 100 élèves :} $100 \textbf{ div } 50 = 2$ et $100 \textbf{ mod } 50 = 0$. Affichage : 2 cars complets, 0 élève restant. \textbf{Remarque :} dans ce cas, il n'y a \emph{pas} de dernier car partiellement rempli : tous les cars sont complets, le reste est nul.
\item Toutes les variables sont de type \texttt{Entier} car on compte des élèves et des cars, qui sont des quantités entières (on ne peut pas avoir $2,5$ cars ou $37,5$ élèves), et les opérateurs \texttt{div} et \texttt{mod} ne s'appliquent qu'à des entiers.
\end{enumerate}
\textbf{Barème indicatif (10 pts) :} définition + instructions de base (2 pts) ; entrées/sorties avec noms et types (2 pts) ; algorithme correct avec \texttt{div} et \texttt{mod} (3 pts) ; deux traces vérifiées + remarque (2 pts) ; justification du type (1 pt).
\textbf{Points de vigilance :} bien distinguer \texttt{div} (quotient) et \texttt{mod} (reste) ; vérifier chaque trace par la relation $\text{nbEleves} = \text{nbCars} \times 50 + \text{restants}$ ; traiter le cas particulier où le reste vaut $0$.
\end{corrige}

\newpage

\coursec{Fiche bilan}

\begin{popfigure}
\begin{tikzpicture}[mindmap, grow cyclic, every node/.style={concept, text=white, font=\small\bfseries},
  root concept/.append style={concept color=popink, font=\bfseries, minimum size=3.2cm},
  level 1/.append style={level distance=3.6cm, sibling angle=60, minimum size=2.4cm},
  level 2/.append style={level distance=2.2cm, sibling angle=45, minimum size=1.8cm, font=\scriptsize\bfseries}]
\node[root concept]{Langage de définition des algorithmes}
  child[concept color=popblue]{ node{Notion de langage}
    child{ node{Pseudo-code} }
    child{ node{Organigramme} }
  }
  child[concept color=poporange]{ node{Structure de base}
    child{ node{En-tête} }
    child{ node{Déclarations} }
    child{ node{Corps Début/Fin} }
  }
  child[concept color=popgreen]{ node{Variables et constantes}
    child{ node{Nom (identificateur)} }
    child{ node{Valeur} }
  }
  child[concept color=poppurple]{ node{Types de données}
    child{ node{Entier, Réel} }
    child{ node{Caractère, Chaîne} }
    child{ node{Booléen} }
  }
  child[concept color=poppink]{ node{Instructions de base}
    child{ node{Lire (entrée)} }
    child{ node{Écrire (sortie)} }
    child{ node{Affectation $\leftarrow$} }
  }
  child[concept color=popgold]{ node{Rédaction d'algorithmes}
    child{ node{Situations concrètes} }
    child{ node{Justifier la démarche} }
  };
\end{tikzpicture}
\legende{Carte mentale de la leçon : le langage de définition des algorithmes.}
\end{popfigure}

\begin{bilan}
\textbf{Définitions clés}
\begin{itemize}
  \item Un \cle{algorithme} est une suite finie et ordonnée d'instructions permettant de résoudre un problème.
  \item Le \cle{langage de définition des algorithmes} (ou pseudo-code) est un langage simple, en français, qui permet d'écrire un algorithme indépendamment de tout langage de programmation.
  \item Une \cle{variable} est une zone mémoire nommée dont la valeur peut changer ; une \cle{constante} garde toujours la même valeur.
  \item L'\cle{affectation} (symbole $\leftarrow$) donne une valeur à une variable : \texttt{x $\leftarrow$ 5}.
\end{itemize}

\textbf{Structure de base d'un algorithme (à connaître par cœur)}
\begin{lstlisting}
Algorithme nom_de_l_algorithme
Variable x : Entier          (* déclarations *)
Constante TVA : Reel <- 19,25
Debut
    Ecrire("Entrez une valeur : ")
    Lire(x)                  (* entrée *)
    x <- x + 1               (* traitement *)
    Ecrire("Résultat : ", x) (* sortie *)
Fin
\end{lstlisting}

\textbf{Types de données et instructions}
\begin{center}
\begin{tabularx}{\linewidth}{|l|X|X|}
\hline
\rowcolor{popblueL} \textbf{Type} & \textbf{Exemples de valeurs} & \textbf{Usage} \\
\hline
Entier & $-3$, $0$, $25$ & âges, quantités, notes entières \\
\hline
Réel & $3{,}14$ ; $19{,}25$ & prix, moyennes, mesures \\
\hline
Caractère & 'A', 'z' & une seule lettre ou un symbole \\
\hline
Chaîne & "Yaoundé" & noms, mots, phrases \\
\hline
Booléen & Vrai, Faux & réponses oui/non, tests \\
\hline
\end{tabularx}
\end{center}

\textbf{Méthode express pour rédiger un algorithme}
\begin{enumerate}
  \item Comprendre le problème : quelles sont les \cle{données} (entrées) et le \cle{résultat} attendu (sortie) ?
  \item Choisir et déclarer les variables avec le bon type.
  \item Écrire les instructions dans l'ordre : \texttt{Ecrire}/\texttt{Lire} (entrées), calculs par affectation (traitement), \texttt{Ecrire} (sortie).
  \item Vérifier en exécutant l'algorithme « à la main » avec des valeurs d'exemple.
\end{enumerate}
\end{bilan}

\begin{aretenir}[Checklist avant le BEPC]
\begin{itemize}
  \item[$\square$] Je sais définir un algorithme et dire à quoi sert le pseudo-code.
  \item[$\square$] Je connais les trois parties d'un algorithme : en-tête, déclarations, corps entre \texttt{Debut} et \texttt{Fin}.
  \item[$\square$] Je sais distinguer une variable d'une constante.
  \item[$\square$] Je connais les cinq types de données : Entier, Réel, Caractère, Chaîne, Booléen.
  \item[$\square$] Je sais choisir le bon type pour une donnée (ex. une moyenne $\rightarrow$ Réel).
  \item[$\square$] Je sais utiliser \texttt{Lire} pour une entrée au clavier.
  \item[$\square$] Je sais utiliser \texttt{Ecrire} pour afficher un message ou un résultat.
  \item[$\square$] Je sais écrire une affectation avec le symbole $\leftarrow$ (jamais $=$ en pseudo-code).
  \item[$\square$] Je sais que l'affectation remplace l'ancienne valeur de la variable.
  \item[$\square$] Je sais rédiger un algorithme complet sur une situation de la vie courante (calcul d'une moyenne, d'un montant, d'un périmètre...).
  \item[$\square$] Je sais justifier ma démarche et tester mon algorithme avec des valeurs concrètes.
\end{itemize}
\end{aretenir}

\findecours

\end{document}
